\documentclass[a4paper,12pt]{article}

%Class
	%\documentclass[a4paper,12pt]{article}

%Packages
	\usepackage{amsmath}					% Standard maths
	\usepackage{amssymb}					% Standard maths
	\usepackage{amsthm}						% Standard maths
	\usepackage{thmtools, thm-restate}% Theorem styles
	\usepackage{mathtools}
	\usepackage[newzealand]{babel}% Date/time in British format (yes, I know it says NZ)
	\usepackage{datetime}					% Date/time
	\usepackage{multirow}					% For stacked subscripts
	\usepackage{xcolor}						% For colour!
	\usepackage[normalem]{ulem}		% For strikeout
	\usepackage{isomath}
	\usepackage{fancyhdr}					% For the headers
	\usepackage{mfirstuc}					% For the headers
	\usepackage[hmargin=1in,top=0.8in,bottom=1in]{geometry} % Page margins more reasonable
	\usepackage{graphicx}					% For including graphics
	%\usepackage{float}						% For something
	\usepackage{enumitem}					% For numbering tweaks
	\usepackage[margin=15pt,font={small},labelfont=bf,format=hang]{caption} % Changing caption style
	\usepackage{url}
	\usepackage[hidelinks]{hyperref}	% For hyperlinks
	\usepackage{breakurl}
	\usepackage{cleveref}
	\usepackage[scaled=.92]{helvet}
	\usepackage[]{FiraSans}
	\usepackage[T1]{fontenc}
	\usepackage{comment}
	\newbool{show_question}
	\newbool{show_solution}
	\newbool{show_commentary}
	\newcommand{\solutioninheading}{\ifbool{show_solution}{: SOLUTIONS\ifbool{show_commentary}{ \& FEEDBACK}{}}{}}
	
	\usepackage{pgfplots}
\newcommand{\ep}{\varepsilon}	

    \newcommand{\question}[3]{
    	\item %\textbf{#1} \par \medskip % Title
    	%
    	\ifbool{show_question}{#2}{} % Question
    	%
    	%
    	\ifbool{show_question}{\ifbool{show_solution}{\par\medskip\textbf{Solution:}}{}}{} % Both
    	\ifbool{show_solution}{#3}{} % Answer
    }

\usepackage[most]{tcolorbox}

%MATHEMATICAL SHORTCUTS
%Two and three-part definitions
	\newcommand{\twopartdef}[4]
	{	\left\{
			\begin{array}{ll}
				#1 & \mbox{if } #2 \\
				#3 & \mbox{if } #4
			\end{array}
		\right.}
%	\newcommand{\threepartdef}[6]
%	{	\left\{
%			\begin{array}{ll}
%				#1 & \mbox{if } #2 \\
%				#3 & \mbox{if } #4 \\
%				#5 & \mbox{if } #6
%			\end{array}
%		\right.}
%Blackboard bold	
	\newcommand{\N}{{\mathbb N}} 
	\newcommand{\R}{{\mathbb R}} 
	\newcommand{\C}{{\mathbb C}} 
	%Curly letters
%	\newcommand{\E}{{\mathcal E}}	
%	\newcommand{\F}{{\mathcal F}} 
%	\newcommand{\Null}{{\mathcal N}} 
%	\newcommand{\R}{\mathcal{R}} 
%	\newcommand{\Z}{\mathcal{Z}} 
%Uprights
	\newcommand{\D}{{\mathrm D}}
	\renewcommand{\d}{{\mathrm d}}
%Easy Greeks
%	\def\a{\alpha}
%	\def\g{\gamma}
%	\newcommand{\e}{{\varepsilon}}
%	\newcommand{\la}{{\lambda}}  
%	\newcommand{\om}{{\Omega}}
	\renewcommand{\epsilon}{\varepsilon}
	
%Vector operators
	\def \p {\partial}
	\newcommand{\del}{{\nabla}}
	\newcommand{\grad}{{\nabla}}
	\newcommand{\cross}{{\times}}
	\newcommand{\Lap}{{\nabla^2}} 
	\newcommand{\vnabla}{\boldsymbol{\nabla}}
%Arrows
%	\newcommand{\goesto}[1]{\xrightarrow[#1]{}}
%hats
%	\newcommand{\nhat}{\widehat{\v{n}}}
%	\newcommand{\shat}{\widehat{\v{s}}}
	\renewcommand{\hat}{\widehat}
%Note to self
%	\newcommand{\note}[1]{[\textbf{\textsl{\MakeUppercase{#1}}}]}
	
%Stress tensor
%	\newcommand{\T}{\vg{\sigma}}

% Vectors, quaternions etc.
\renewcommand{\v}[1]{\mathbfit{#1}}      % Generic vector
\newcommand{\vc}[1]{\bm{\mathcal{#1}}}      % vector mathcal
\newcommand{\uv}[1]{\widehat{\mathbfit{#1}}}      % Generic unit vector
\newcommand{\q}[1]{\mathsfbfit{#1}}        % Generic quaternion
\renewcommand{\t}[1]{\mathsfbfit{#1}}	 % Generic tensor
\newcommand{\m}[1]{\mathsfbfit{#1}}	 % Generic matrix
\newcommand{\vu}[1]{\mathbf{#1}}         % Upright vector (for zero vector)
\newcommand{\qu}[1]{\mathbf{#1}}       % Upright quaternion (for zero quaternion)
\newcommand{\tu}[1]{\bm{\mathsf{#1}}}	 % Upright tensor (for zero tensor)

%Real and imaginary
\renewcommand{\Re}{\operatorname{Re}}
\renewcommand{\Im}{\operatorname{Im}}
	
%Non-dim numbers
%	\renewcommand{\Re}{\operatorname{Re}}
%	\newcommand{\Bi}{\operatorname{Bi}}
%	\newcommand{\Fo}{\operatorname{Fo}}
	
%Misc
\DeclareMathSymbol{\Delta}{\mathalpha}{operators}{1} % Keeps Delta upright
\DeclareMathSymbol{\Gamma}{\mathalpha}{operators}{0} % Keeps Gamma upright
	\newcommand{\cosec}{\operatorname{cosec}}
	\newcommand{\cosech}{\operatorname{cosech}}
	\newcommand{\sech}{\operatorname{sech}}
	\newcommand{\tr}{\operatorname{tr}}
	\newcommand{\erf}{\operatorname{erf}}
	\newcommand{\order}{\mathcal{O}}
%	\newcommand{\V}{\mathcal{V}} % 	CurlyV = (Bi V) / (1 + Bi)

% Differentiating 
\newcommand{\fd}[2]{\mathchoice{\frac{\d #1}{\d #2}}{\d #1/\d #2}{\d #1/\d #2}{\d #1/\d #2}}
\newcommand{\fdd}[2]{\mathchoice{\frac{\d^2 #1}{\d #2^2}}{\d^2 #1/\d #2^2}{\d^2 #1/\d #2^2}{\d^2 #1/\d #2^2}}
\newcommand{\fddd}[2]{\mathchoice{\frac{\d^3 #1}{\d #2^3}}{\d^3 #1/\d #2^3}{\d^3 #1/\d #2^3}{\d^3 #1/\d #2^3}}
\newcommand{\fdddd}[2]{\mathchoice{\frac{\d^4 #1}{\d #2^4}}{\d^4 #1/\d #2^4}{\d^4 #1/\d #2^4}{\d^4 #1/\d #2^4}}
\newcommand{\fdn}[3]{\mathchoice{\frac{\d^{#1} #2}{\d #3^{#1}}}{\d^{#1} #2/\d #3^{#1}}{\d^{#1} #2/\d #3^{#1}}{\d^{#1} #2/\d #3^{#1}}}
\newcommand{\pd}[2]{\mathchoice{\frac{\p #1}{\p #2}}{\p #1/\p #2}{\p #1/\p #2}{\p #1/\p #2}}
\newcommand{\pdd}[2]{\mathchoice{\frac{\p^2 #1}{\p #2^2}}{\p^2 #1/\p #2^2}{\p^2 #1/\p #2^2}{\p^2 #1/\p #2^2}}
\newcommand{\pddmixed}[3]{\frac{\p^2 #1}{\p #2 \p #3}}
\newcommand{\pddd}[2]{\frac{\p^3 #1}{\p #2^3}}
\newcommand{\pdn}[3]{\mathchoice{\frac{\p^{#1} #2}{\p #3^{#1}}}{\p^{#1} #2/\p #3^{#1}}{\p^{#1} #2/\p #3^{#1}}{\p^{#1} #2/\p #3^{#1}}}
	
	% Misc
\newcommand{\e}{\mathrm{e}}
\renewcommand{\i}{\mathrm{i}}
\newcommand{\dx}{\,\mathrm{d}x}
\newcommand{\dy}{\,\mathrm{d}y}
\renewcommand{\epsilon}{\varepsilon}
\renewcommand{\Re}{\operatorname{Re}}
\newcommand{\Four}{\mathcal{F}}
	\definecolor{nem-red}{HTML}{FFEDF0}
\definecolor{nem-dark-red}{HTML}{d21f3c}
	\newenvironment{nem}{\begin{tcolorbox}[colback=nem-red,colframe=nem-dark-red,breakable,enhanced,parbox=false,title={\sf\small{Non-examinable}}]}{\end{tcolorbox}}
	
\newenvironment{colour-box-breakable}[3]{
\setlist[itemize]{topsep=-\parskip,before=\vspace{-2mm},itemsep=3pt}
%\setlength{\parskip}{0pt}
\begin{tcolorbox}[
	enhanced,
    breakable,
	colback=#2,
	colframe=#1,
	colbacktitle=#2, 
	frame hidden,
    boxrule=0pt,
    leftrule=2pt,
	toptitle=2.5mm,
	bottomtitle=-.5mm, 
    pad before break=1pt,
	pad after break=6pt,
	parbox=false,
	bottom=10pt,
	%grow to right by=0.5mm,
    %grow to left by=-0.2mm,
	sharp corners=west,
  	fonttitle={\sf\small},
	fontupper={\small},
	title={\color{#1}#3}, 
  	overlay first={
	  \fill[#1]
      ([xshift=2pt]frame.north west) --
      ([xshift=2pt]frame.south west) --
      ([xshift=0pt]frame.south west) --
      ([yshift=-2pt]frame.north west) arc
      [radius=2pt,start angle=180,end angle=90]--
      cycle;
    },
     overlay unbroken={
	  \fill[#1]
      ([xshift=2pt]frame.north west) --
      ([xshift=2pt]frame.south west) arc
      [radius=2pt,start angle=-90,end angle=-180] --
      ([yshift=-2pt]frame.north west) arc
      [radius=2pt,start angle=180,end angle=90]--
      cycle;
    },
  	overlay last={
	  \fill[#1]
      ([xshift=2pt]frame.north west) --
      ([xshift=2pt]frame.south west) arc
      [radius=2pt,start angle=-90,end angle=-180] --
      ([yshift=-2pt]frame.north west) --
      ([yshift=0pt]frame.north west) --
      cycle;
    },
]}{\end{tcolorbox}}

\definecolor{commentary-green}{HTML}{FBF3EE}
\definecolor{commentary-dark-green}{HTML}{c57a44}

\newcommand{\commentary}[1]
{
\ifbool{show_commentary}{
	\begin{colour-box-breakable}
		{commentary-dark-green}
		{commentary-green}
		{\footnotesize\raisebox{-2pt}{\includegraphics[width=1em]{../common/squirrel.png}}\,\,Feedback}
	#1
\end{colour-box-breakable}
	}{}
}
%}{
%  \newenvironment{commentary}{\commentaryA}{\endcommentary}
%}

%Depth of display breaks to allow	
	\allowdisplaybreaks[1]

%How deep do we want the Table of Contents to go?	
%	\setcounter{tocdepth}{1}

%Sections
\makeatletter
\renewcommand{\section}{\@startsection
{section}%                   % the name
{1}%                         % the level
{0mm}%                       % the indent
{-\baselineskip}%            % the before skip
{0.5\baselineskip}%          % the after skip
{\normalfont\bfseries}} % the style
\makeatother

% Wavy underline
\makeatletter
\font\sixly=lasyb10 scaled 1200
\def\tinners{\bgroup \markoverwith{\lower8.5\p@\hbox{\sixly \char58}}\ULon}
\makeatother
\newcommand{\answer}[1]{\tinners{\; #1 \;}}

%Setting the footnotes to use *,**,+ etc rather than 1,2,3	
\renewcommand{\thefootnote}{\fnsymbol{footnote}}

%Italic quote environment
%\newenvironment{quoteitalic}{\begin{quote}\itshape}{\end{quote}}

%Heading
\newcommand{\heading}[2]{
\setlength{\parindent}{0pt}
\textbf{MATH3171 (Mathematical Biology III)}

\textbf{YEAR 2025--26, \MakeUppercase{#1} TERM}
\setlength{\parskip}{3ex}

\textbf{\MakeUppercase{#2}}
%
\setlength{\parskip}{2ex}

}

%========================================================================================================================
%DOCUMENT BEGINS HERE!

\usepackage{pgfplots}
\usepackage[export]{adjustbox}
\newcommand{\xmas}[1]{%
  \includegraphics[width=1.5em,valign=c]{images/#1}%
}
\setbool{show_question}{true}
\setbool{show_solution}{false}
\setbool{show_commentary}{false}
\begin{document}
%
\heading{Michaelmas}{Problem sheet 4\solutioninheading}
%
%
\ifbool{show_solution}{}{
  This problem sheet is due in by \textbf{5pm} on \textbf{Wednesday 10 December 2025}. You may submit it on paper to my office (\textsc{mcs}\,3008) or electronically on Gradescope.

  \textbf{Topics covered}: Chapter 6.

  This whole problem sheet concerns the version of the spatial predator--prey system for two populations, $u(x,t)$ and $v(x,t)$, that we saw in class:
  \begin{subequations}
    \begin{align}
      \pd{u}{t} &= u(1-u-v), \label{u-eq} \tag{\xmas{santa}}\\
      \pd{v}{t} &= \pdd{v}{x} + \alpha v(u-\beta),
      \tag{\xmas{christmas-tree}}
    \end{align}
  \end{subequations}
  where the constants $\alpha,\beta > 0$.

  In class we stated that the homogeneous system has three equilibria,
  \begin{itemize}
    \item $(0,0)$, a saddle point,
    \item $(1,0)$, a saddle point,
    \item $(\beta, 1-\beta)$, a stable node when $4\alpha < \beta/(1-\beta)$, and a stable spiral for $4\alpha > \beta/(1-\beta)$,
  \end{itemize}
  and that this suggests possible trajectories between our equilibria:
  \begin{center}
    from $(0,0)$ to $(\beta,1-\beta)$ \qquad and \qquad
    from $(1,0)$ to $(\beta,1-\beta)$.
  \end{center}
}

\begin{enumerate}[leftmargin=0.6cm]
    % Q1
    \question{}{ % Question
      \begin{enumerate}
        \item Draw a phase plane of the homogeneous system: plot $v$ against $u$, marking on the equilibria, nullclines, and some sample trajectories.
        \item Recall that $(0,0)$ is a saddle point. Look at trajectories leaving $(0,0)$. What does this suggest about the existence of a trajectory from $(0,0)$ to $(\beta,1-\beta)$?
      \end{enumerate}
    }
    { % Solution
      \begin{enumerate}
        \item The homogeneous system is
          \begin{subequations}
            \begin{align}
              \fd{u}{t} &= u(1-u-v),\\
              \fd{v}{t} &= \alpha v (u-\beta),
            \end{align}
          \end{subequations}
          with nullclines
          \begin{subequations}
            \begin{align}
              \fd{u}{t} = 0 &\implies u = 0, \quad v = 1 - u\\
              \fd{v}{t} = 0 &\implies u = \beta, \quad v = 0.
            \end{align}
          \end{subequations}
          The phase portrait looks like:
          \begin{center}
            \hspace{-2cm}
            %% Creator: Matplotlib, PGF backend
%%
%% To include the figure in your LaTeX document, write
%%   \input{<filename>.pgf}
%%
%% Make sure the required packages are loaded in your preamble
%%   \usepackage{pgf}
%%
%% and, on pdftex
%%   \usepackage[utf8]{inputenc}\DeclareUnicodeCharacter{2212}{-}
%%
%% or, on luatex and xetex
%%   \usepackage{unicode-math}
%%
%% Figures using additional raster images can only be included by \input if
%% they are in the same directory as the main LaTeX file. For loading figures
%% from other directories you can use the `import` package
%%   \usepackage{import}
%%
%% and then include the figures with
%%   \import{<path to file>}{<filename>.pgf}
%%
%% Matplotlib used the following preamble
%%   \usepackage{fontspec}
%%   \setmainfont{DejaVuSerif.ttf}[Path=/Users/adam/opt/anaconda3/lib/python3.8/site-packages/matplotlib/mpl-data/fonts/ttf/]
%%   \setsansfont{DejaVuSans.ttf}[Path=/Users/adam/opt/anaconda3/lib/python3.8/site-packages/matplotlib/mpl-data/fonts/ttf/]
%%   \setmonofont{DejaVuSansMono.ttf}[Path=/Users/adam/opt/anaconda3/lib/python3.8/site-packages/matplotlib/mpl-data/fonts/ttf/]
%%
\begingroup%
\makeatletter%
\begin{pgfpicture}%
\pgfpathrectangle{\pgfpointorigin}{\pgfqpoint{3.000000in}{2.200000in}}%
\pgfusepath{use as bounding box, clip}%
\begin{pgfscope}%
\pgfsetbuttcap%
\pgfsetmiterjoin%
\definecolor{currentfill}{rgb}{1.000000,1.000000,1.000000}%
\pgfsetfillcolor{currentfill}%
\pgfsetlinewidth{0.000000pt}%
\definecolor{currentstroke}{rgb}{1.000000,1.000000,1.000000}%
\pgfsetstrokecolor{currentstroke}%
\pgfsetdash{}{0pt}%
\pgfpathmoveto{\pgfqpoint{0.000000in}{0.000000in}}%
\pgfpathlineto{\pgfqpoint{3.000000in}{0.000000in}}%
\pgfpathlineto{\pgfqpoint{3.000000in}{2.200000in}}%
\pgfpathlineto{\pgfqpoint{0.000000in}{2.200000in}}%
\pgfpathclose%
\pgfusepath{fill}%
\end{pgfscope}%
\begin{pgfscope}%
\pgfsetbuttcap%
\pgfsetmiterjoin%
\definecolor{currentfill}{rgb}{1.000000,1.000000,1.000000}%
\pgfsetfillcolor{currentfill}%
\pgfsetlinewidth{0.000000pt}%
\definecolor{currentstroke}{rgb}{0.000000,0.000000,0.000000}%
\pgfsetstrokecolor{currentstroke}%
\pgfsetstrokeopacity{0.000000}%
\pgfsetdash{}{0pt}%
\pgfpathmoveto{\pgfqpoint{0.506574in}{0.278150in}}%
\pgfpathlineto{\pgfqpoint{2.966667in}{0.278150in}}%
\pgfpathlineto{\pgfqpoint{2.966667in}{2.166667in}}%
\pgfpathlineto{\pgfqpoint{0.506574in}{2.166667in}}%
\pgfpathclose%
\pgfusepath{fill}%
\end{pgfscope}%
\begin{pgfscope}%
\pgfsetbuttcap%
\pgfsetroundjoin%
\definecolor{currentfill}{rgb}{0.000000,0.000000,0.000000}%
\pgfsetfillcolor{currentfill}%
\pgfsetlinewidth{0.803000pt}%
\definecolor{currentstroke}{rgb}{0.000000,0.000000,0.000000}%
\pgfsetstrokecolor{currentstroke}%
\pgfsetdash{}{0pt}%
\pgfsys@defobject{currentmarker}{\pgfqpoint{0.000000in}{-0.048611in}}{\pgfqpoint{0.000000in}{0.000000in}}{%
\pgfpathmoveto{\pgfqpoint{0.000000in}{0.000000in}}%
\pgfpathlineto{\pgfqpoint{0.000000in}{-0.048611in}}%
\pgfusepath{stroke,fill}%
}%
\begin{pgfscope}%
\pgfsys@transformshift{0.506574in}{0.278150in}%
\pgfsys@useobject{currentmarker}{}%
\end{pgfscope}%
\end{pgfscope}%
\begin{pgfscope}%
\definecolor{textcolor}{rgb}{0.000000,0.000000,0.000000}%
\pgfsetstrokecolor{textcolor}%
\pgfsetfillcolor{textcolor}%
\pgftext[x=0.506574in,y=0.180927in,,top]{\color{textcolor}\rmfamily\fontsize{12.000000}{14.400000}\selectfont 0}%
\end{pgfscope}%
\begin{pgfscope}%
\pgfsetbuttcap%
\pgfsetroundjoin%
\definecolor{currentfill}{rgb}{0.000000,0.000000,0.000000}%
\pgfsetfillcolor{currentfill}%
\pgfsetlinewidth{0.803000pt}%
\definecolor{currentstroke}{rgb}{0.000000,0.000000,0.000000}%
\pgfsetstrokecolor{currentstroke}%
\pgfsetdash{}{0pt}%
\pgfsys@defobject{currentmarker}{\pgfqpoint{0.000000in}{-0.048611in}}{\pgfqpoint{0.000000in}{0.000000in}}{%
\pgfpathmoveto{\pgfqpoint{0.000000in}{0.000000in}}%
\pgfpathlineto{\pgfqpoint{0.000000in}{-0.048611in}}%
\pgfusepath{stroke,fill}%
}%
\begin{pgfscope}%
\pgfsys@transformshift{1.275353in}{0.278150in}%
\pgfsys@useobject{currentmarker}{}%
\end{pgfscope}%
\end{pgfscope}%
\begin{pgfscope}%
\definecolor{textcolor}{rgb}{0.000000,0.000000,0.000000}%
\pgfsetstrokecolor{textcolor}%
\pgfsetfillcolor{textcolor}%
\pgftext[x=1.275353in,y=0.180927in,,top]{\color{textcolor}\rmfamily\fontsize{12.000000}{14.400000}\selectfont \(\displaystyle \beta\)}%
\end{pgfscope}%
\begin{pgfscope}%
\pgfsetbuttcap%
\pgfsetroundjoin%
\definecolor{currentfill}{rgb}{0.000000,0.000000,0.000000}%
\pgfsetfillcolor{currentfill}%
\pgfsetlinewidth{0.803000pt}%
\definecolor{currentstroke}{rgb}{0.000000,0.000000,0.000000}%
\pgfsetstrokecolor{currentstroke}%
\pgfsetdash{}{0pt}%
\pgfsys@defobject{currentmarker}{\pgfqpoint{0.000000in}{-0.048611in}}{\pgfqpoint{0.000000in}{0.000000in}}{%
\pgfpathmoveto{\pgfqpoint{0.000000in}{0.000000in}}%
\pgfpathlineto{\pgfqpoint{0.000000in}{-0.048611in}}%
\pgfusepath{stroke,fill}%
}%
\begin{pgfscope}%
\pgfsys@transformshift{2.044132in}{0.278150in}%
\pgfsys@useobject{currentmarker}{}%
\end{pgfscope}%
\end{pgfscope}%
\begin{pgfscope}%
\definecolor{textcolor}{rgb}{0.000000,0.000000,0.000000}%
\pgfsetstrokecolor{textcolor}%
\pgfsetfillcolor{textcolor}%
\pgftext[x=2.044132in,y=0.180927in,,top]{\color{textcolor}\rmfamily\fontsize{12.000000}{14.400000}\selectfont 1}%
\end{pgfscope}%
\begin{pgfscope}%
\pgfsetbuttcap%
\pgfsetroundjoin%
\definecolor{currentfill}{rgb}{0.000000,0.000000,0.000000}%
\pgfsetfillcolor{currentfill}%
\pgfsetlinewidth{0.803000pt}%
\definecolor{currentstroke}{rgb}{0.000000,0.000000,0.000000}%
\pgfsetstrokecolor{currentstroke}%
\pgfsetdash{}{0pt}%
\pgfsys@defobject{currentmarker}{\pgfqpoint{-0.048611in}{0.000000in}}{\pgfqpoint{-0.000000in}{0.000000in}}{%
\pgfpathmoveto{\pgfqpoint{-0.000000in}{0.000000in}}%
\pgfpathlineto{\pgfqpoint{-0.048611in}{0.000000in}}%
\pgfusepath{stroke,fill}%
}%
\begin{pgfscope}%
\pgfsys@transformshift{0.506574in}{0.278150in}%
\pgfsys@useobject{currentmarker}{}%
\end{pgfscope}%
\end{pgfscope}%
\begin{pgfscope}%
\definecolor{textcolor}{rgb}{0.000000,0.000000,0.000000}%
\pgfsetstrokecolor{textcolor}%
\pgfsetfillcolor{textcolor}%
\pgftext[x=0.303313in, y=0.214836in, left, base]{\color{textcolor}\rmfamily\fontsize{12.000000}{14.400000}\selectfont 0}%
\end{pgfscope}%
\begin{pgfscope}%
\pgfsetbuttcap%
\pgfsetroundjoin%
\definecolor{currentfill}{rgb}{0.000000,0.000000,0.000000}%
\pgfsetfillcolor{currentfill}%
\pgfsetlinewidth{0.803000pt}%
\definecolor{currentstroke}{rgb}{0.000000,0.000000,0.000000}%
\pgfsetstrokecolor{currentstroke}%
\pgfsetdash{}{0pt}%
\pgfsys@defobject{currentmarker}{\pgfqpoint{-0.048611in}{0.000000in}}{\pgfqpoint{-0.000000in}{0.000000in}}{%
\pgfpathmoveto{\pgfqpoint{-0.000000in}{0.000000in}}%
\pgfpathlineto{\pgfqpoint{-0.048611in}{0.000000in}}%
\pgfusepath{stroke,fill}%
}%
\begin{pgfscope}%
\pgfsys@transformshift{0.506574in}{0.868311in}%
\pgfsys@useobject{currentmarker}{}%
\end{pgfscope}%
\end{pgfscope}%
\begin{pgfscope}%
\definecolor{textcolor}{rgb}{0.000000,0.000000,0.000000}%
\pgfsetstrokecolor{textcolor}%
\pgfsetfillcolor{textcolor}%
\pgftext[x=0.022683in, y=0.804997in, left, base]{\color{textcolor}\rmfamily\fontsize{12.000000}{14.400000}\selectfont \(\displaystyle 1-\beta\)}%
\end{pgfscope}%
\begin{pgfscope}%
\pgfsetbuttcap%
\pgfsetroundjoin%
\definecolor{currentfill}{rgb}{0.000000,0.000000,0.000000}%
\pgfsetfillcolor{currentfill}%
\pgfsetlinewidth{0.803000pt}%
\definecolor{currentstroke}{rgb}{0.000000,0.000000,0.000000}%
\pgfsetstrokecolor{currentstroke}%
\pgfsetdash{}{0pt}%
\pgfsys@defobject{currentmarker}{\pgfqpoint{-0.048611in}{0.000000in}}{\pgfqpoint{-0.000000in}{0.000000in}}{%
\pgfpathmoveto{\pgfqpoint{-0.000000in}{0.000000in}}%
\pgfpathlineto{\pgfqpoint{-0.048611in}{0.000000in}}%
\pgfusepath{stroke,fill}%
}%
\begin{pgfscope}%
\pgfsys@transformshift{0.506574in}{1.458473in}%
\pgfsys@useobject{currentmarker}{}%
\end{pgfscope}%
\end{pgfscope}%
\begin{pgfscope}%
\definecolor{textcolor}{rgb}{0.000000,0.000000,0.000000}%
\pgfsetstrokecolor{textcolor}%
\pgfsetfillcolor{textcolor}%
\pgftext[x=0.303313in, y=1.395159in, left, base]{\color{textcolor}\rmfamily\fontsize{12.000000}{14.400000}\selectfont 1}%
\end{pgfscope}%
\begin{pgfscope}%
\pgfsetbuttcap%
\pgfsetroundjoin%
\pgfsetlinewidth{0.702625pt}%
\definecolor{currentstroke}{rgb}{0.121569,0.466667,0.705882}%
\pgfsetstrokecolor{currentstroke}%
\pgfsetdash{{2.590000pt}{1.120000pt}}{0.000000pt}%
\pgfpathmoveto{\pgfqpoint{2.044132in}{0.278150in}}%
\pgfpathlineto{\pgfqpoint{0.506574in}{1.458473in}}%
\pgfusepath{stroke}%
\end{pgfscope}%
\begin{pgfscope}%
\pgfsetrectcap%
\pgfsetmiterjoin%
\pgfsetlinewidth{0.803000pt}%
\definecolor{currentstroke}{rgb}{0.000000,0.000000,0.000000}%
\pgfsetstrokecolor{currentstroke}%
\pgfsetdash{}{0pt}%
\pgfpathmoveto{\pgfqpoint{0.506574in}{0.278150in}}%
\pgfpathlineto{\pgfqpoint{0.506574in}{2.166667in}}%
\pgfusepath{stroke}%
\end{pgfscope}%
\begin{pgfscope}%
\pgfsetrectcap%
\pgfsetmiterjoin%
\pgfsetlinewidth{0.803000pt}%
\definecolor{currentstroke}{rgb}{0.000000,0.000000,0.000000}%
\pgfsetstrokecolor{currentstroke}%
\pgfsetdash{}{0pt}%
\pgfpathmoveto{\pgfqpoint{0.506574in}{0.278150in}}%
\pgfpathlineto{\pgfqpoint{2.966667in}{0.278150in}}%
\pgfusepath{stroke}%
\end{pgfscope}%
\begin{pgfscope}%
\pgfsetrectcap%
\pgfsetroundjoin%
\pgfsetlinewidth{1.505625pt}%
\definecolor{currentstroke}{rgb}{0.172549,0.627451,0.172549}%
\pgfsetstrokecolor{currentstroke}%
\pgfsetdash{}{0pt}%
\pgfpathmoveto{\pgfqpoint{0.521949in}{0.289953in}}%
\pgfpathlineto{\pgfqpoint{0.528452in}{0.288049in}}%
\pgfpathlineto{\pgfqpoint{0.538587in}{0.286349in}}%
\pgfpathlineto{\pgfqpoint{0.554666in}{0.284867in}}%
\pgfpathlineto{\pgfqpoint{0.582675in}{0.283534in}}%
\pgfpathlineto{\pgfqpoint{0.632487in}{0.282405in}}%
\pgfpathlineto{\pgfqpoint{0.729670in}{0.281459in}}%
\pgfpathlineto{\pgfqpoint{0.935277in}{0.280747in}}%
\pgfpathlineto{\pgfqpoint{1.369678in}{0.280496in}}%
\pgfpathlineto{\pgfqpoint{1.769880in}{0.281199in}}%
\pgfpathlineto{\pgfqpoint{1.915120in}{0.282397in}}%
\pgfpathlineto{\pgfqpoint{1.972558in}{0.283854in}}%
\pgfpathlineto{\pgfqpoint{1.998658in}{0.285526in}}%
\pgfpathlineto{\pgfqpoint{2.011291in}{0.287351in}}%
\pgfpathlineto{\pgfqpoint{2.017519in}{0.289251in}}%
\pgfpathlineto{\pgfqpoint{2.020630in}{0.291270in}}%
\pgfpathlineto{\pgfqpoint{2.021920in}{0.293328in}}%
\pgfpathlineto{\pgfqpoint{2.022137in}{0.295711in}}%
\pgfpathlineto{\pgfqpoint{2.021263in}{0.298917in}}%
\pgfpathlineto{\pgfqpoint{2.018657in}{0.303793in}}%
\pgfpathlineto{\pgfqpoint{2.012882in}{0.311905in}}%
\pgfpathlineto{\pgfqpoint{2.000474in}{0.327096in}}%
\pgfpathlineto{\pgfqpoint{1.972337in}{0.359476in}}%
\pgfpathlineto{\pgfqpoint{1.921860in}{0.415816in}}%
\pgfpathlineto{\pgfqpoint{1.893420in}{0.446951in}}%
\pgfpathlineto{\pgfqpoint{1.893420in}{0.446951in}}%
\pgfusepath{stroke}%
\end{pgfscope}%
\begin{pgfscope}%
\pgfsetrectcap%
\pgfsetroundjoin%
\pgfsetlinewidth{1.505625pt}%
\definecolor{currentstroke}{rgb}{0.172549,0.627451,0.172549}%
\pgfsetstrokecolor{currentstroke}%
\pgfsetdash{}{0pt}%
\pgfpathmoveto{\pgfqpoint{2.028756in}{0.289953in}}%
\pgfpathlineto{\pgfqpoint{2.028179in}{0.292546in}}%
\pgfpathlineto{\pgfqpoint{2.025970in}{0.296628in}}%
\pgfpathlineto{\pgfqpoint{2.020538in}{0.304009in}}%
\pgfpathlineto{\pgfqpoint{2.006831in}{0.320390in}}%
\pgfpathlineto{\pgfqpoint{1.969870in}{0.362474in}}%
\pgfpathlineto{\pgfqpoint{1.912692in}{0.425961in}}%
\pgfpathlineto{\pgfqpoint{1.856272in}{0.487015in}}%
\pgfpathlineto{\pgfqpoint{1.802828in}{0.543260in}}%
\pgfpathlineto{\pgfqpoint{1.752624in}{0.594522in}}%
\pgfpathlineto{\pgfqpoint{1.705451in}{0.641128in}}%
\pgfpathlineto{\pgfqpoint{1.662147in}{0.682416in}}%
\pgfpathlineto{\pgfqpoint{1.620919in}{0.720228in}}%
\pgfpathlineto{\pgfqpoint{1.582648in}{0.753859in}}%
\pgfpathlineto{\pgfqpoint{1.546278in}{0.784339in}}%
\pgfpathlineto{\pgfqpoint{1.512241in}{0.811386in}}%
\pgfpathlineto{\pgfqpoint{1.480823in}{0.834916in}}%
\pgfpathlineto{\pgfqpoint{1.452173in}{0.855005in}}%
\pgfpathlineto{\pgfqpoint{1.424932in}{0.872726in}}%
\pgfpathlineto{\pgfqpoint{1.400603in}{0.887249in}}%
\pgfpathlineto{\pgfqpoint{1.377850in}{0.899549in}}%
\pgfpathlineto{\pgfqpoint{1.356771in}{0.909659in}}%
\pgfpathlineto{\pgfqpoint{1.338460in}{0.917267in}}%
\pgfpathlineto{\pgfqpoint{1.321732in}{0.923087in}}%
\pgfpathlineto{\pgfqpoint{1.306594in}{0.927242in}}%
\pgfpathlineto{\pgfqpoint{1.293035in}{0.929869in}}%
\pgfpathlineto{\pgfqpoint{1.281025in}{0.931120in}}%
\pgfpathlineto{\pgfqpoint{1.270521in}{0.931152in}}%
\pgfpathlineto{\pgfqpoint{1.261466in}{0.930126in}}%
\pgfpathlineto{\pgfqpoint{1.253795in}{0.928198in}}%
\pgfpathlineto{\pgfqpoint{1.247434in}{0.925523in}}%
\pgfpathlineto{\pgfqpoint{1.242037in}{0.922035in}}%
\pgfpathlineto{\pgfqpoint{1.237919in}{0.918033in}}%
\pgfpathlineto{\pgfqpoint{1.234974in}{0.913661in}}%
\pgfpathlineto{\pgfqpoint{1.233020in}{0.908792in}}%
\pgfpathlineto{\pgfqpoint{1.232267in}{0.905117in}}%
\pgfpathlineto{\pgfqpoint{1.232267in}{0.905117in}}%
\pgfusepath{stroke}%
\end{pgfscope}%
\begin{pgfscope}%
\pgfsetrectcap%
\pgfsetroundjoin%
\pgfsetlinewidth{1.505625pt}%
\definecolor{currentstroke}{rgb}{0.172549,0.627451,0.172549}%
\pgfsetstrokecolor{currentstroke}%
\pgfsetstrokeopacity{0.200000}%
\pgfsetdash{}{0pt}%
\pgfpathmoveto{\pgfqpoint{0.660329in}{0.868311in}}%
\pgfpathlineto{\pgfqpoint{0.668078in}{0.840793in}}%
\pgfpathlineto{\pgfqpoint{0.676570in}{0.814899in}}%
\pgfpathlineto{\pgfqpoint{0.685843in}{0.790551in}}%
\pgfpathlineto{\pgfqpoint{0.694625in}{0.770459in}}%
\pgfpathlineto{\pgfqpoint{0.704059in}{0.751449in}}%
\pgfpathlineto{\pgfqpoint{0.714168in}{0.733477in}}%
\pgfpathlineto{\pgfqpoint{0.724979in}{0.716501in}}%
\pgfpathlineto{\pgfqpoint{0.736515in}{0.700480in}}%
\pgfpathlineto{\pgfqpoint{0.748799in}{0.685376in}}%
\pgfpathlineto{\pgfqpoint{0.761851in}{0.671152in}}%
\pgfpathlineto{\pgfqpoint{0.775690in}{0.657774in}}%
\pgfpathlineto{\pgfqpoint{0.790330in}{0.645208in}}%
\pgfpathlineto{\pgfqpoint{0.805783in}{0.633424in}}%
\pgfpathlineto{\pgfqpoint{0.822057in}{0.622391in}}%
\pgfpathlineto{\pgfqpoint{0.839152in}{0.612081in}}%
\pgfpathlineto{\pgfqpoint{0.857065in}{0.602468in}}%
\pgfpathlineto{\pgfqpoint{0.878524in}{0.592302in}}%
\pgfpathlineto{\pgfqpoint{0.901012in}{0.582978in}}%
\pgfpathlineto{\pgfqpoint{0.924487in}{0.574462in}}%
\pgfpathlineto{\pgfqpoint{0.948898in}{0.566723in}}%
\pgfpathlineto{\pgfqpoint{0.977393in}{0.558911in}}%
\pgfpathlineto{\pgfqpoint{1.006870in}{0.552006in}}%
\pgfpathlineto{\pgfqpoint{1.037189in}{0.545975in}}%
\pgfpathlineto{\pgfqpoint{1.071671in}{0.540257in}}%
\pgfpathlineto{\pgfqpoint{1.106749in}{0.535535in}}%
\pgfpathlineto{\pgfqpoint{1.145704in}{0.531444in}}%
\pgfpathlineto{\pgfqpoint{1.184677in}{0.528465in}}%
\pgfpathlineto{\pgfqpoint{1.226755in}{0.526432in}}%
\pgfpathlineto{\pgfqpoint{1.267885in}{0.525621in}}%
\pgfpathlineto{\pgfqpoint{1.307585in}{0.525984in}}%
\pgfpathlineto{\pgfqpoint{1.345413in}{0.527479in}}%
\pgfpathlineto{\pgfqpoint{1.380981in}{0.530065in}}%
\pgfpathlineto{\pgfqpoint{1.411326in}{0.533360in}}%
\pgfpathlineto{\pgfqpoint{1.439307in}{0.537510in}}%
\pgfpathlineto{\pgfqpoint{1.462574in}{0.542001in}}%
\pgfpathlineto{\pgfqpoint{1.483700in}{0.547151in}}%
\pgfpathlineto{\pgfqpoint{1.502646in}{0.552939in}}%
\pgfpathlineto{\pgfqpoint{1.517828in}{0.558675in}}%
\pgfpathlineto{\pgfqpoint{1.531250in}{0.564890in}}%
\pgfpathlineto{\pgfqpoint{1.542940in}{0.571566in}}%
\pgfpathlineto{\pgfqpoint{1.551910in}{0.577872in}}%
\pgfpathlineto{\pgfqpoint{1.559578in}{0.584513in}}%
\pgfpathlineto{\pgfqpoint{1.565986in}{0.591472in}}%
\pgfpathlineto{\pgfqpoint{1.571180in}{0.598735in}}%
\pgfpathlineto{\pgfqpoint{1.575208in}{0.606285in}}%
\pgfpathlineto{\pgfqpoint{1.578123in}{0.614103in}}%
\pgfpathlineto{\pgfqpoint{1.579977in}{0.622172in}}%
\pgfpathlineto{\pgfqpoint{1.580826in}{0.630471in}}%
\pgfpathlineto{\pgfqpoint{1.580649in}{0.640058in}}%
\pgfpathlineto{\pgfqpoint{1.579350in}{0.649882in}}%
\pgfpathlineto{\pgfqpoint{1.576685in}{0.661037in}}%
\pgfpathlineto{\pgfqpoint{1.572838in}{0.672401in}}%
\pgfpathlineto{\pgfqpoint{1.567366in}{0.685089in}}%
\pgfpathlineto{\pgfqpoint{1.560070in}{0.699080in}}%
\pgfpathlineto{\pgfqpoint{1.550794in}{0.714323in}}%
\pgfpathlineto{\pgfqpoint{1.539440in}{0.730716in}}%
\pgfpathlineto{\pgfqpoint{1.525977in}{0.748110in}}%
\pgfpathlineto{\pgfqpoint{1.510446in}{0.766292in}}%
\pgfpathlineto{\pgfqpoint{1.492969in}{0.784990in}}%
\pgfpathlineto{\pgfqpoint{1.473752in}{0.803869in}}%
\pgfpathlineto{\pgfqpoint{1.454171in}{0.821592in}}%
\pgfpathlineto{\pgfqpoint{1.433476in}{0.838847in}}%
\pgfpathlineto{\pgfqpoint{1.413090in}{0.854446in}}%
\pgfpathlineto{\pgfqpoint{1.392287in}{0.868950in}}%
\pgfpathlineto{\pgfqpoint{1.372430in}{0.881417in}}%
\pgfpathlineto{\pgfqpoint{1.353744in}{0.891830in}}%
\pgfpathlineto{\pgfqpoint{1.336398in}{0.900229in}}%
\pgfpathlineto{\pgfqpoint{1.320511in}{0.906699in}}%
\pgfpathlineto{\pgfqpoint{1.306159in}{0.911364in}}%
\pgfpathlineto{\pgfqpoint{1.293379in}{0.914374in}}%
\pgfpathlineto{\pgfqpoint{1.282173in}{0.915903in}}%
\pgfpathlineto{\pgfqpoint{1.272517in}{0.916136in}}%
\pgfpathlineto{\pgfqpoint{1.264363in}{0.915260in}}%
\pgfpathlineto{\pgfqpoint{1.257642in}{0.913462in}}%
\pgfpathlineto{\pgfqpoint{1.252273in}{0.910923in}}%
\pgfpathlineto{\pgfqpoint{1.248163in}{0.907812in}}%
\pgfpathlineto{\pgfqpoint{1.245089in}{0.904101in}}%
\pgfpathlineto{\pgfqpoint{1.243166in}{0.900094in}}%
\pgfpathlineto{\pgfqpoint{1.242225in}{0.895535in}}%
\pgfpathlineto{\pgfqpoint{1.242363in}{0.890757in}}%
\pgfpathlineto{\pgfqpoint{1.243272in}{0.886677in}}%
\pgfpathlineto{\pgfqpoint{1.243272in}{0.886677in}}%
\pgfusepath{stroke}%
\end{pgfscope}%
\begin{pgfscope}%
\pgfsetrectcap%
\pgfsetroundjoin%
\pgfsetlinewidth{1.505625pt}%
\definecolor{currentstroke}{rgb}{0.172549,0.627451,0.172549}%
\pgfsetstrokecolor{currentstroke}%
\pgfsetstrokeopacity{0.200000}%
\pgfsetdash{}{0pt}%
\pgfpathmoveto{\pgfqpoint{0.660329in}{1.458473in}}%
\pgfpathlineto{\pgfqpoint{0.658823in}{1.396282in}}%
\pgfpathlineto{\pgfqpoint{0.658408in}{1.330917in}}%
\pgfpathlineto{\pgfqpoint{0.659218in}{1.269392in}}%
\pgfpathlineto{\pgfqpoint{0.661181in}{1.211591in}}%
\pgfpathlineto{\pgfqpoint{0.663894in}{1.162640in}}%
\pgfpathlineto{\pgfqpoint{0.667483in}{1.116488in}}%
\pgfpathlineto{\pgfqpoint{0.671942in}{1.073025in}}%
\pgfpathlineto{\pgfqpoint{0.677274in}{1.032141in}}%
\pgfpathlineto{\pgfqpoint{0.683489in}{0.993722in}}%
\pgfpathlineto{\pgfqpoint{0.689770in}{0.961552in}}%
\pgfpathlineto{\pgfqpoint{0.696780in}{0.931168in}}%
\pgfpathlineto{\pgfqpoint{0.704539in}{0.902495in}}%
\pgfpathlineto{\pgfqpoint{0.713067in}{0.875460in}}%
\pgfpathlineto{\pgfqpoint{0.722388in}{0.849994in}}%
\pgfpathlineto{\pgfqpoint{0.731213in}{0.828946in}}%
\pgfpathlineto{\pgfqpoint{0.740679in}{0.809005in}}%
\pgfpathlineto{\pgfqpoint{0.750804in}{0.790128in}}%
\pgfpathlineto{\pgfqpoint{0.761600in}{0.772275in}}%
\pgfpathlineto{\pgfqpoint{0.773082in}{0.755408in}}%
\pgfpathlineto{\pgfqpoint{0.785263in}{0.739490in}}%
\pgfpathlineto{\pgfqpoint{0.798152in}{0.724485in}}%
\pgfpathlineto{\pgfqpoint{0.811757in}{0.710360in}}%
\pgfpathlineto{\pgfqpoint{0.826081in}{0.697081in}}%
\pgfpathlineto{\pgfqpoint{0.841125in}{0.684618in}}%
\pgfpathlineto{\pgfqpoint{0.856885in}{0.672941in}}%
\pgfpathlineto{\pgfqpoint{0.873352in}{0.662022in}}%
\pgfpathlineto{\pgfqpoint{0.890513in}{0.651834in}}%
\pgfpathlineto{\pgfqpoint{0.910949in}{0.641054in}}%
\pgfpathlineto{\pgfqpoint{0.932227in}{0.631160in}}%
\pgfpathlineto{\pgfqpoint{0.954297in}{0.622117in}}%
\pgfpathlineto{\pgfqpoint{0.977100in}{0.613895in}}%
\pgfpathlineto{\pgfqpoint{1.003540in}{0.605586in}}%
\pgfpathlineto{\pgfqpoint{1.030700in}{0.598237in}}%
\pgfpathlineto{\pgfqpoint{1.058446in}{0.591809in}}%
\pgfpathlineto{\pgfqpoint{1.089783in}{0.585704in}}%
\pgfpathlineto{\pgfqpoint{1.121441in}{0.580647in}}%
\pgfpathlineto{\pgfqpoint{1.156357in}{0.576246in}}%
\pgfpathlineto{\pgfqpoint{1.191057in}{0.573010in}}%
\pgfpathlineto{\pgfqpoint{1.225214in}{0.570890in}}%
\pgfpathlineto{\pgfqpoint{1.261486in}{0.569798in}}%
\pgfpathlineto{\pgfqpoint{1.296344in}{0.569924in}}%
\pgfpathlineto{\pgfqpoint{1.329437in}{0.571216in}}%
\pgfpathlineto{\pgfqpoint{1.357959in}{0.573380in}}%
\pgfpathlineto{\pgfqpoint{1.384543in}{0.576443in}}%
\pgfpathlineto{\pgfqpoint{1.409034in}{0.580365in}}%
\pgfpathlineto{\pgfqpoint{1.431318in}{0.585110in}}%
\pgfpathlineto{\pgfqpoint{1.449600in}{0.590104in}}%
\pgfpathlineto{\pgfqpoint{1.465969in}{0.595716in}}%
\pgfpathlineto{\pgfqpoint{1.480423in}{0.601917in}}%
\pgfpathlineto{\pgfqpoint{1.491807in}{0.607974in}}%
\pgfpathlineto{\pgfqpoint{1.501679in}{0.614458in}}%
\pgfpathlineto{\pgfqpoint{1.510075in}{0.621345in}}%
\pgfpathlineto{\pgfqpoint{1.517040in}{0.628609in}}%
\pgfpathlineto{\pgfqpoint{1.522072in}{0.635362in}}%
\pgfpathlineto{\pgfqpoint{1.526056in}{0.642373in}}%
\pgfpathlineto{\pgfqpoint{1.529035in}{0.649623in}}%
\pgfpathlineto{\pgfqpoint{1.531243in}{0.658038in}}%
\pgfpathlineto{\pgfqpoint{1.532304in}{0.666700in}}%
\pgfpathlineto{\pgfqpoint{1.532290in}{0.675577in}}%
\pgfpathlineto{\pgfqpoint{1.531097in}{0.685656in}}%
\pgfpathlineto{\pgfqpoint{1.528759in}{0.695917in}}%
\pgfpathlineto{\pgfqpoint{1.524979in}{0.707363in}}%
\pgfpathlineto{\pgfqpoint{1.519556in}{0.719969in}}%
\pgfpathlineto{\pgfqpoint{1.512331in}{0.733678in}}%
\pgfpathlineto{\pgfqpoint{1.503196in}{0.748392in}}%
\pgfpathlineto{\pgfqpoint{1.492111in}{0.763963in}}%
\pgfpathlineto{\pgfqpoint{1.479103in}{0.780195in}}%
\pgfpathlineto{\pgfqpoint{1.464276in}{0.796835in}}%
\pgfpathlineto{\pgfqpoint{1.447813in}{0.813579in}}%
\pgfpathlineto{\pgfqpoint{1.429969in}{0.830077in}}%
\pgfpathlineto{\pgfqpoint{1.412014in}{0.845194in}}%
\pgfpathlineto{\pgfqpoint{1.393346in}{0.859469in}}%
\pgfpathlineto{\pgfqpoint{1.375228in}{0.871963in}}%
\pgfpathlineto{\pgfqpoint{1.357080in}{0.883113in}}%
\pgfpathlineto{\pgfqpoint{1.340061in}{0.892252in}}%
\pgfpathlineto{\pgfqpoint{1.324347in}{0.899439in}}%
\pgfpathlineto{\pgfqpoint{1.310059in}{0.904779in}}%
\pgfpathlineto{\pgfqpoint{1.297271in}{0.908411in}}%
\pgfpathlineto{\pgfqpoint{1.286014in}{0.910504in}}%
\pgfpathlineto{\pgfqpoint{1.276285in}{0.911240in}}%
\pgfpathlineto{\pgfqpoint{1.268050in}{0.910810in}}%
\pgfpathlineto{\pgfqpoint{1.261253in}{0.909408in}}%
\pgfpathlineto{\pgfqpoint{1.255816in}{0.907221in}}%
\pgfpathlineto{\pgfqpoint{1.251648in}{0.904425in}}%
\pgfpathlineto{\pgfqpoint{1.248532in}{0.901021in}}%
\pgfpathlineto{\pgfqpoint{1.247692in}{0.899696in}}%
\pgfpathlineto{\pgfqpoint{1.247692in}{0.899696in}}%
\pgfusepath{stroke}%
\end{pgfscope}%
\begin{pgfscope}%
\pgfsetrectcap%
\pgfsetroundjoin%
\pgfsetlinewidth{1.505625pt}%
\definecolor{currentstroke}{rgb}{0.172549,0.627451,0.172549}%
\pgfsetstrokecolor{currentstroke}%
\pgfsetstrokeopacity{0.200000}%
\pgfsetdash{}{0pt}%
\pgfpathmoveto{\pgfqpoint{2.812911in}{1.458473in}}%
\pgfpathlineto{\pgfqpoint{2.712532in}{1.493250in}}%
\pgfpathlineto{\pgfqpoint{2.618691in}{1.526683in}}%
\pgfpathlineto{\pgfqpoint{2.488808in}{1.574292in}}%
\pgfpathlineto{\pgfqpoint{2.370402in}{1.618836in}}%
\pgfpathlineto{\pgfqpoint{2.194715in}{1.686275in}}%
\pgfpathlineto{\pgfqpoint{1.959099in}{1.776786in}}%
\pgfpathlineto{\pgfqpoint{1.858316in}{1.814275in}}%
\pgfpathlineto{\pgfqpoint{1.789006in}{1.839130in}}%
\pgfpathlineto{\pgfqpoint{1.724543in}{1.861284in}}%
\pgfpathlineto{\pgfqpoint{1.664510in}{1.880829in}}%
\pgfpathlineto{\pgfqpoint{1.608537in}{1.897862in}}%
\pgfpathlineto{\pgfqpoint{1.556299in}{1.912484in}}%
\pgfpathlineto{\pgfqpoint{1.507504in}{1.924801in}}%
\pgfpathlineto{\pgfqpoint{1.461888in}{1.934918in}}%
\pgfpathlineto{\pgfqpoint{1.419215in}{1.942944in}}%
\pgfpathlineto{\pgfqpoint{1.379269in}{1.948987in}}%
\pgfpathlineto{\pgfqpoint{1.341855in}{1.953153in}}%
\pgfpathlineto{\pgfqpoint{1.306794in}{1.955549in}}%
\pgfpathlineto{\pgfqpoint{1.273922in}{1.956278in}}%
\pgfpathlineto{\pgfqpoint{1.243089in}{1.955442in}}%
\pgfpathlineto{\pgfqpoint{1.214156in}{1.953139in}}%
\pgfpathlineto{\pgfqpoint{1.186997in}{1.949467in}}%
\pgfpathlineto{\pgfqpoint{1.161492in}{1.944517in}}%
\pgfpathlineto{\pgfqpoint{1.137534in}{1.938379in}}%
\pgfpathlineto{\pgfqpoint{1.115021in}{1.931139in}}%
\pgfpathlineto{\pgfqpoint{1.093861in}{1.922877in}}%
\pgfpathlineto{\pgfqpoint{1.073967in}{1.913674in}}%
\pgfpathlineto{\pgfqpoint{1.055259in}{1.903603in}}%
\pgfpathlineto{\pgfqpoint{1.037662in}{1.892735in}}%
\pgfpathlineto{\pgfqpoint{1.021107in}{1.881138in}}%
\pgfpathlineto{\pgfqpoint{1.005530in}{1.868876in}}%
\pgfpathlineto{\pgfqpoint{0.990871in}{1.856009in}}%
\pgfpathlineto{\pgfqpoint{0.977075in}{1.842594in}}%
\pgfpathlineto{\pgfqpoint{0.964089in}{1.828685in}}%
\pgfpathlineto{\pgfqpoint{0.947954in}{1.809457in}}%
\pgfpathlineto{\pgfqpoint{0.933070in}{1.789553in}}%
\pgfpathlineto{\pgfqpoint{0.919343in}{1.769074in}}%
\pgfpathlineto{\pgfqpoint{0.906684in}{1.748116in}}%
\pgfpathlineto{\pgfqpoint{0.895017in}{1.726764in}}%
\pgfpathlineto{\pgfqpoint{0.884268in}{1.705098in}}%
\pgfpathlineto{\pgfqpoint{0.872024in}{1.677682in}}%
\pgfpathlineto{\pgfqpoint{0.860998in}{1.650014in}}%
\pgfpathlineto{\pgfqpoint{0.851089in}{1.622207in}}%
\pgfpathlineto{\pgfqpoint{0.842203in}{1.594359in}}%
\pgfpathlineto{\pgfqpoint{0.832779in}{1.561013in}}%
\pgfpathlineto{\pgfqpoint{0.824589in}{1.527866in}}%
\pgfpathlineto{\pgfqpoint{0.817525in}{1.495033in}}%
\pgfpathlineto{\pgfqpoint{0.810582in}{1.457250in}}%
\pgfpathlineto{\pgfqpoint{0.804916in}{1.420149in}}%
\pgfpathlineto{\pgfqpoint{0.800416in}{1.383831in}}%
\pgfpathlineto{\pgfqpoint{0.796582in}{1.343389in}}%
\pgfpathlineto{\pgfqpoint{0.794032in}{1.304171in}}%
\pgfpathlineto{\pgfqpoint{0.792669in}{1.266246in}}%
\pgfpathlineto{\pgfqpoint{0.792413in}{1.229663in}}%
\pgfpathlineto{\pgfqpoint{0.793195in}{1.194454in}}%
\pgfpathlineto{\pgfqpoint{0.794962in}{1.160636in}}%
\pgfpathlineto{\pgfqpoint{0.797669in}{1.128215in}}%
\pgfpathlineto{\pgfqpoint{0.801279in}{1.097187in}}%
\pgfpathlineto{\pgfqpoint{0.805761in}{1.067542in}}%
\pgfpathlineto{\pgfqpoint{0.811091in}{1.039261in}}%
\pgfpathlineto{\pgfqpoint{0.817249in}{1.012323in}}%
\pgfpathlineto{\pgfqpoint{0.824217in}{0.986701in}}%
\pgfpathlineto{\pgfqpoint{0.831980in}{0.962368in}}%
\pgfpathlineto{\pgfqpoint{0.840526in}{0.939291in}}%
\pgfpathlineto{\pgfqpoint{0.849841in}{0.917438in}}%
\pgfpathlineto{\pgfqpoint{0.859911in}{0.896776in}}%
\pgfpathlineto{\pgfqpoint{0.870724in}{0.877270in}}%
\pgfpathlineto{\pgfqpoint{0.882264in}{0.858887in}}%
\pgfpathlineto{\pgfqpoint{0.894513in}{0.841591in}}%
\pgfpathlineto{\pgfqpoint{0.907451in}{0.825348in}}%
\pgfpathlineto{\pgfqpoint{0.921054in}{0.810124in}}%
\pgfpathlineto{\pgfqpoint{0.935295in}{0.795887in}}%
\pgfpathlineto{\pgfqpoint{0.950142in}{0.782605in}}%
\pgfpathlineto{\pgfqpoint{0.965560in}{0.770244in}}%
\pgfpathlineto{\pgfqpoint{0.981507in}{0.758775in}}%
\pgfpathlineto{\pgfqpoint{1.000023in}{0.746900in}}%
\pgfpathlineto{\pgfqpoint{1.019081in}{0.736075in}}%
\pgfpathlineto{\pgfqpoint{1.038601in}{0.726259in}}%
\pgfpathlineto{\pgfqpoint{1.058498in}{0.717414in}}%
\pgfpathlineto{\pgfqpoint{1.080934in}{0.708680in}}%
\pgfpathlineto{\pgfqpoint{1.103588in}{0.701050in}}%
\pgfpathlineto{\pgfqpoint{1.128593in}{0.693877in}}%
\pgfpathlineto{\pgfqpoint{1.153503in}{0.687923in}}%
\pgfpathlineto{\pgfqpoint{1.178125in}{0.683131in}}%
\pgfpathlineto{\pgfqpoint{1.204437in}{0.679163in}}%
\pgfpathlineto{\pgfqpoint{1.229948in}{0.676443in}}%
\pgfpathlineto{\pgfqpoint{1.254439in}{0.674902in}}%
\pgfpathlineto{\pgfqpoint{1.277712in}{0.674476in}}%
\pgfpathlineto{\pgfqpoint{1.299595in}{0.675100in}}%
\pgfpathlineto{\pgfqpoint{1.319943in}{0.676711in}}%
\pgfpathlineto{\pgfqpoint{1.338641in}{0.679245in}}%
\pgfpathlineto{\pgfqpoint{1.355608in}{0.682639in}}%
\pgfpathlineto{\pgfqpoint{1.370788in}{0.686832in}}%
\pgfpathlineto{\pgfqpoint{1.383112in}{0.691322in}}%
\pgfpathlineto{\pgfqpoint{1.393915in}{0.696381in}}%
\pgfpathlineto{\pgfqpoint{1.403214in}{0.701959in}}%
\pgfpathlineto{\pgfqpoint{1.411042in}{0.708006in}}%
\pgfpathlineto{\pgfqpoint{1.416918in}{0.713869in}}%
\pgfpathlineto{\pgfqpoint{1.421657in}{0.720041in}}%
\pgfpathlineto{\pgfqpoint{1.425307in}{0.726484in}}%
\pgfpathlineto{\pgfqpoint{1.427920in}{0.733159in}}%
\pgfpathlineto{\pgfqpoint{1.429665in}{0.740726in}}%
\pgfpathlineto{\pgfqpoint{1.430305in}{0.748478in}}%
\pgfpathlineto{\pgfqpoint{1.429838in}{0.757089in}}%
\pgfpathlineto{\pgfqpoint{1.428264in}{0.765798in}}%
\pgfpathlineto{\pgfqpoint{1.425434in}{0.775273in}}%
\pgfpathlineto{\pgfqpoint{1.421234in}{0.785440in}}%
\pgfpathlineto{\pgfqpoint{1.415595in}{0.796190in}}%
\pgfpathlineto{\pgfqpoint{1.408499in}{0.807380in}}%
\pgfpathlineto{\pgfqpoint{1.399462in}{0.819490in}}%
\pgfpathlineto{\pgfqpoint{1.389034in}{0.831571in}}%
\pgfpathlineto{\pgfqpoint{1.377418in}{0.843359in}}%
\pgfpathlineto{\pgfqpoint{1.364275in}{0.855077in}}%
\pgfpathlineto{\pgfqpoint{1.350525in}{0.865805in}}%
\pgfpathlineto{\pgfqpoint{1.336552in}{0.875266in}}%
\pgfpathlineto{\pgfqpoint{1.323286in}{0.882941in}}%
\pgfpathlineto{\pgfqpoint{1.310487in}{0.889091in}}%
\pgfpathlineto{\pgfqpoint{1.298923in}{0.893468in}}%
\pgfpathlineto{\pgfqpoint{1.289484in}{0.896065in}}%
\pgfpathlineto{\pgfqpoint{1.289484in}{0.896065in}}%
\pgfusepath{stroke}%
\end{pgfscope}%
\begin{pgfscope}%
\pgfsetrectcap%
\pgfsetroundjoin%
\pgfsetlinewidth{1.505625pt}%
\definecolor{currentstroke}{rgb}{0.172549,0.627451,0.172549}%
\pgfsetstrokecolor{currentstroke}%
\pgfsetstrokeopacity{0.200000}%
\pgfsetdash{}{0pt}%
\pgfpathmoveto{\pgfqpoint{2.812911in}{0.868311in}}%
\pgfpathlineto{\pgfqpoint{2.713476in}{0.894591in}}%
\pgfpathlineto{\pgfqpoint{2.622134in}{0.920247in}}%
\pgfpathlineto{\pgfqpoint{2.537828in}{0.945252in}}%
\pgfpathlineto{\pgfqpoint{2.459685in}{0.969586in}}%
\pgfpathlineto{\pgfqpoint{2.386976in}{0.993227in}}%
\pgfpathlineto{\pgfqpoint{2.297446in}{1.023637in}}%
\pgfpathlineto{\pgfqpoint{2.215311in}{1.052742in}}%
\pgfpathlineto{\pgfqpoint{2.121590in}{1.087237in}}%
\pgfpathlineto{\pgfqpoint{2.020312in}{1.125800in}}%
\pgfpathlineto{\pgfqpoint{1.873691in}{1.183077in}}%
\pgfpathlineto{\pgfqpoint{1.737871in}{1.235934in}}%
\pgfpathlineto{\pgfqpoint{1.662665in}{1.264118in}}%
\pgfpathlineto{\pgfqpoint{1.604107in}{1.284996in}}%
\pgfpathlineto{\pgfqpoint{1.550398in}{1.302941in}}%
\pgfpathlineto{\pgfqpoint{1.508999in}{1.315735in}}%
\pgfpathlineto{\pgfqpoint{1.470396in}{1.326640in}}%
\pgfpathlineto{\pgfqpoint{1.434379in}{1.335732in}}%
\pgfpathlineto{\pgfqpoint{1.400760in}{1.343092in}}%
\pgfpathlineto{\pgfqpoint{1.369369in}{1.348802in}}%
\pgfpathlineto{\pgfqpoint{1.340053in}{1.352950in}}%
\pgfpathlineto{\pgfqpoint{1.312673in}{1.355623in}}%
\pgfpathlineto{\pgfqpoint{1.287103in}{1.356906in}}%
\pgfpathlineto{\pgfqpoint{1.263227in}{1.356889in}}%
\pgfpathlineto{\pgfqpoint{1.240937in}{1.355658in}}%
\pgfpathlineto{\pgfqpoint{1.220136in}{1.353296in}}%
\pgfpathlineto{\pgfqpoint{1.200732in}{1.349887in}}%
\pgfpathlineto{\pgfqpoint{1.182642in}{1.345513in}}%
\pgfpathlineto{\pgfqpoint{1.165787in}{1.340249in}}%
\pgfpathlineto{\pgfqpoint{1.150096in}{1.334173in}}%
\pgfpathlineto{\pgfqpoint{1.135500in}{1.327356in}}%
\pgfpathlineto{\pgfqpoint{1.121938in}{1.319866in}}%
\pgfpathlineto{\pgfqpoint{1.109352in}{1.311770in}}%
\pgfpathlineto{\pgfqpoint{1.097688in}{1.303131in}}%
\pgfpathlineto{\pgfqpoint{1.084837in}{1.292128in}}%
\pgfpathlineto{\pgfqpoint{1.073161in}{1.280524in}}%
\pgfpathlineto{\pgfqpoint{1.062585in}{1.268407in}}%
\pgfpathlineto{\pgfqpoint{1.053043in}{1.255860in}}%
\pgfpathlineto{\pgfqpoint{1.044472in}{1.242960in}}%
\pgfpathlineto{\pgfqpoint{1.035624in}{1.227559in}}%
\pgfpathlineto{\pgfqpoint{1.027937in}{1.211876in}}%
\pgfpathlineto{\pgfqpoint{1.021335in}{1.196004in}}%
\pgfpathlineto{\pgfqpoint{1.015032in}{1.177740in}}%
\pgfpathlineto{\pgfqpoint{1.009964in}{1.159448in}}%
\pgfpathlineto{\pgfqpoint{1.006049in}{1.141226in}}%
\pgfpathlineto{\pgfqpoint{1.003212in}{1.123158in}}%
\pgfpathlineto{\pgfqpoint{1.001224in}{1.103107in}}%
\pgfpathlineto{\pgfqpoint{1.000426in}{1.083434in}}%
\pgfpathlineto{\pgfqpoint{1.000738in}{1.064213in}}%
\pgfpathlineto{\pgfqpoint{1.002090in}{1.045508in}}%
\pgfpathlineto{\pgfqpoint{1.004413in}{1.027370in}}%
\pgfpathlineto{\pgfqpoint{1.008061in}{1.007930in}}%
\pgfpathlineto{\pgfqpoint{1.012754in}{0.989284in}}%
\pgfpathlineto{\pgfqpoint{1.018418in}{0.971466in}}%
\pgfpathlineto{\pgfqpoint{1.024983in}{0.954499in}}%
\pgfpathlineto{\pgfqpoint{1.032378in}{0.938397in}}%
\pgfpathlineto{\pgfqpoint{1.040537in}{0.923168in}}%
\pgfpathlineto{\pgfqpoint{1.049392in}{0.908815in}}%
\pgfpathlineto{\pgfqpoint{1.059857in}{0.894033in}}%
\pgfpathlineto{\pgfqpoint{1.070994in}{0.880297in}}%
\pgfpathlineto{\pgfqpoint{1.082714in}{0.867590in}}%
\pgfpathlineto{\pgfqpoint{1.094925in}{0.855895in}}%
\pgfpathlineto{\pgfqpoint{1.108697in}{0.844262in}}%
\pgfpathlineto{\pgfqpoint{1.122824in}{0.833771in}}%
\pgfpathlineto{\pgfqpoint{1.138389in}{0.823653in}}%
\pgfpathlineto{\pgfqpoint{1.154078in}{0.814784in}}%
\pgfpathlineto{\pgfqpoint{1.170942in}{0.806569in}}%
\pgfpathlineto{\pgfqpoint{1.187598in}{0.799675in}}%
\pgfpathlineto{\pgfqpoint{1.205022in}{0.793675in}}%
\pgfpathlineto{\pgfqpoint{1.221823in}{0.789024in}}%
\pgfpathlineto{\pgfqpoint{1.238863in}{0.785450in}}%
\pgfpathlineto{\pgfqpoint{1.254819in}{0.783203in}}%
\pgfpathlineto{\pgfqpoint{1.270421in}{0.782147in}}%
\pgfpathlineto{\pgfqpoint{1.284499in}{0.782334in}}%
\pgfpathlineto{\pgfqpoint{1.296964in}{0.783634in}}%
\pgfpathlineto{\pgfqpoint{1.307177in}{0.785759in}}%
\pgfpathlineto{\pgfqpoint{1.315908in}{0.788647in}}%
\pgfpathlineto{\pgfqpoint{1.323177in}{0.792194in}}%
\pgfpathlineto{\pgfqpoint{1.329025in}{0.796295in}}%
\pgfpathlineto{\pgfqpoint{1.333517in}{0.800851in}}%
\pgfpathlineto{\pgfqpoint{1.336732in}{0.805764in}}%
\pgfpathlineto{\pgfqpoint{1.338763in}{0.810941in}}%
\pgfpathlineto{\pgfqpoint{1.339738in}{0.816631in}}%
\pgfpathlineto{\pgfqpoint{1.339581in}{0.822762in}}%
\pgfpathlineto{\pgfqpoint{1.338262in}{0.829233in}}%
\pgfpathlineto{\pgfqpoint{1.335643in}{0.836249in}}%
\pgfpathlineto{\pgfqpoint{1.331656in}{0.843616in}}%
\pgfpathlineto{\pgfqpoint{1.326340in}{0.851080in}}%
\pgfpathlineto{\pgfqpoint{1.319594in}{0.858597in}}%
\pgfpathlineto{\pgfqpoint{1.311912in}{0.865511in}}%
\pgfpathlineto{\pgfqpoint{1.303437in}{0.871658in}}%
\pgfpathlineto{\pgfqpoint{1.294914in}{0.876498in}}%
\pgfpathlineto{\pgfqpoint{1.287050in}{0.879776in}}%
\pgfpathlineto{\pgfqpoint{1.287050in}{0.879776in}}%
\pgfusepath{stroke}%
\end{pgfscope}%
\begin{pgfscope}%
\pgfsetrectcap%
\pgfsetroundjoin%
\pgfsetlinewidth{1.505625pt}%
\definecolor{currentstroke}{rgb}{0.172549,0.627451,0.172549}%
\pgfsetstrokecolor{currentstroke}%
\pgfsetstrokeopacity{0.200000}%
\pgfsetdash{}{0pt}%
\pgfpathmoveto{\pgfqpoint{2.812911in}{0.514214in}}%
\pgfpathlineto{\pgfqpoint{2.721198in}{0.528370in}}%
\pgfpathlineto{\pgfqpoint{2.638829in}{0.542475in}}%
\pgfpathlineto{\pgfqpoint{2.564340in}{0.556523in}}%
\pgfpathlineto{\pgfqpoint{2.496552in}{0.570505in}}%
\pgfpathlineto{\pgfqpoint{2.434510in}{0.584412in}}%
\pgfpathlineto{\pgfqpoint{2.377427in}{0.598237in}}%
\pgfpathlineto{\pgfqpoint{2.324656in}{0.611968in}}%
\pgfpathlineto{\pgfqpoint{2.263937in}{0.628989in}}%
\pgfpathlineto{\pgfqpoint{2.208241in}{0.645832in}}%
\pgfpathlineto{\pgfqpoint{2.156860in}{0.662479in}}%
\pgfpathlineto{\pgfqpoint{2.100090in}{0.682169in}}%
\pgfpathlineto{\pgfqpoint{2.047889in}{0.701515in}}%
\pgfpathlineto{\pgfqpoint{1.991899in}{0.723608in}}%
\pgfpathlineto{\pgfqpoint{1.940415in}{0.745137in}}%
\pgfpathlineto{\pgfqpoint{1.886265in}{0.768988in}}%
\pgfpathlineto{\pgfqpoint{1.830421in}{0.794775in}}%
\pgfpathlineto{\pgfqpoint{1.768306in}{0.824659in}}%
\pgfpathlineto{\pgfqpoint{1.683203in}{0.866992in}}%
\pgfpathlineto{\pgfqpoint{1.561875in}{0.927423in}}%
\pgfpathlineto{\pgfqpoint{1.512025in}{0.951025in}}%
\pgfpathlineto{\pgfqpoint{1.470087in}{0.969676in}}%
\pgfpathlineto{\pgfqpoint{1.434779in}{0.984155in}}%
\pgfpathlineto{\pgfqpoint{1.405084in}{0.995189in}}%
\pgfpathlineto{\pgfqpoint{1.377784in}{1.004171in}}%
\pgfpathlineto{\pgfqpoint{1.352743in}{1.011203in}}%
\pgfpathlineto{\pgfqpoint{1.331834in}{1.016000in}}%
\pgfpathlineto{\pgfqpoint{1.312602in}{1.019375in}}%
\pgfpathlineto{\pgfqpoint{1.294964in}{1.021426in}}%
\pgfpathlineto{\pgfqpoint{1.278844in}{1.022256in}}%
\pgfpathlineto{\pgfqpoint{1.264166in}{1.021968in}}%
\pgfpathlineto{\pgfqpoint{1.250859in}{1.020665in}}%
\pgfpathlineto{\pgfqpoint{1.238854in}{1.018449in}}%
\pgfpathlineto{\pgfqpoint{1.228085in}{1.015421in}}%
\pgfpathlineto{\pgfqpoint{1.218489in}{1.011675in}}%
\pgfpathlineto{\pgfqpoint{1.209215in}{1.006836in}}%
\pgfpathlineto{\pgfqpoint{1.201208in}{1.001358in}}%
\pgfpathlineto{\pgfqpoint{1.194391in}{0.995351in}}%
\pgfpathlineto{\pgfqpoint{1.188693in}{0.988916in}}%
\pgfpathlineto{\pgfqpoint{1.183671in}{0.981519in}}%
\pgfpathlineto{\pgfqpoint{1.179811in}{0.973838in}}%
\pgfpathlineto{\pgfqpoint{1.176843in}{0.965313in}}%
\pgfpathlineto{\pgfqpoint{1.175038in}{0.956687in}}%
\pgfpathlineto{\pgfqpoint{1.174282in}{0.947401in}}%
\pgfpathlineto{\pgfqpoint{1.174642in}{0.938225in}}%
\pgfpathlineto{\pgfqpoint{1.176138in}{0.928622in}}%
\pgfpathlineto{\pgfqpoint{1.178849in}{0.918742in}}%
\pgfpathlineto{\pgfqpoint{1.182559in}{0.909331in}}%
\pgfpathlineto{\pgfqpoint{1.187424in}{0.899926in}}%
\pgfpathlineto{\pgfqpoint{1.193422in}{0.890705in}}%
\pgfpathlineto{\pgfqpoint{1.200486in}{0.881853in}}%
\pgfpathlineto{\pgfqpoint{1.208498in}{0.873551in}}%
\pgfpathlineto{\pgfqpoint{1.217717in}{0.865639in}}%
\pgfpathlineto{\pgfqpoint{1.227508in}{0.858721in}}%
\pgfpathlineto{\pgfqpoint{1.237581in}{0.852928in}}%
\pgfpathlineto{\pgfqpoint{1.248009in}{0.848204in}}%
\pgfpathlineto{\pgfqpoint{1.258006in}{0.844868in}}%
\pgfpathlineto{\pgfqpoint{1.267225in}{0.842914in}}%
\pgfpathlineto{\pgfqpoint{1.275362in}{0.842285in}}%
\pgfpathlineto{\pgfqpoint{1.281969in}{0.842835in}}%
\pgfpathlineto{\pgfqpoint{1.287042in}{0.844321in}}%
\pgfpathlineto{\pgfqpoint{1.290649in}{0.846507in}}%
\pgfpathlineto{\pgfqpoint{1.292854in}{0.849077in}}%
\pgfpathlineto{\pgfqpoint{1.293970in}{0.852019in}}%
\pgfpathlineto{\pgfqpoint{1.294041in}{0.855391in}}%
\pgfpathlineto{\pgfqpoint{1.292975in}{0.859105in}}%
\pgfpathlineto{\pgfqpoint{1.290662in}{0.863059in}}%
\pgfpathlineto{\pgfqpoint{1.289441in}{0.864574in}}%
\pgfpathlineto{\pgfqpoint{1.289441in}{0.864574in}}%
\pgfusepath{stroke}%
\end{pgfscope}%
\begin{pgfscope}%
\pgfsetbuttcap%
\pgfsetroundjoin%
\pgfsetlinewidth{0.702625pt}%
\definecolor{currentstroke}{rgb}{1.000000,0.498039,0.054902}%
\pgfsetstrokecolor{currentstroke}%
\pgfsetdash{{2.590000pt}{1.120000pt}}{0.000000pt}%
\pgfpathmoveto{\pgfqpoint{0.506574in}{0.278150in}}%
\pgfpathlineto{\pgfqpoint{2.966667in}{0.278150in}}%
\pgfusepath{stroke}%
\end{pgfscope}%
\begin{pgfscope}%
\definecolor{textcolor}{rgb}{0.000000,0.000000,0.000000}%
\pgfsetstrokecolor{textcolor}%
\pgfsetfillcolor{textcolor}%
\pgftext[x=2.966667in,y=0.180927in,right,top]{\color{textcolor}\rmfamily\fontsize{12.000000}{14.400000}\selectfont \(\displaystyle u\)}%
\end{pgfscope}%
\begin{pgfscope}%
\definecolor{textcolor}{rgb}{0.000000,0.000000,0.000000}%
\pgfsetstrokecolor{textcolor}%
\pgfsetfillcolor{textcolor}%
\pgftext[x=0.409351in,y=2.166667in,right,top]{\color{textcolor}\rmfamily\fontsize{12.000000}{14.400000}\selectfont \(\displaystyle v\)}%
\end{pgfscope}%
\begin{pgfscope}%
\pgfsetroundcap%
\pgfsetroundjoin%
\pgfsetlinewidth{1.003750pt}%
\definecolor{currentstroke}{rgb}{0.172549,0.627451,0.172549}%
\pgfsetstrokecolor{currentstroke}%
\pgfsetdash{}{0pt}%
\pgfpathmoveto{\pgfqpoint{0.732539in}{0.281442in}}%
\pgfpathquadraticcurveto{\pgfqpoint{0.733989in}{0.281433in}}{\pgfqpoint{0.719912in}{0.281517in}}%
\pgfusepath{stroke}%
\end{pgfscope}%
\begin{pgfscope}%
\pgfsetroundcap%
\pgfsetroundjoin%
\definecolor{currentfill}{rgb}{0.172549,0.627451,0.172549}%
\pgfsetfillcolor{currentfill}%
\pgfsetlinewidth{1.003750pt}%
\definecolor{currentstroke}{rgb}{0.172549,0.627451,0.172549}%
\pgfsetstrokecolor{currentstroke}%
\pgfsetdash{}{0pt}%
\pgfpathmoveto{\pgfqpoint{0.620213in}{0.332115in}}%
\pgfpathlineto{\pgfqpoint{0.719912in}{0.281517in}}%
\pgfpathlineto{\pgfqpoint{0.619614in}{0.232117in}}%
\pgfpathlineto{\pgfqpoint{0.620213in}{0.332115in}}%
\pgfpathclose%
\pgfusepath{stroke,fill}%
\end{pgfscope}%
\begin{pgfscope}%
\pgfsetroundcap%
\pgfsetroundjoin%
\pgfsetlinewidth{1.003750pt}%
\definecolor{currentstroke}{rgb}{0.172549,0.627451,0.172549}%
\pgfsetstrokecolor{currentstroke}%
\pgfsetdash{}{0pt}%
\pgfpathmoveto{\pgfqpoint{1.605949in}{0.733561in}}%
\pgfpathquadraticcurveto{\pgfqpoint{1.605117in}{0.734296in}}{\pgfqpoint{1.615928in}{0.724756in}}%
\pgfusepath{stroke}%
\end{pgfscope}%
\begin{pgfscope}%
\pgfsetroundcap%
\pgfsetroundjoin%
\definecolor{currentfill}{rgb}{0.172549,0.627451,0.172549}%
\pgfsetfillcolor{currentfill}%
\pgfsetlinewidth{1.003750pt}%
\definecolor{currentstroke}{rgb}{0.172549,0.627451,0.172549}%
\pgfsetstrokecolor{currentstroke}%
\pgfsetdash{}{0pt}%
\pgfpathmoveto{\pgfqpoint{1.657832in}{0.621103in}}%
\pgfpathlineto{\pgfqpoint{1.615928in}{0.724756in}}%
\pgfpathlineto{\pgfqpoint{1.723993in}{0.696087in}}%
\pgfpathlineto{\pgfqpoint{1.657832in}{0.621103in}}%
\pgfpathclose%
\pgfusepath{stroke,fill}%
\end{pgfscope}%
\begin{pgfscope}%
\pgfsetroundcap%
\pgfsetroundjoin%
\pgfsetlinewidth{1.003750pt}%
\definecolor{currentstroke}{rgb}{0.172549,0.627451,0.172549}%
\pgfsetstrokecolor{currentstroke}%
\pgfsetstrokeopacity{0.200000}%
\pgfsetdash{}{0pt}%
\pgfpathmoveto{\pgfqpoint{0.693331in}{0.773262in}}%
\pgfpathquadraticcurveto{\pgfqpoint{0.693978in}{0.771860in}}{\pgfqpoint{0.688117in}{0.784557in}}%
\pgfusepath{stroke}%
\end{pgfscope}%
\begin{pgfscope}%
\pgfsetroundcap%
\pgfsetroundjoin%
\definecolor{currentfill}{rgb}{0.172549,0.627451,0.172549}%
\pgfsetfillcolor{currentfill}%
\pgfsetfillopacity{0.200000}%
\pgfsetlinewidth{1.003750pt}%
\definecolor{currentstroke}{rgb}{0.172549,0.627451,0.172549}%
\pgfsetstrokecolor{currentstroke}%
\pgfsetstrokeopacity{0.200000}%
\pgfsetdash{}{0pt}%
\pgfpathmoveto{\pgfqpoint{0.691602in}{0.896306in}}%
\pgfpathlineto{\pgfqpoint{0.688117in}{0.784557in}}%
\pgfpathlineto{\pgfqpoint{0.600809in}{0.854395in}}%
\pgfpathlineto{\pgfqpoint{0.691602in}{0.896306in}}%
\pgfpathclose%
\pgfusepath{stroke,fill}%
\end{pgfscope}%
\begin{pgfscope}%
\pgfsetroundcap%
\pgfsetroundjoin%
\pgfsetlinewidth{1.003750pt}%
\definecolor{currentstroke}{rgb}{0.172549,0.627451,0.172549}%
\pgfsetstrokecolor{currentstroke}%
\pgfsetstrokeopacity{0.200000}%
\pgfsetdash{}{0pt}%
\pgfpathmoveto{\pgfqpoint{0.659364in}{1.263447in}}%
\pgfpathquadraticcurveto{\pgfqpoint{0.659442in}{1.260493in}}{\pgfqpoint{0.659108in}{1.273061in}}%
\pgfusepath{stroke}%
\end{pgfscope}%
\begin{pgfscope}%
\pgfsetroundcap%
\pgfsetroundjoin%
\definecolor{currentfill}{rgb}{0.172549,0.627451,0.172549}%
\pgfsetfillcolor{currentfill}%
\pgfsetfillopacity{0.200000}%
\pgfsetlinewidth{1.003750pt}%
\definecolor{currentstroke}{rgb}{0.172549,0.627451,0.172549}%
\pgfsetstrokecolor{currentstroke}%
\pgfsetstrokeopacity{0.200000}%
\pgfsetdash{}{0pt}%
\pgfpathmoveto{\pgfqpoint{0.706438in}{1.374352in}}%
\pgfpathlineto{\pgfqpoint{0.659108in}{1.273061in}}%
\pgfpathlineto{\pgfqpoint{0.606473in}{1.371700in}}%
\pgfpathlineto{\pgfqpoint{0.706438in}{1.374352in}}%
\pgfpathclose%
\pgfusepath{stroke,fill}%
\end{pgfscope}%
\begin{pgfscope}%
\pgfsetroundcap%
\pgfsetroundjoin%
\pgfsetlinewidth{1.003750pt}%
\definecolor{currentstroke}{rgb}{0.172549,0.627451,0.172549}%
\pgfsetstrokecolor{currentstroke}%
\pgfsetstrokeopacity{0.200000}%
\pgfsetdash{}{0pt}%
\pgfpathmoveto{\pgfqpoint{1.789006in}{1.839130in}}%
\pgfpathquadraticcurveto{\pgfqpoint{1.778004in}{1.842970in}}{\pgfqpoint{1.781663in}{1.841693in}}%
\pgfusepath{stroke}%
\end{pgfscope}%
\begin{pgfscope}%
\pgfsetroundcap%
\pgfsetroundjoin%
\definecolor{currentfill}{rgb}{0.172549,0.627451,0.172549}%
\pgfsetfillcolor{currentfill}%
\pgfsetfillopacity{0.200000}%
\pgfsetlinewidth{1.003750pt}%
\definecolor{currentstroke}{rgb}{0.172549,0.627451,0.172549}%
\pgfsetstrokecolor{currentstroke}%
\pgfsetstrokeopacity{0.200000}%
\pgfsetdash{}{0pt}%
\pgfpathmoveto{\pgfqpoint{1.859600in}{1.761532in}}%
\pgfpathlineto{\pgfqpoint{1.781663in}{1.841693in}}%
\pgfpathlineto{\pgfqpoint{1.892554in}{1.855946in}}%
\pgfpathlineto{\pgfqpoint{1.859600in}{1.761532in}}%
\pgfpathclose%
\pgfusepath{stroke,fill}%
\end{pgfscope}%
\begin{pgfscope}%
\pgfsetroundcap%
\pgfsetroundjoin%
\pgfsetlinewidth{1.003750pt}%
\definecolor{currentstroke}{rgb}{0.172549,0.627451,0.172549}%
\pgfsetstrokecolor{currentstroke}%
\pgfsetstrokeopacity{0.200000}%
\pgfsetdash{}{0pt}%
\pgfpathmoveto{\pgfqpoint{2.086571in}{1.100438in}}%
\pgfpathquadraticcurveto{\pgfqpoint{2.078059in}{1.103673in}}{\pgfqpoint{2.084062in}{1.101391in}}%
\pgfusepath{stroke}%
\end{pgfscope}%
\begin{pgfscope}%
\pgfsetroundcap%
\pgfsetroundjoin%
\definecolor{currentfill}{rgb}{0.172549,0.627451,0.172549}%
\pgfsetfillcolor{currentfill}%
\pgfsetfillopacity{0.200000}%
\pgfsetlinewidth{1.003750pt}%
\definecolor{currentstroke}{rgb}{0.172549,0.627451,0.172549}%
\pgfsetstrokecolor{currentstroke}%
\pgfsetstrokeopacity{0.200000}%
\pgfsetdash{}{0pt}%
\pgfpathmoveto{\pgfqpoint{2.159771in}{1.019122in}}%
\pgfpathlineto{\pgfqpoint{2.084062in}{1.101391in}}%
\pgfpathlineto{\pgfqpoint{2.195303in}{1.112597in}}%
\pgfpathlineto{\pgfqpoint{2.159771in}{1.019122in}}%
\pgfpathclose%
\pgfusepath{stroke,fill}%
\end{pgfscope}%
\begin{pgfscope}%
\pgfsetroundcap%
\pgfsetroundjoin%
\pgfsetlinewidth{1.003750pt}%
\definecolor{currentstroke}{rgb}{0.172549,0.627451,0.172549}%
\pgfsetstrokecolor{currentstroke}%
\pgfsetstrokeopacity{0.200000}%
\pgfsetdash{}{0pt}%
\pgfpathmoveto{\pgfqpoint{2.299714in}{0.618797in}}%
\pgfpathquadraticcurveto{\pgfqpoint{2.293645in}{0.620499in}}{\pgfqpoint{2.302528in}{0.618007in}}%
\pgfusepath{stroke}%
\end{pgfscope}%
\begin{pgfscope}%
\pgfsetroundcap%
\pgfsetroundjoin%
\definecolor{currentfill}{rgb}{0.172549,0.627451,0.172549}%
\pgfsetfillcolor{currentfill}%
\pgfsetfillopacity{0.200000}%
\pgfsetlinewidth{1.003750pt}%
\definecolor{currentstroke}{rgb}{0.172549,0.627451,0.172549}%
\pgfsetstrokecolor{currentstroke}%
\pgfsetstrokeopacity{0.200000}%
\pgfsetdash{}{0pt}%
\pgfpathmoveto{\pgfqpoint{2.385309in}{0.542858in}}%
\pgfpathlineto{\pgfqpoint{2.302528in}{0.618007in}}%
\pgfpathlineto{\pgfqpoint{2.412316in}{0.639143in}}%
\pgfpathlineto{\pgfqpoint{2.385309in}{0.542858in}}%
\pgfpathclose%
\pgfusepath{stroke,fill}%
\end{pgfscope}%
\begin{pgfscope}%
\pgfsetbuttcap%
\pgfsetroundjoin%
\pgfsetlinewidth{0.702625pt}%
\definecolor{currentstroke}{rgb}{0.121569,0.466667,0.705882}%
\pgfsetstrokecolor{currentstroke}%
\pgfsetdash{{2.590000pt}{1.120000pt}}{0.000000pt}%
\pgfpathmoveto{\pgfqpoint{0.506574in}{0.278150in}}%
\pgfpathlineto{\pgfqpoint{0.506574in}{1.458473in}}%
\pgfusepath{stroke}%
\end{pgfscope}%
\begin{pgfscope}%
\pgfsetbuttcap%
\pgfsetroundjoin%
\pgfsetlinewidth{0.702625pt}%
\definecolor{currentstroke}{rgb}{1.000000,0.498039,0.054902}%
\pgfsetstrokecolor{currentstroke}%
\pgfsetdash{{2.590000pt}{1.120000pt}}{0.000000pt}%
\pgfpathmoveto{\pgfqpoint{1.275353in}{0.278150in}}%
\pgfpathlineto{\pgfqpoint{1.275353in}{2.166667in}}%
\pgfusepath{stroke}%
\end{pgfscope}%
\begin{pgfscope}%
\pgfsetbuttcap%
\pgfsetroundjoin%
\definecolor{currentfill}{rgb}{1.000000,1.000000,1.000000}%
\pgfsetfillcolor{currentfill}%
\pgfsetlinewidth{1.003750pt}%
\definecolor{currentstroke}{rgb}{0.000000,0.000000,0.000000}%
\pgfsetstrokecolor{currentstroke}%
\pgfsetdash{}{0pt}%
\pgfsys@defobject{currentmarker}{\pgfqpoint{-0.041667in}{-0.041667in}}{\pgfqpoint{0.041667in}{0.041667in}}{%
\pgfpathmoveto{\pgfqpoint{0.000000in}{-0.041667in}}%
\pgfpathcurveto{\pgfqpoint{0.011050in}{-0.041667in}}{\pgfqpoint{0.021649in}{-0.037276in}}{\pgfqpoint{0.029463in}{-0.029463in}}%
\pgfpathcurveto{\pgfqpoint{0.037276in}{-0.021649in}}{\pgfqpoint{0.041667in}{-0.011050in}}{\pgfqpoint{0.041667in}{0.000000in}}%
\pgfpathcurveto{\pgfqpoint{0.041667in}{0.011050in}}{\pgfqpoint{0.037276in}{0.021649in}}{\pgfqpoint{0.029463in}{0.029463in}}%
\pgfpathcurveto{\pgfqpoint{0.021649in}{0.037276in}}{\pgfqpoint{0.011050in}{0.041667in}}{\pgfqpoint{0.000000in}{0.041667in}}%
\pgfpathcurveto{\pgfqpoint{-0.011050in}{0.041667in}}{\pgfqpoint{-0.021649in}{0.037276in}}{\pgfqpoint{-0.029463in}{0.029463in}}%
\pgfpathcurveto{\pgfqpoint{-0.037276in}{0.021649in}}{\pgfqpoint{-0.041667in}{0.011050in}}{\pgfqpoint{-0.041667in}{0.000000in}}%
\pgfpathcurveto{\pgfqpoint{-0.041667in}{-0.011050in}}{\pgfqpoint{-0.037276in}{-0.021649in}}{\pgfqpoint{-0.029463in}{-0.029463in}}%
\pgfpathcurveto{\pgfqpoint{-0.021649in}{-0.037276in}}{\pgfqpoint{-0.011050in}{-0.041667in}}{\pgfqpoint{0.000000in}{-0.041667in}}%
\pgfpathclose%
\pgfusepath{stroke,fill}%
}%
\begin{pgfscope}%
\pgfsys@transformshift{0.506574in}{0.278150in}%
\pgfsys@useobject{currentmarker}{}%
\end{pgfscope}%
\end{pgfscope}%
\begin{pgfscope}%
\pgfsetbuttcap%
\pgfsetroundjoin%
\definecolor{currentfill}{rgb}{1.000000,1.000000,1.000000}%
\pgfsetfillcolor{currentfill}%
\pgfsetlinewidth{1.003750pt}%
\definecolor{currentstroke}{rgb}{0.000000,0.000000,0.000000}%
\pgfsetstrokecolor{currentstroke}%
\pgfsetdash{}{0pt}%
\pgfsys@defobject{currentmarker}{\pgfqpoint{-0.041667in}{-0.041667in}}{\pgfqpoint{0.041667in}{0.041667in}}{%
\pgfpathmoveto{\pgfqpoint{0.000000in}{-0.041667in}}%
\pgfpathcurveto{\pgfqpoint{0.011050in}{-0.041667in}}{\pgfqpoint{0.021649in}{-0.037276in}}{\pgfqpoint{0.029463in}{-0.029463in}}%
\pgfpathcurveto{\pgfqpoint{0.037276in}{-0.021649in}}{\pgfqpoint{0.041667in}{-0.011050in}}{\pgfqpoint{0.041667in}{0.000000in}}%
\pgfpathcurveto{\pgfqpoint{0.041667in}{0.011050in}}{\pgfqpoint{0.037276in}{0.021649in}}{\pgfqpoint{0.029463in}{0.029463in}}%
\pgfpathcurveto{\pgfqpoint{0.021649in}{0.037276in}}{\pgfqpoint{0.011050in}{0.041667in}}{\pgfqpoint{0.000000in}{0.041667in}}%
\pgfpathcurveto{\pgfqpoint{-0.011050in}{0.041667in}}{\pgfqpoint{-0.021649in}{0.037276in}}{\pgfqpoint{-0.029463in}{0.029463in}}%
\pgfpathcurveto{\pgfqpoint{-0.037276in}{0.021649in}}{\pgfqpoint{-0.041667in}{0.011050in}}{\pgfqpoint{-0.041667in}{0.000000in}}%
\pgfpathcurveto{\pgfqpoint{-0.041667in}{-0.011050in}}{\pgfqpoint{-0.037276in}{-0.021649in}}{\pgfqpoint{-0.029463in}{-0.029463in}}%
\pgfpathcurveto{\pgfqpoint{-0.021649in}{-0.037276in}}{\pgfqpoint{-0.011050in}{-0.041667in}}{\pgfqpoint{0.000000in}{-0.041667in}}%
\pgfpathclose%
\pgfusepath{stroke,fill}%
}%
\begin{pgfscope}%
\pgfsys@transformshift{2.044132in}{0.278150in}%
\pgfsys@useobject{currentmarker}{}%
\end{pgfscope}%
\end{pgfscope}%
\begin{pgfscope}%
\pgfsetbuttcap%
\pgfsetroundjoin%
\definecolor{currentfill}{rgb}{0.000000,0.000000,0.000000}%
\pgfsetfillcolor{currentfill}%
\pgfsetlinewidth{1.003750pt}%
\definecolor{currentstroke}{rgb}{0.000000,0.000000,0.000000}%
\pgfsetstrokecolor{currentstroke}%
\pgfsetdash{}{0pt}%
\pgfsys@defobject{currentmarker}{\pgfqpoint{-0.041667in}{-0.041667in}}{\pgfqpoint{0.041667in}{0.041667in}}{%
\pgfpathmoveto{\pgfqpoint{0.000000in}{-0.041667in}}%
\pgfpathcurveto{\pgfqpoint{0.011050in}{-0.041667in}}{\pgfqpoint{0.021649in}{-0.037276in}}{\pgfqpoint{0.029463in}{-0.029463in}}%
\pgfpathcurveto{\pgfqpoint{0.037276in}{-0.021649in}}{\pgfqpoint{0.041667in}{-0.011050in}}{\pgfqpoint{0.041667in}{0.000000in}}%
\pgfpathcurveto{\pgfqpoint{0.041667in}{0.011050in}}{\pgfqpoint{0.037276in}{0.021649in}}{\pgfqpoint{0.029463in}{0.029463in}}%
\pgfpathcurveto{\pgfqpoint{0.021649in}{0.037276in}}{\pgfqpoint{0.011050in}{0.041667in}}{\pgfqpoint{0.000000in}{0.041667in}}%
\pgfpathcurveto{\pgfqpoint{-0.011050in}{0.041667in}}{\pgfqpoint{-0.021649in}{0.037276in}}{\pgfqpoint{-0.029463in}{0.029463in}}%
\pgfpathcurveto{\pgfqpoint{-0.037276in}{0.021649in}}{\pgfqpoint{-0.041667in}{0.011050in}}{\pgfqpoint{-0.041667in}{0.000000in}}%
\pgfpathcurveto{\pgfqpoint{-0.041667in}{-0.011050in}}{\pgfqpoint{-0.037276in}{-0.021649in}}{\pgfqpoint{-0.029463in}{-0.029463in}}%
\pgfpathcurveto{\pgfqpoint{-0.021649in}{-0.037276in}}{\pgfqpoint{-0.011050in}{-0.041667in}}{\pgfqpoint{0.000000in}{-0.041667in}}%
\pgfpathclose%
\pgfusepath{stroke,fill}%
}%
\begin{pgfscope}%
\pgfsys@transformshift{1.275353in}{0.868311in}%
\pgfsys@useobject{currentmarker}{}%
\end{pgfscope}%
\end{pgfscope}%
\end{pgfpicture}%
\makeatother%
\endgroup%

            %% Creator: Matplotlib, PGF backend
%%
%% To include the figure in your LaTeX document, write
%%   \input{<filename>.pgf}
%%
%% Make sure the required packages are loaded in your preamble
%%   \usepackage{pgf}
%%
%% and, on pdftex
%%   \usepackage[utf8]{inputenc}\DeclareUnicodeCharacter{2212}{-}
%%
%% or, on luatex and xetex
%%   \usepackage{unicode-math}
%%
%% Figures using additional raster images can only be included by \input if
%% they are in the same directory as the main LaTeX file. For loading figures
%% from other directories you can use the `import` package
%%   \usepackage{import}
%%
%% and then include the figures with
%%   \import{<path to file>}{<filename>.pgf}
%%
%% Matplotlib used the following preamble
%%   \usepackage{fontspec}
%%   \setmainfont{DejaVuSerif.ttf}[Path=/Users/adam/opt/anaconda3/lib/python3.8/site-packages/matplotlib/mpl-data/fonts/ttf/]
%%   \setsansfont{DejaVuSans.ttf}[Path=/Users/adam/opt/anaconda3/lib/python3.8/site-packages/matplotlib/mpl-data/fonts/ttf/]
%%   \setmonofont{DejaVuSansMono.ttf}[Path=/Users/adam/opt/anaconda3/lib/python3.8/site-packages/matplotlib/mpl-data/fonts/ttf/]
%%
\begingroup%
\makeatletter%
\begin{pgfpicture}%
\pgfpathrectangle{\pgfpointorigin}{\pgfqpoint{3.000000in}{2.200000in}}%
\pgfusepath{use as bounding box, clip}%
\begin{pgfscope}%
\pgfsetbuttcap%
\pgfsetmiterjoin%
\definecolor{currentfill}{rgb}{1.000000,1.000000,1.000000}%
\pgfsetfillcolor{currentfill}%
\pgfsetlinewidth{0.000000pt}%
\definecolor{currentstroke}{rgb}{1.000000,1.000000,1.000000}%
\pgfsetstrokecolor{currentstroke}%
\pgfsetdash{}{0pt}%
\pgfpathmoveto{\pgfqpoint{0.000000in}{0.000000in}}%
\pgfpathlineto{\pgfqpoint{3.000000in}{0.000000in}}%
\pgfpathlineto{\pgfqpoint{3.000000in}{2.200000in}}%
\pgfpathlineto{\pgfqpoint{0.000000in}{2.200000in}}%
\pgfpathclose%
\pgfusepath{fill}%
\end{pgfscope}%
\begin{pgfscope}%
\pgfsetbuttcap%
\pgfsetmiterjoin%
\definecolor{currentfill}{rgb}{1.000000,1.000000,1.000000}%
\pgfsetfillcolor{currentfill}%
\pgfsetlinewidth{0.000000pt}%
\definecolor{currentstroke}{rgb}{0.000000,0.000000,0.000000}%
\pgfsetstrokecolor{currentstroke}%
\pgfsetstrokeopacity{0.000000}%
\pgfsetdash{}{0pt}%
\pgfpathmoveto{\pgfqpoint{0.506574in}{0.278150in}}%
\pgfpathlineto{\pgfqpoint{2.966667in}{0.278150in}}%
\pgfpathlineto{\pgfqpoint{2.966667in}{2.166667in}}%
\pgfpathlineto{\pgfqpoint{0.506574in}{2.166667in}}%
\pgfpathclose%
\pgfusepath{fill}%
\end{pgfscope}%
\begin{pgfscope}%
\pgfsetbuttcap%
\pgfsetroundjoin%
\definecolor{currentfill}{rgb}{0.000000,0.000000,0.000000}%
\pgfsetfillcolor{currentfill}%
\pgfsetlinewidth{0.803000pt}%
\definecolor{currentstroke}{rgb}{0.000000,0.000000,0.000000}%
\pgfsetstrokecolor{currentstroke}%
\pgfsetdash{}{0pt}%
\pgfsys@defobject{currentmarker}{\pgfqpoint{0.000000in}{-0.048611in}}{\pgfqpoint{0.000000in}{0.000000in}}{%
\pgfpathmoveto{\pgfqpoint{0.000000in}{0.000000in}}%
\pgfpathlineto{\pgfqpoint{0.000000in}{-0.048611in}}%
\pgfusepath{stroke,fill}%
}%
\begin{pgfscope}%
\pgfsys@transformshift{0.506574in}{0.278150in}%
\pgfsys@useobject{currentmarker}{}%
\end{pgfscope}%
\end{pgfscope}%
\begin{pgfscope}%
\definecolor{textcolor}{rgb}{0.000000,0.000000,0.000000}%
\pgfsetstrokecolor{textcolor}%
\pgfsetfillcolor{textcolor}%
\pgftext[x=0.506574in,y=0.180927in,,top]{\color{textcolor}\rmfamily\fontsize{12.000000}{14.400000}\selectfont 0}%
\end{pgfscope}%
\begin{pgfscope}%
\pgfsetbuttcap%
\pgfsetroundjoin%
\definecolor{currentfill}{rgb}{0.000000,0.000000,0.000000}%
\pgfsetfillcolor{currentfill}%
\pgfsetlinewidth{0.803000pt}%
\definecolor{currentstroke}{rgb}{0.000000,0.000000,0.000000}%
\pgfsetstrokecolor{currentstroke}%
\pgfsetdash{}{0pt}%
\pgfsys@defobject{currentmarker}{\pgfqpoint{0.000000in}{-0.048611in}}{\pgfqpoint{0.000000in}{0.000000in}}{%
\pgfpathmoveto{\pgfqpoint{0.000000in}{0.000000in}}%
\pgfpathlineto{\pgfqpoint{0.000000in}{-0.048611in}}%
\pgfusepath{stroke,fill}%
}%
\begin{pgfscope}%
\pgfsys@transformshift{1.275353in}{0.278150in}%
\pgfsys@useobject{currentmarker}{}%
\end{pgfscope}%
\end{pgfscope}%
\begin{pgfscope}%
\definecolor{textcolor}{rgb}{0.000000,0.000000,0.000000}%
\pgfsetstrokecolor{textcolor}%
\pgfsetfillcolor{textcolor}%
\pgftext[x=1.275353in,y=0.180927in,,top]{\color{textcolor}\rmfamily\fontsize{12.000000}{14.400000}\selectfont \(\displaystyle \beta\)}%
\end{pgfscope}%
\begin{pgfscope}%
\pgfsetbuttcap%
\pgfsetroundjoin%
\definecolor{currentfill}{rgb}{0.000000,0.000000,0.000000}%
\pgfsetfillcolor{currentfill}%
\pgfsetlinewidth{0.803000pt}%
\definecolor{currentstroke}{rgb}{0.000000,0.000000,0.000000}%
\pgfsetstrokecolor{currentstroke}%
\pgfsetdash{}{0pt}%
\pgfsys@defobject{currentmarker}{\pgfqpoint{0.000000in}{-0.048611in}}{\pgfqpoint{0.000000in}{0.000000in}}{%
\pgfpathmoveto{\pgfqpoint{0.000000in}{0.000000in}}%
\pgfpathlineto{\pgfqpoint{0.000000in}{-0.048611in}}%
\pgfusepath{stroke,fill}%
}%
\begin{pgfscope}%
\pgfsys@transformshift{2.044132in}{0.278150in}%
\pgfsys@useobject{currentmarker}{}%
\end{pgfscope}%
\end{pgfscope}%
\begin{pgfscope}%
\definecolor{textcolor}{rgb}{0.000000,0.000000,0.000000}%
\pgfsetstrokecolor{textcolor}%
\pgfsetfillcolor{textcolor}%
\pgftext[x=2.044132in,y=0.180927in,,top]{\color{textcolor}\rmfamily\fontsize{12.000000}{14.400000}\selectfont 1}%
\end{pgfscope}%
\begin{pgfscope}%
\pgfsetbuttcap%
\pgfsetroundjoin%
\definecolor{currentfill}{rgb}{0.000000,0.000000,0.000000}%
\pgfsetfillcolor{currentfill}%
\pgfsetlinewidth{0.803000pt}%
\definecolor{currentstroke}{rgb}{0.000000,0.000000,0.000000}%
\pgfsetstrokecolor{currentstroke}%
\pgfsetdash{}{0pt}%
\pgfsys@defobject{currentmarker}{\pgfqpoint{-0.048611in}{0.000000in}}{\pgfqpoint{-0.000000in}{0.000000in}}{%
\pgfpathmoveto{\pgfqpoint{-0.000000in}{0.000000in}}%
\pgfpathlineto{\pgfqpoint{-0.048611in}{0.000000in}}%
\pgfusepath{stroke,fill}%
}%
\begin{pgfscope}%
\pgfsys@transformshift{0.506574in}{0.278150in}%
\pgfsys@useobject{currentmarker}{}%
\end{pgfscope}%
\end{pgfscope}%
\begin{pgfscope}%
\definecolor{textcolor}{rgb}{0.000000,0.000000,0.000000}%
\pgfsetstrokecolor{textcolor}%
\pgfsetfillcolor{textcolor}%
\pgftext[x=0.303313in, y=0.214836in, left, base]{\color{textcolor}\rmfamily\fontsize{12.000000}{14.400000}\selectfont 0}%
\end{pgfscope}%
\begin{pgfscope}%
\pgfsetbuttcap%
\pgfsetroundjoin%
\definecolor{currentfill}{rgb}{0.000000,0.000000,0.000000}%
\pgfsetfillcolor{currentfill}%
\pgfsetlinewidth{0.803000pt}%
\definecolor{currentstroke}{rgb}{0.000000,0.000000,0.000000}%
\pgfsetstrokecolor{currentstroke}%
\pgfsetdash{}{0pt}%
\pgfsys@defobject{currentmarker}{\pgfqpoint{-0.048611in}{0.000000in}}{\pgfqpoint{-0.000000in}{0.000000in}}{%
\pgfpathmoveto{\pgfqpoint{-0.000000in}{0.000000in}}%
\pgfpathlineto{\pgfqpoint{-0.048611in}{0.000000in}}%
\pgfusepath{stroke,fill}%
}%
\begin{pgfscope}%
\pgfsys@transformshift{0.506574in}{0.868311in}%
\pgfsys@useobject{currentmarker}{}%
\end{pgfscope}%
\end{pgfscope}%
\begin{pgfscope}%
\definecolor{textcolor}{rgb}{0.000000,0.000000,0.000000}%
\pgfsetstrokecolor{textcolor}%
\pgfsetfillcolor{textcolor}%
\pgftext[x=0.022683in, y=0.804997in, left, base]{\color{textcolor}\rmfamily\fontsize{12.000000}{14.400000}\selectfont \(\displaystyle 1-\beta\)}%
\end{pgfscope}%
\begin{pgfscope}%
\pgfsetbuttcap%
\pgfsetroundjoin%
\definecolor{currentfill}{rgb}{0.000000,0.000000,0.000000}%
\pgfsetfillcolor{currentfill}%
\pgfsetlinewidth{0.803000pt}%
\definecolor{currentstroke}{rgb}{0.000000,0.000000,0.000000}%
\pgfsetstrokecolor{currentstroke}%
\pgfsetdash{}{0pt}%
\pgfsys@defobject{currentmarker}{\pgfqpoint{-0.048611in}{0.000000in}}{\pgfqpoint{-0.000000in}{0.000000in}}{%
\pgfpathmoveto{\pgfqpoint{-0.000000in}{0.000000in}}%
\pgfpathlineto{\pgfqpoint{-0.048611in}{0.000000in}}%
\pgfusepath{stroke,fill}%
}%
\begin{pgfscope}%
\pgfsys@transformshift{0.506574in}{1.458473in}%
\pgfsys@useobject{currentmarker}{}%
\end{pgfscope}%
\end{pgfscope}%
\begin{pgfscope}%
\definecolor{textcolor}{rgb}{0.000000,0.000000,0.000000}%
\pgfsetstrokecolor{textcolor}%
\pgfsetfillcolor{textcolor}%
\pgftext[x=0.303313in, y=1.395159in, left, base]{\color{textcolor}\rmfamily\fontsize{12.000000}{14.400000}\selectfont 1}%
\end{pgfscope}%
\begin{pgfscope}%
\pgfsetbuttcap%
\pgfsetroundjoin%
\pgfsetlinewidth{0.702625pt}%
\definecolor{currentstroke}{rgb}{0.121569,0.466667,0.705882}%
\pgfsetstrokecolor{currentstroke}%
\pgfsetdash{{2.590000pt}{1.120000pt}}{0.000000pt}%
\pgfpathmoveto{\pgfqpoint{2.044132in}{0.278150in}}%
\pgfpathlineto{\pgfqpoint{0.506574in}{1.458473in}}%
\pgfusepath{stroke}%
\end{pgfscope}%
\begin{pgfscope}%
\pgfsetrectcap%
\pgfsetmiterjoin%
\pgfsetlinewidth{0.803000pt}%
\definecolor{currentstroke}{rgb}{0.000000,0.000000,0.000000}%
\pgfsetstrokecolor{currentstroke}%
\pgfsetdash{}{0pt}%
\pgfpathmoveto{\pgfqpoint{0.506574in}{0.278150in}}%
\pgfpathlineto{\pgfqpoint{0.506574in}{2.166667in}}%
\pgfusepath{stroke}%
\end{pgfscope}%
\begin{pgfscope}%
\pgfsetrectcap%
\pgfsetmiterjoin%
\pgfsetlinewidth{0.803000pt}%
\definecolor{currentstroke}{rgb}{0.000000,0.000000,0.000000}%
\pgfsetstrokecolor{currentstroke}%
\pgfsetdash{}{0pt}%
\pgfpathmoveto{\pgfqpoint{0.506574in}{0.278150in}}%
\pgfpathlineto{\pgfqpoint{2.966667in}{0.278150in}}%
\pgfusepath{stroke}%
\end{pgfscope}%
\begin{pgfscope}%
\pgfsetrectcap%
\pgfsetroundjoin%
\pgfsetlinewidth{1.505625pt}%
\definecolor{currentstroke}{rgb}{0.172549,0.627451,0.172549}%
\pgfsetstrokecolor{currentstroke}%
\pgfsetdash{}{0pt}%
\pgfpathmoveto{\pgfqpoint{0.521949in}{0.289953in}}%
\pgfpathlineto{\pgfqpoint{0.544100in}{0.289194in}}%
\pgfpathlineto{\pgfqpoint{0.596311in}{0.288516in}}%
\pgfpathlineto{\pgfqpoint{0.771087in}{0.287795in}}%
\pgfpathlineto{\pgfqpoint{1.251240in}{0.287393in}}%
\pgfpathlineto{\pgfqpoint{1.809835in}{0.287879in}}%
\pgfpathlineto{\pgfqpoint{1.975886in}{0.288843in}}%
\pgfpathlineto{\pgfqpoint{2.017141in}{0.290065in}}%
\pgfpathlineto{\pgfqpoint{2.024572in}{0.291173in}}%
\pgfpathlineto{\pgfqpoint{2.025719in}{0.292075in}}%
\pgfpathlineto{\pgfqpoint{2.025455in}{0.293041in}}%
\pgfpathlineto{\pgfqpoint{2.022943in}{0.295558in}}%
\pgfpathlineto{\pgfqpoint{2.001426in}{0.313326in}}%
\pgfpathlineto{\pgfqpoint{1.872380in}{0.418723in}}%
\pgfpathlineto{\pgfqpoint{1.758637in}{0.510243in}}%
\pgfpathlineto{\pgfqpoint{1.661728in}{0.586886in}}%
\pgfpathlineto{\pgfqpoint{1.576315in}{0.653116in}}%
\pgfpathlineto{\pgfqpoint{1.501454in}{0.709846in}}%
\pgfpathlineto{\pgfqpoint{1.434481in}{0.759239in}}%
\pgfpathlineto{\pgfqpoint{1.375532in}{0.801322in}}%
\pgfpathlineto{\pgfqpoint{1.325937in}{0.835379in}}%
\pgfpathlineto{\pgfqpoint{1.284019in}{0.862838in}}%
\pgfpathlineto{\pgfqpoint{1.275353in}{0.868311in}}%
\pgfpathlineto{\pgfqpoint{1.275353in}{0.868311in}}%
\pgfusepath{stroke}%
\end{pgfscope}%
\begin{pgfscope}%
\pgfsetrectcap%
\pgfsetroundjoin%
\pgfsetlinewidth{1.505625pt}%
\definecolor{currentstroke}{rgb}{0.172549,0.627451,0.172549}%
\pgfsetstrokecolor{currentstroke}%
\pgfsetdash{}{0pt}%
\pgfpathmoveto{\pgfqpoint{2.028756in}{0.289953in}}%
\pgfpathlineto{\pgfqpoint{2.028158in}{0.291061in}}%
\pgfpathlineto{\pgfqpoint{2.024093in}{0.294662in}}%
\pgfpathlineto{\pgfqpoint{1.942086in}{0.361971in}}%
\pgfpathlineto{\pgfqpoint{1.819369in}{0.461560in}}%
\pgfpathlineto{\pgfqpoint{1.712780in}{0.546685in}}%
\pgfpathlineto{\pgfqpoint{1.618063in}{0.620921in}}%
\pgfpathlineto{\pgfqpoint{1.535131in}{0.684500in}}%
\pgfpathlineto{\pgfqpoint{1.465721in}{0.736382in}}%
\pgfpathlineto{\pgfqpoint{1.403476in}{0.781561in}}%
\pgfpathlineto{\pgfqpoint{1.350408in}{0.818755in}}%
\pgfpathlineto{\pgfqpoint{1.304563in}{0.849563in}}%
\pgfpathlineto{\pgfqpoint{1.275353in}{0.868311in}}%
\pgfpathlineto{\pgfqpoint{1.275353in}{0.868311in}}%
\pgfusepath{stroke}%
\end{pgfscope}%
\begin{pgfscope}%
\pgfsetrectcap%
\pgfsetroundjoin%
\pgfsetlinewidth{1.505625pt}%
\definecolor{currentstroke}{rgb}{0.172549,0.627451,0.172549}%
\pgfsetstrokecolor{currentstroke}%
\pgfsetstrokeopacity{0.200000}%
\pgfsetdash{}{0pt}%
\pgfpathmoveto{\pgfqpoint{0.660329in}{0.868311in}}%
\pgfpathlineto{\pgfqpoint{0.680204in}{0.857767in}}%
\pgfpathlineto{\pgfqpoint{0.702350in}{0.847768in}}%
\pgfpathlineto{\pgfqpoint{0.726855in}{0.838330in}}%
\pgfpathlineto{\pgfqpoint{0.753750in}{0.829470in}}%
\pgfpathlineto{\pgfqpoint{0.782998in}{0.821204in}}%
\pgfpathlineto{\pgfqpoint{0.814481in}{0.813547in}}%
\pgfpathlineto{\pgfqpoint{0.847991in}{0.806509in}}%
\pgfpathlineto{\pgfqpoint{0.883225in}{0.800098in}}%
\pgfpathlineto{\pgfqpoint{0.919792in}{0.794317in}}%
\pgfpathlineto{\pgfqpoint{0.957224in}{0.789163in}}%
\pgfpathlineto{\pgfqpoint{0.994998in}{0.784629in}}%
\pgfpathlineto{\pgfqpoint{1.069375in}{0.777353in}}%
\pgfpathlineto{\pgfqpoint{1.138755in}{0.772302in}}%
\pgfpathlineto{\pgfqpoint{1.199909in}{0.769214in}}%
\pgfpathlineto{\pgfqpoint{1.251046in}{0.767782in}}%
\pgfpathlineto{\pgfqpoint{1.291806in}{0.767693in}}%
\pgfpathlineto{\pgfqpoint{1.322901in}{0.768649in}}%
\pgfpathlineto{\pgfqpoint{1.345649in}{0.770392in}}%
\pgfpathlineto{\pgfqpoint{1.361571in}{0.772707in}}%
\pgfpathlineto{\pgfqpoint{1.372134in}{0.775421in}}%
\pgfpathlineto{\pgfqpoint{1.378611in}{0.778402in}}%
\pgfpathlineto{\pgfqpoint{1.382044in}{0.781546in}}%
\pgfpathlineto{\pgfqpoint{1.383253in}{0.784779in}}%
\pgfpathlineto{\pgfqpoint{1.382860in}{0.788043in}}%
\pgfpathlineto{\pgfqpoint{1.381336in}{0.791298in}}%
\pgfpathlineto{\pgfqpoint{1.377657in}{0.796101in}}%
\pgfpathlineto{\pgfqpoint{1.371312in}{0.802260in}}%
\pgfpathlineto{\pgfqpoint{1.360730in}{0.810825in}}%
\pgfpathlineto{\pgfqpoint{1.342169in}{0.824235in}}%
\pgfpathlineto{\pgfqpoint{1.308810in}{0.846769in}}%
\pgfpathlineto{\pgfqpoint{1.275353in}{0.868311in}}%
\pgfpathlineto{\pgfqpoint{1.275353in}{0.868311in}}%
\pgfusepath{stroke}%
\end{pgfscope}%
\begin{pgfscope}%
\pgfsetrectcap%
\pgfsetroundjoin%
\pgfsetlinewidth{1.505625pt}%
\definecolor{currentstroke}{rgb}{0.172549,0.627451,0.172549}%
\pgfsetstrokecolor{currentstroke}%
\pgfsetstrokeopacity{0.200000}%
\pgfsetdash{}{0pt}%
\pgfpathmoveto{\pgfqpoint{0.660329in}{1.458473in}}%
\pgfpathlineto{\pgfqpoint{0.656186in}{1.436973in}}%
\pgfpathlineto{\pgfqpoint{0.653067in}{1.415743in}}%
\pgfpathlineto{\pgfqpoint{0.650873in}{1.394813in}}%
\pgfpathlineto{\pgfqpoint{0.649528in}{1.374211in}}%
\pgfpathlineto{\pgfqpoint{0.648972in}{1.353959in}}%
\pgfpathlineto{\pgfqpoint{0.649162in}{1.334075in}}%
\pgfpathlineto{\pgfqpoint{0.650065in}{1.314576in}}%
\pgfpathlineto{\pgfqpoint{0.651661in}{1.295474in}}%
\pgfpathlineto{\pgfqpoint{0.653937in}{1.276783in}}%
\pgfpathlineto{\pgfqpoint{0.656887in}{1.258511in}}%
\pgfpathlineto{\pgfqpoint{0.660515in}{1.240667in}}%
\pgfpathlineto{\pgfqpoint{0.664825in}{1.223260in}}%
\pgfpathlineto{\pgfqpoint{0.669830in}{1.206297in}}%
\pgfpathlineto{\pgfqpoint{0.675544in}{1.189784in}}%
\pgfpathlineto{\pgfqpoint{0.681984in}{1.173726in}}%
\pgfpathlineto{\pgfqpoint{0.689170in}{1.158130in}}%
\pgfpathlineto{\pgfqpoint{0.697121in}{1.143000in}}%
\pgfpathlineto{\pgfqpoint{0.705855in}{1.128341in}}%
\pgfpathlineto{\pgfqpoint{0.715389in}{1.114159in}}%
\pgfpathlineto{\pgfqpoint{0.725737in}{1.100456in}}%
\pgfpathlineto{\pgfqpoint{0.736907in}{1.087237in}}%
\pgfpathlineto{\pgfqpoint{0.748903in}{1.074506in}}%
\pgfpathlineto{\pgfqpoint{0.761718in}{1.062264in}}%
\pgfpathlineto{\pgfqpoint{0.775340in}{1.050515in}}%
\pgfpathlineto{\pgfqpoint{0.789744in}{1.039259in}}%
\pgfpathlineto{\pgfqpoint{0.804893in}{1.028498in}}%
\pgfpathlineto{\pgfqpoint{0.820740in}{1.018230in}}%
\pgfpathlineto{\pgfqpoint{0.837225in}{1.008455in}}%
\pgfpathlineto{\pgfqpoint{0.854273in}{0.999169in}}%
\pgfpathlineto{\pgfqpoint{0.871800in}{0.990368in}}%
\pgfpathlineto{\pgfqpoint{0.907899in}{0.974195in}}%
\pgfpathlineto{\pgfqpoint{0.944659in}{0.959870in}}%
\pgfpathlineto{\pgfqpoint{0.981154in}{0.947301in}}%
\pgfpathlineto{\pgfqpoint{1.016481in}{0.936374in}}%
\pgfpathlineto{\pgfqpoint{1.065603in}{0.922762in}}%
\pgfpathlineto{\pgfqpoint{1.121225in}{0.909017in}}%
\pgfpathlineto{\pgfqpoint{1.188918in}{0.893779in}}%
\pgfpathlineto{\pgfqpoint{1.243872in}{0.881145in}}%
\pgfpathlineto{\pgfqpoint{1.259735in}{0.876383in}}%
\pgfpathlineto{\pgfqpoint{1.267945in}{0.872760in}}%
\pgfpathlineto{\pgfqpoint{1.275353in}{0.868311in}}%
\pgfpathlineto{\pgfqpoint{1.275353in}{0.868311in}}%
\pgfusepath{stroke}%
\end{pgfscope}%
\begin{pgfscope}%
\pgfsetrectcap%
\pgfsetroundjoin%
\pgfsetlinewidth{1.505625pt}%
\definecolor{currentstroke}{rgb}{0.172549,0.627451,0.172549}%
\pgfsetstrokecolor{currentstroke}%
\pgfsetstrokeopacity{0.200000}%
\pgfsetdash{}{0pt}%
\pgfpathmoveto{\pgfqpoint{2.812911in}{1.458473in}}%
\pgfpathlineto{\pgfqpoint{1.691329in}{1.520519in}}%
\pgfpathlineto{\pgfqpoint{1.450478in}{1.531186in}}%
\pgfpathlineto{\pgfqpoint{1.287419in}{1.534504in}}%
\pgfpathlineto{\pgfqpoint{1.170401in}{1.532655in}}%
\pgfpathlineto{\pgfqpoint{1.082948in}{1.527037in}}%
\pgfpathlineto{\pgfqpoint{1.015637in}{1.518605in}}%
\pgfpathlineto{\pgfqpoint{0.962678in}{1.508039in}}%
\pgfpathlineto{\pgfqpoint{0.920314in}{1.495842in}}%
\pgfpathlineto{\pgfqpoint{0.885999in}{1.482396in}}%
\pgfpathlineto{\pgfqpoint{0.857947in}{1.467998in}}%
\pgfpathlineto{\pgfqpoint{0.834872in}{1.452880in}}%
\pgfpathlineto{\pgfqpoint{0.815822in}{1.437229in}}%
\pgfpathlineto{\pgfqpoint{0.800084in}{1.421197in}}%
\pgfpathlineto{\pgfqpoint{0.787111in}{1.404905in}}%
\pgfpathlineto{\pgfqpoint{0.776480in}{1.388457in}}%
\pgfpathlineto{\pgfqpoint{0.767859in}{1.371935in}}%
\pgfpathlineto{\pgfqpoint{0.760988in}{1.355412in}}%
\pgfpathlineto{\pgfqpoint{0.755657in}{1.338946in}}%
\pgfpathlineto{\pgfqpoint{0.751701in}{1.322587in}}%
\pgfpathlineto{\pgfqpoint{0.748986in}{1.306379in}}%
\pgfpathlineto{\pgfqpoint{0.747406in}{1.290358in}}%
\pgfpathlineto{\pgfqpoint{0.746872in}{1.274555in}}%
\pgfpathlineto{\pgfqpoint{0.747316in}{1.258998in}}%
\pgfpathlineto{\pgfqpoint{0.748682in}{1.243710in}}%
\pgfpathlineto{\pgfqpoint{0.750923in}{1.228711in}}%
\pgfpathlineto{\pgfqpoint{0.754003in}{1.214020in}}%
\pgfpathlineto{\pgfqpoint{0.757891in}{1.199651in}}%
\pgfpathlineto{\pgfqpoint{0.762563in}{1.185619in}}%
\pgfpathlineto{\pgfqpoint{0.767994in}{1.171935in}}%
\pgfpathlineto{\pgfqpoint{0.774167in}{1.158609in}}%
\pgfpathlineto{\pgfqpoint{0.781062in}{1.145651in}}%
\pgfpathlineto{\pgfqpoint{0.788660in}{1.133068in}}%
\pgfpathlineto{\pgfqpoint{0.796941in}{1.120867in}}%
\pgfpathlineto{\pgfqpoint{0.805882in}{1.109054in}}%
\pgfpathlineto{\pgfqpoint{0.815459in}{1.097631in}}%
\pgfpathlineto{\pgfqpoint{0.825642in}{1.086604in}}%
\pgfpathlineto{\pgfqpoint{0.836399in}{1.075974in}}%
\pgfpathlineto{\pgfqpoint{0.847691in}{1.065743in}}%
\pgfpathlineto{\pgfqpoint{0.859478in}{1.055910in}}%
\pgfpathlineto{\pgfqpoint{0.884336in}{1.037435in}}%
\pgfpathlineto{\pgfqpoint{0.910540in}{1.020528in}}%
\pgfpathlineto{\pgfqpoint{0.937593in}{1.005153in}}%
\pgfpathlineto{\pgfqpoint{0.964964in}{0.991252in}}%
\pgfpathlineto{\pgfqpoint{0.992117in}{0.978757in}}%
\pgfpathlineto{\pgfqpoint{1.018551in}{0.967581in}}%
\pgfpathlineto{\pgfqpoint{1.055918in}{0.953083in}}%
\pgfpathlineto{\pgfqpoint{1.099916in}{0.937440in}}%
\pgfpathlineto{\pgfqpoint{1.176959in}{0.911810in}}%
\pgfpathlineto{\pgfqpoint{1.214336in}{0.898732in}}%
\pgfpathlineto{\pgfqpoint{1.234772in}{0.890500in}}%
\pgfpathlineto{\pgfqpoint{1.250444in}{0.882993in}}%
\pgfpathlineto{\pgfqpoint{1.265229in}{0.874571in}}%
\pgfpathlineto{\pgfqpoint{1.275353in}{0.868311in}}%
\pgfpathlineto{\pgfqpoint{1.275353in}{0.868311in}}%
\pgfusepath{stroke}%
\end{pgfscope}%
\begin{pgfscope}%
\pgfsetrectcap%
\pgfsetroundjoin%
\pgfsetlinewidth{1.505625pt}%
\definecolor{currentstroke}{rgb}{0.172549,0.627451,0.172549}%
\pgfsetstrokecolor{currentstroke}%
\pgfsetstrokeopacity{0.200000}%
\pgfsetdash{}{0pt}%
\pgfpathmoveto{\pgfqpoint{2.812911in}{1.163392in}}%
\pgfpathlineto{\pgfqpoint{2.186539in}{1.195305in}}%
\pgfpathlineto{\pgfqpoint{1.620018in}{1.227848in}}%
\pgfpathlineto{\pgfqpoint{1.468209in}{1.235359in}}%
\pgfpathlineto{\pgfqpoint{1.357529in}{1.239210in}}%
\pgfpathlineto{\pgfqpoint{1.273660in}{1.240310in}}%
\pgfpathlineto{\pgfqpoint{1.208305in}{1.239289in}}%
\pgfpathlineto{\pgfqpoint{1.156314in}{1.236602in}}%
\pgfpathlineto{\pgfqpoint{1.114317in}{1.232592in}}%
\pgfpathlineto{\pgfqpoint{1.080015in}{1.227524in}}%
\pgfpathlineto{\pgfqpoint{1.051783in}{1.221607in}}%
\pgfpathlineto{\pgfqpoint{1.028439in}{1.215009in}}%
\pgfpathlineto{\pgfqpoint{1.009103in}{1.207868in}}%
\pgfpathlineto{\pgfqpoint{0.993105in}{1.200298in}}%
\pgfpathlineto{\pgfqpoint{0.979926in}{1.192393in}}%
\pgfpathlineto{\pgfqpoint{0.969159in}{1.184233in}}%
\pgfpathlineto{\pgfqpoint{0.960476in}{1.175887in}}%
\pgfpathlineto{\pgfqpoint{0.953615in}{1.167412in}}%
\pgfpathlineto{\pgfqpoint{0.948358in}{1.158859in}}%
\pgfpathlineto{\pgfqpoint{0.944527in}{1.150270in}}%
\pgfpathlineto{\pgfqpoint{0.941969in}{1.141684in}}%
\pgfpathlineto{\pgfqpoint{0.940557in}{1.133131in}}%
\pgfpathlineto{\pgfqpoint{0.940181in}{1.124641in}}%
\pgfpathlineto{\pgfqpoint{0.940745in}{1.116238in}}%
\pgfpathlineto{\pgfqpoint{0.942162in}{1.107942in}}%
\pgfpathlineto{\pgfqpoint{0.944359in}{1.099774in}}%
\pgfpathlineto{\pgfqpoint{0.947265in}{1.091747in}}%
\pgfpathlineto{\pgfqpoint{0.950819in}{1.083877in}}%
\pgfpathlineto{\pgfqpoint{0.959636in}{1.068649in}}%
\pgfpathlineto{\pgfqpoint{0.970383in}{1.054163in}}%
\pgfpathlineto{\pgfqpoint{0.982672in}{1.040466in}}%
\pgfpathlineto{\pgfqpoint{0.996141in}{1.027587in}}%
\pgfpathlineto{\pgfqpoint{1.010455in}{1.015537in}}%
\pgfpathlineto{\pgfqpoint{1.025302in}{1.004312in}}%
\pgfpathlineto{\pgfqpoint{1.047953in}{0.988986in}}%
\pgfpathlineto{\pgfqpoint{1.070332in}{0.975388in}}%
\pgfpathlineto{\pgfqpoint{1.098641in}{0.959727in}}%
\pgfpathlineto{\pgfqpoint{1.130208in}{0.943626in}}%
\pgfpathlineto{\pgfqpoint{1.189341in}{0.915158in}}%
\pgfpathlineto{\pgfqpoint{1.226182in}{0.896779in}}%
\pgfpathlineto{\pgfqpoint{1.250670in}{0.883278in}}%
\pgfpathlineto{\pgfqpoint{1.275353in}{0.868311in}}%
\pgfpathlineto{\pgfqpoint{1.275353in}{0.868311in}}%
\pgfusepath{stroke}%
\end{pgfscope}%
\begin{pgfscope}%
\pgfsetrectcap%
\pgfsetroundjoin%
\pgfsetlinewidth{1.505625pt}%
\definecolor{currentstroke}{rgb}{0.172549,0.627451,0.172549}%
\pgfsetstrokecolor{currentstroke}%
\pgfsetstrokeopacity{0.200000}%
\pgfsetdash{}{0pt}%
\pgfpathmoveto{\pgfqpoint{2.812911in}{0.868311in}}%
\pgfpathlineto{\pgfqpoint{2.299104in}{0.890766in}}%
\pgfpathlineto{\pgfqpoint{2.004334in}{0.906746in}}%
\pgfpathlineto{\pgfqpoint{1.813699in}{0.918605in}}%
\pgfpathlineto{\pgfqpoint{1.583371in}{0.934491in}}%
\pgfpathlineto{\pgfqpoint{1.405139in}{0.946828in}}%
\pgfpathlineto{\pgfqpoint{1.337029in}{0.950562in}}%
\pgfpathlineto{\pgfqpoint{1.290467in}{0.952045in}}%
\pgfpathlineto{\pgfqpoint{1.257848in}{0.951959in}}%
\pgfpathlineto{\pgfqpoint{1.234765in}{0.950768in}}%
\pgfpathlineto{\pgfqpoint{1.218471in}{0.948799in}}%
\pgfpathlineto{\pgfqpoint{1.207150in}{0.946291in}}%
\pgfpathlineto{\pgfqpoint{1.199549in}{0.943420in}}%
\pgfpathlineto{\pgfqpoint{1.194765in}{0.940321in}}%
\pgfpathlineto{\pgfqpoint{1.192133in}{0.937092in}}%
\pgfpathlineto{\pgfqpoint{1.191146in}{0.933809in}}%
\pgfpathlineto{\pgfqpoint{1.191412in}{0.930530in}}%
\pgfpathlineto{\pgfqpoint{1.192627in}{0.927297in}}%
\pgfpathlineto{\pgfqpoint{1.195712in}{0.922599in}}%
\pgfpathlineto{\pgfqpoint{1.201275in}{0.916717in}}%
\pgfpathlineto{\pgfqpoint{1.209247in}{0.910082in}}%
\pgfpathlineto{\pgfqpoint{1.221954in}{0.901162in}}%
\pgfpathlineto{\pgfqpoint{1.243025in}{0.888024in}}%
\pgfpathlineto{\pgfqpoint{1.275353in}{0.868311in}}%
\pgfpathlineto{\pgfqpoint{1.275353in}{0.868311in}}%
\pgfusepath{stroke}%
\end{pgfscope}%
\begin{pgfscope}%
\pgfsetrectcap%
\pgfsetroundjoin%
\pgfsetlinewidth{1.505625pt}%
\definecolor{currentstroke}{rgb}{0.172549,0.627451,0.172549}%
\pgfsetstrokecolor{currentstroke}%
\pgfsetstrokeopacity{0.200000}%
\pgfsetdash{}{0pt}%
\pgfpathmoveto{\pgfqpoint{2.812911in}{0.573230in}}%
\pgfpathlineto{\pgfqpoint{2.418413in}{0.585073in}}%
\pgfpathlineto{\pgfqpoint{2.185296in}{0.594517in}}%
\pgfpathlineto{\pgfqpoint{2.033387in}{0.602426in}}%
\pgfpathlineto{\pgfqpoint{1.927907in}{0.609281in}}%
\pgfpathlineto{\pgfqpoint{1.851320in}{0.615376in}}%
\pgfpathlineto{\pgfqpoint{1.793820in}{0.620904in}}%
\pgfpathlineto{\pgfqpoint{1.749497in}{0.625998in}}%
\pgfpathlineto{\pgfqpoint{1.714580in}{0.630755in}}%
\pgfpathlineto{\pgfqpoint{1.686554in}{0.635244in}}%
\pgfpathlineto{\pgfqpoint{1.644718in}{0.643615in}}%
\pgfpathlineto{\pgfqpoint{1.615173in}{0.651399in}}%
\pgfpathlineto{\pgfqpoint{1.593124in}{0.658771in}}%
\pgfpathlineto{\pgfqpoint{1.575801in}{0.665837in}}%
\pgfpathlineto{\pgfqpoint{1.561533in}{0.672666in}}%
\pgfpathlineto{\pgfqpoint{1.543707in}{0.682552in}}%
\pgfpathlineto{\pgfqpoint{1.523776in}{0.695182in}}%
\pgfpathlineto{\pgfqpoint{1.498026in}{0.713088in}}%
\pgfpathlineto{\pgfqpoint{1.339660in}{0.826103in}}%
\pgfpathlineto{\pgfqpoint{1.295297in}{0.855598in}}%
\pgfpathlineto{\pgfqpoint{1.275353in}{0.868311in}}%
\pgfpathlineto{\pgfqpoint{1.275353in}{0.868311in}}%
\pgfusepath{stroke}%
\end{pgfscope}%
\begin{pgfscope}%
\pgfsetrectcap%
\pgfsetroundjoin%
\pgfsetlinewidth{1.505625pt}%
\definecolor{currentstroke}{rgb}{0.172549,0.627451,0.172549}%
\pgfsetstrokecolor{currentstroke}%
\pgfsetstrokeopacity{0.200000}%
\pgfsetdash{}{0pt}%
\pgfpathmoveto{\pgfqpoint{0.660329in}{0.573230in}}%
\pgfpathlineto{\pgfqpoint{0.693607in}{0.568014in}}%
\pgfpathlineto{\pgfqpoint{0.732752in}{0.563196in}}%
\pgfpathlineto{\pgfqpoint{0.778128in}{0.558812in}}%
\pgfpathlineto{\pgfqpoint{0.829827in}{0.554895in}}%
\pgfpathlineto{\pgfqpoint{0.887569in}{0.551478in}}%
\pgfpathlineto{\pgfqpoint{0.950627in}{0.548589in}}%
\pgfpathlineto{\pgfqpoint{1.017806in}{0.546252in}}%
\pgfpathlineto{\pgfqpoint{1.087497in}{0.544479in}}%
\pgfpathlineto{\pgfqpoint{1.226749in}{0.542617in}}%
\pgfpathlineto{\pgfqpoint{1.353423in}{0.542874in}}%
\pgfpathlineto{\pgfqpoint{1.457054in}{0.544965in}}%
\pgfpathlineto{\pgfqpoint{1.498906in}{0.546584in}}%
\pgfpathlineto{\pgfqpoint{1.534224in}{0.548524in}}%
\pgfpathlineto{\pgfqpoint{1.563443in}{0.550739in}}%
\pgfpathlineto{\pgfqpoint{1.587154in}{0.553189in}}%
\pgfpathlineto{\pgfqpoint{1.606023in}{0.555836in}}%
\pgfpathlineto{\pgfqpoint{1.620724in}{0.558649in}}%
\pgfpathlineto{\pgfqpoint{1.631896in}{0.561599in}}%
\pgfpathlineto{\pgfqpoint{1.640122in}{0.564662in}}%
\pgfpathlineto{\pgfqpoint{1.645912in}{0.567818in}}%
\pgfpathlineto{\pgfqpoint{1.649704in}{0.571051in}}%
\pgfpathlineto{\pgfqpoint{1.651870in}{0.574345in}}%
\pgfpathlineto{\pgfqpoint{1.652717in}{0.577690in}}%
\pgfpathlineto{\pgfqpoint{1.652500in}{0.581076in}}%
\pgfpathlineto{\pgfqpoint{1.651425in}{0.584494in}}%
\pgfpathlineto{\pgfqpoint{1.647345in}{0.591399in}}%
\pgfpathlineto{\pgfqpoint{1.641477in}{0.598364in}}%
\pgfpathlineto{\pgfqpoint{1.630673in}{0.608858in}}%
\pgfpathlineto{\pgfqpoint{1.614357in}{0.622828in}}%
\pgfpathlineto{\pgfqpoint{1.584105in}{0.646912in}}%
\pgfpathlineto{\pgfqpoint{1.528800in}{0.689271in}}%
\pgfpathlineto{\pgfqpoint{1.460385in}{0.740309in}}%
\pgfpathlineto{\pgfqpoint{1.399415in}{0.784455in}}%
\pgfpathlineto{\pgfqpoint{1.346072in}{0.821728in}}%
\pgfpathlineto{\pgfqpoint{1.300640in}{0.852128in}}%
\pgfpathlineto{\pgfqpoint{1.275353in}{0.868311in}}%
\pgfpathlineto{\pgfqpoint{1.275353in}{0.868311in}}%
\pgfusepath{stroke}%
\end{pgfscope}%
\begin{pgfscope}%
\pgfsetbuttcap%
\pgfsetroundjoin%
\pgfsetlinewidth{0.702625pt}%
\definecolor{currentstroke}{rgb}{1.000000,0.498039,0.054902}%
\pgfsetstrokecolor{currentstroke}%
\pgfsetdash{{2.590000pt}{1.120000pt}}{0.000000pt}%
\pgfpathmoveto{\pgfqpoint{0.506574in}{0.278150in}}%
\pgfpathlineto{\pgfqpoint{2.966667in}{0.278150in}}%
\pgfusepath{stroke}%
\end{pgfscope}%
\begin{pgfscope}%
\definecolor{textcolor}{rgb}{0.000000,0.000000,0.000000}%
\pgfsetstrokecolor{textcolor}%
\pgfsetfillcolor{textcolor}%
\pgftext[x=2.966667in,y=0.180927in,right,top]{\color{textcolor}\rmfamily\fontsize{12.000000}{14.400000}\selectfont \(\displaystyle u\)}%
\end{pgfscope}%
\begin{pgfscope}%
\definecolor{textcolor}{rgb}{0.000000,0.000000,0.000000}%
\pgfsetstrokecolor{textcolor}%
\pgfsetfillcolor{textcolor}%
\pgftext[x=0.409351in,y=2.166667in,right,top]{\color{textcolor}\rmfamily\fontsize{12.000000}{14.400000}\selectfont \(\displaystyle v\)}%
\end{pgfscope}%
\begin{pgfscope}%
\pgfsetroundcap%
\pgfsetroundjoin%
\pgfsetlinewidth{1.003750pt}%
\definecolor{currentstroke}{rgb}{0.172549,0.627451,0.172549}%
\pgfsetstrokecolor{currentstroke}%
\pgfsetdash{}{0pt}%
\pgfpathmoveto{\pgfqpoint{0.771087in}{0.287795in}}%
\pgfpathquadraticcurveto{\pgfqpoint{0.807516in}{0.287728in}}{\pgfqpoint{0.828416in}{0.287690in}}%
\pgfusepath{stroke}%
\end{pgfscope}%
\begin{pgfscope}%
\pgfsetroundcap%
\pgfsetroundjoin%
\definecolor{currentfill}{rgb}{0.172549,0.627451,0.172549}%
\pgfsetfillcolor{currentfill}%
\pgfsetlinewidth{1.003750pt}%
\definecolor{currentstroke}{rgb}{0.172549,0.627451,0.172549}%
\pgfsetstrokecolor{currentstroke}%
\pgfsetdash{}{0pt}%
\pgfpathmoveto{\pgfqpoint{0.728508in}{0.337873in}}%
\pgfpathlineto{\pgfqpoint{0.828416in}{0.287690in}}%
\pgfpathlineto{\pgfqpoint{0.728325in}{0.237873in}}%
\pgfpathlineto{\pgfqpoint{0.728508in}{0.337873in}}%
\pgfpathclose%
\pgfusepath{stroke,fill}%
\end{pgfscope}%
\begin{pgfscope}%
\pgfsetroundcap%
\pgfsetroundjoin%
\pgfsetlinewidth{1.003750pt}%
\definecolor{currentstroke}{rgb}{0.172549,0.627451,0.172549}%
\pgfsetstrokecolor{currentstroke}%
\pgfsetdash{}{0pt}%
\pgfpathmoveto{\pgfqpoint{1.577931in}{0.651877in}}%
\pgfpathquadraticcurveto{\pgfqpoint{1.575740in}{0.653556in}}{\pgfqpoint{1.585873in}{0.645788in}}%
\pgfusepath{stroke}%
\end{pgfscope}%
\begin{pgfscope}%
\pgfsetroundcap%
\pgfsetroundjoin%
\definecolor{currentfill}{rgb}{0.172549,0.627451,0.172549}%
\pgfsetfillcolor{currentfill}%
\pgfsetlinewidth{1.003750pt}%
\definecolor{currentstroke}{rgb}{0.172549,0.627451,0.172549}%
\pgfsetstrokecolor{currentstroke}%
\pgfsetdash{}{0pt}%
\pgfpathmoveto{\pgfqpoint{1.634813in}{0.545265in}}%
\pgfpathlineto{\pgfqpoint{1.585873in}{0.645788in}}%
\pgfpathlineto{\pgfqpoint{1.695656in}{0.624626in}}%
\pgfpathlineto{\pgfqpoint{1.634813in}{0.545265in}}%
\pgfpathclose%
\pgfusepath{stroke,fill}%
\end{pgfscope}%
\begin{pgfscope}%
\pgfsetroundcap%
\pgfsetroundjoin%
\pgfsetlinewidth{1.003750pt}%
\definecolor{currentstroke}{rgb}{0.172549,0.627451,0.172549}%
\pgfsetstrokecolor{currentstroke}%
\pgfsetstrokeopacity{0.200000}%
\pgfsetdash{}{0pt}%
\pgfpathmoveto{\pgfqpoint{0.957224in}{0.789163in}}%
\pgfpathquadraticcurveto{\pgfqpoint{0.976111in}{0.786896in}}{\pgfqpoint{0.979580in}{0.786480in}}%
\pgfusepath{stroke}%
\end{pgfscope}%
\begin{pgfscope}%
\pgfsetroundcap%
\pgfsetroundjoin%
\definecolor{currentfill}{rgb}{0.172549,0.627451,0.172549}%
\pgfsetfillcolor{currentfill}%
\pgfsetfillopacity{0.200000}%
\pgfsetlinewidth{1.003750pt}%
\definecolor{currentstroke}{rgb}{0.172549,0.627451,0.172549}%
\pgfsetstrokecolor{currentstroke}%
\pgfsetstrokeopacity{0.200000}%
\pgfsetdash{}{0pt}%
\pgfpathmoveto{\pgfqpoint{0.886252in}{0.848042in}}%
\pgfpathlineto{\pgfqpoint{0.979580in}{0.786480in}}%
\pgfpathlineto{\pgfqpoint{0.874334in}{0.748755in}}%
\pgfpathlineto{\pgfqpoint{0.886252in}{0.848042in}}%
\pgfpathclose%
\pgfusepath{stroke,fill}%
\end{pgfscope}%
\begin{pgfscope}%
\pgfsetroundcap%
\pgfsetroundjoin%
\pgfsetlinewidth{1.003750pt}%
\definecolor{currentstroke}{rgb}{0.172549,0.627451,0.172549}%
\pgfsetstrokecolor{currentstroke}%
\pgfsetstrokeopacity{0.200000}%
\pgfsetdash{}{0pt}%
\pgfpathmoveto{\pgfqpoint{0.871800in}{0.990368in}}%
\pgfpathquadraticcurveto{\pgfqpoint{0.880755in}{0.986207in}}{\pgfqpoint{0.875628in}{0.988589in}}%
\pgfusepath{stroke}%
\end{pgfscope}%
\begin{pgfscope}%
\pgfsetroundcap%
\pgfsetroundjoin%
\definecolor{currentfill}{rgb}{0.172549,0.627451,0.172549}%
\pgfsetfillcolor{currentfill}%
\pgfsetfillopacity{0.200000}%
\pgfsetlinewidth{1.003750pt}%
\definecolor{currentstroke}{rgb}{0.172549,0.627451,0.172549}%
\pgfsetstrokecolor{currentstroke}%
\pgfsetstrokeopacity{0.200000}%
\pgfsetdash{}{0pt}%
\pgfpathmoveto{\pgfqpoint{0.806010in}{1.076072in}}%
\pgfpathlineto{\pgfqpoint{0.875628in}{0.988589in}}%
\pgfpathlineto{\pgfqpoint{0.763871in}{0.985385in}}%
\pgfpathlineto{\pgfqpoint{0.806010in}{1.076072in}}%
\pgfpathclose%
\pgfusepath{stroke,fill}%
\end{pgfscope}%
\begin{pgfscope}%
\pgfsetroundcap%
\pgfsetroundjoin%
\pgfsetlinewidth{1.003750pt}%
\definecolor{currentstroke}{rgb}{0.172549,0.627451,0.172549}%
\pgfsetstrokecolor{currentstroke}%
\pgfsetstrokeopacity{0.200000}%
\pgfsetdash{}{0pt}%
\pgfpathmoveto{\pgfqpoint{0.885999in}{1.482396in}}%
\pgfpathquadraticcurveto{\pgfqpoint{0.871973in}{1.475197in}}{\pgfqpoint{0.871762in}{1.475089in}}%
\pgfusepath{stroke}%
\end{pgfscope}%
\begin{pgfscope}%
\pgfsetroundcap%
\pgfsetroundjoin%
\definecolor{currentfill}{rgb}{0.172549,0.627451,0.172549}%
\pgfsetfillcolor{currentfill}%
\pgfsetfillopacity{0.200000}%
\pgfsetlinewidth{1.003750pt}%
\definecolor{currentstroke}{rgb}{0.172549,0.627451,0.172549}%
\pgfsetstrokecolor{currentstroke}%
\pgfsetstrokeopacity{0.200000}%
\pgfsetdash{}{0pt}%
\pgfpathmoveto{\pgfqpoint{0.983559in}{1.476270in}}%
\pgfpathlineto{\pgfqpoint{0.871762in}{1.475089in}}%
\pgfpathlineto{\pgfqpoint{0.937895in}{1.565235in}}%
\pgfpathlineto{\pgfqpoint{0.983559in}{1.476270in}}%
\pgfpathclose%
\pgfusepath{stroke,fill}%
\end{pgfscope}%
\begin{pgfscope}%
\pgfsetroundcap%
\pgfsetroundjoin%
\pgfsetlinewidth{1.003750pt}%
\definecolor{currentstroke}{rgb}{0.172549,0.627451,0.172549}%
\pgfsetstrokecolor{currentstroke}%
\pgfsetstrokeopacity{0.200000}%
\pgfsetdash{}{0pt}%
\pgfpathmoveto{\pgfqpoint{1.080015in}{1.227524in}}%
\pgfpathquadraticcurveto{\pgfqpoint{1.065899in}{1.224565in}}{\pgfqpoint{1.066981in}{1.224792in}}%
\pgfusepath{stroke}%
\end{pgfscope}%
\begin{pgfscope}%
\pgfsetroundcap%
\pgfsetroundjoin%
\definecolor{currentfill}{rgb}{0.172549,0.627451,0.172549}%
\pgfsetfillcolor{currentfill}%
\pgfsetfillopacity{0.200000}%
\pgfsetlinewidth{1.003750pt}%
\definecolor{currentstroke}{rgb}{0.172549,0.627451,0.172549}%
\pgfsetstrokecolor{currentstroke}%
\pgfsetstrokeopacity{0.200000}%
\pgfsetdash{}{0pt}%
\pgfpathmoveto{\pgfqpoint{1.175111in}{1.196368in}}%
\pgfpathlineto{\pgfqpoint{1.066981in}{1.224792in}}%
\pgfpathlineto{\pgfqpoint{1.154599in}{1.294241in}}%
\pgfpathlineto{\pgfqpoint{1.175111in}{1.196368in}}%
\pgfpathclose%
\pgfusepath{stroke,fill}%
\end{pgfscope}%
\begin{pgfscope}%
\pgfsetroundcap%
\pgfsetroundjoin%
\pgfsetlinewidth{1.003750pt}%
\definecolor{currentstroke}{rgb}{0.172549,0.627451,0.172549}%
\pgfsetstrokecolor{currentstroke}%
\pgfsetstrokeopacity{0.200000}%
\pgfsetdash{}{0pt}%
\pgfpathmoveto{\pgfqpoint{1.337029in}{0.950562in}}%
\pgfpathquadraticcurveto{\pgfqpoint{1.324331in}{0.951049in}}{\pgfqpoint{1.327149in}{0.950941in}}%
\pgfusepath{stroke}%
\end{pgfscope}%
\begin{pgfscope}%
\pgfsetroundcap%
\pgfsetroundjoin%
\definecolor{currentfill}{rgb}{0.172549,0.627451,0.172549}%
\pgfsetfillcolor{currentfill}%
\pgfsetfillopacity{0.200000}%
\pgfsetlinewidth{1.003750pt}%
\definecolor{currentstroke}{rgb}{0.172549,0.627451,0.172549}%
\pgfsetstrokecolor{currentstroke}%
\pgfsetstrokeopacity{0.200000}%
\pgfsetdash{}{0pt}%
\pgfpathmoveto{\pgfqpoint{1.425160in}{0.897146in}}%
\pgfpathlineto{\pgfqpoint{1.327149in}{0.950941in}}%
\pgfpathlineto{\pgfqpoint{1.428991in}{0.997073in}}%
\pgfpathlineto{\pgfqpoint{1.425160in}{0.897146in}}%
\pgfpathclose%
\pgfusepath{stroke,fill}%
\end{pgfscope}%
\begin{pgfscope}%
\pgfsetroundcap%
\pgfsetroundjoin%
\pgfsetlinewidth{1.003750pt}%
\definecolor{currentstroke}{rgb}{0.172549,0.627451,0.172549}%
\pgfsetstrokecolor{currentstroke}%
\pgfsetstrokeopacity{0.200000}%
\pgfsetdash{}{0pt}%
\pgfpathmoveto{\pgfqpoint{1.663679in}{0.639517in}}%
\pgfpathquadraticcurveto{\pgfqpoint{1.654198in}{0.641566in}}{\pgfqpoint{1.659896in}{0.640335in}}%
\pgfusepath{stroke}%
\end{pgfscope}%
\begin{pgfscope}%
\pgfsetroundcap%
\pgfsetroundjoin%
\definecolor{currentfill}{rgb}{0.172549,0.627451,0.172549}%
\pgfsetfillcolor{currentfill}%
\pgfsetfillopacity{0.200000}%
\pgfsetlinewidth{1.003750pt}%
\definecolor{currentstroke}{rgb}{0.172549,0.627451,0.172549}%
\pgfsetstrokecolor{currentstroke}%
\pgfsetstrokeopacity{0.200000}%
\pgfsetdash{}{0pt}%
\pgfpathmoveto{\pgfqpoint{1.747078in}{0.570340in}}%
\pgfpathlineto{\pgfqpoint{1.659896in}{0.640335in}}%
\pgfpathlineto{\pgfqpoint{1.768201in}{0.668083in}}%
\pgfpathlineto{\pgfqpoint{1.747078in}{0.570340in}}%
\pgfpathclose%
\pgfusepath{stroke,fill}%
\end{pgfscope}%
\begin{pgfscope}%
\pgfsetroundcap%
\pgfsetroundjoin%
\pgfsetlinewidth{1.003750pt}%
\definecolor{currentstroke}{rgb}{0.172549,0.627451,0.172549}%
\pgfsetstrokecolor{currentstroke}%
\pgfsetstrokeopacity{0.200000}%
\pgfsetdash{}{0pt}%
\pgfpathmoveto{\pgfqpoint{0.887569in}{0.551478in}}%
\pgfpathquadraticcurveto{\pgfqpoint{0.919098in}{0.550034in}}{\pgfqpoint{0.935115in}{0.549300in}}%
\pgfusepath{stroke}%
\end{pgfscope}%
\begin{pgfscope}%
\pgfsetroundcap%
\pgfsetroundjoin%
\definecolor{currentfill}{rgb}{0.172549,0.627451,0.172549}%
\pgfsetfillcolor{currentfill}%
\pgfsetfillopacity{0.200000}%
\pgfsetlinewidth{1.003750pt}%
\definecolor{currentstroke}{rgb}{0.172549,0.627451,0.172549}%
\pgfsetstrokecolor{currentstroke}%
\pgfsetstrokeopacity{0.200000}%
\pgfsetdash{}{0pt}%
\pgfpathmoveto{\pgfqpoint{0.837508in}{0.603823in}}%
\pgfpathlineto{\pgfqpoint{0.935115in}{0.549300in}}%
\pgfpathlineto{\pgfqpoint{0.832932in}{0.503928in}}%
\pgfpathlineto{\pgfqpoint{0.837508in}{0.603823in}}%
\pgfpathclose%
\pgfusepath{stroke,fill}%
\end{pgfscope}%
\begin{pgfscope}%
\pgfsetbuttcap%
\pgfsetroundjoin%
\pgfsetlinewidth{0.702625pt}%
\definecolor{currentstroke}{rgb}{0.121569,0.466667,0.705882}%
\pgfsetstrokecolor{currentstroke}%
\pgfsetdash{{2.590000pt}{1.120000pt}}{0.000000pt}%
\pgfpathmoveto{\pgfqpoint{0.506574in}{0.278150in}}%
\pgfpathlineto{\pgfqpoint{0.506574in}{1.458473in}}%
\pgfusepath{stroke}%
\end{pgfscope}%
\begin{pgfscope}%
\pgfsetbuttcap%
\pgfsetroundjoin%
\pgfsetlinewidth{0.702625pt}%
\definecolor{currentstroke}{rgb}{1.000000,0.498039,0.054902}%
\pgfsetstrokecolor{currentstroke}%
\pgfsetdash{{2.590000pt}{1.120000pt}}{0.000000pt}%
\pgfpathmoveto{\pgfqpoint{1.275353in}{0.278150in}}%
\pgfpathlineto{\pgfqpoint{1.275353in}{2.166667in}}%
\pgfusepath{stroke}%
\end{pgfscope}%
\begin{pgfscope}%
\pgfsetbuttcap%
\pgfsetroundjoin%
\definecolor{currentfill}{rgb}{1.000000,1.000000,1.000000}%
\pgfsetfillcolor{currentfill}%
\pgfsetlinewidth{1.003750pt}%
\definecolor{currentstroke}{rgb}{0.000000,0.000000,0.000000}%
\pgfsetstrokecolor{currentstroke}%
\pgfsetdash{}{0pt}%
\pgfsys@defobject{currentmarker}{\pgfqpoint{-0.041667in}{-0.041667in}}{\pgfqpoint{0.041667in}{0.041667in}}{%
\pgfpathmoveto{\pgfqpoint{0.000000in}{-0.041667in}}%
\pgfpathcurveto{\pgfqpoint{0.011050in}{-0.041667in}}{\pgfqpoint{0.021649in}{-0.037276in}}{\pgfqpoint{0.029463in}{-0.029463in}}%
\pgfpathcurveto{\pgfqpoint{0.037276in}{-0.021649in}}{\pgfqpoint{0.041667in}{-0.011050in}}{\pgfqpoint{0.041667in}{0.000000in}}%
\pgfpathcurveto{\pgfqpoint{0.041667in}{0.011050in}}{\pgfqpoint{0.037276in}{0.021649in}}{\pgfqpoint{0.029463in}{0.029463in}}%
\pgfpathcurveto{\pgfqpoint{0.021649in}{0.037276in}}{\pgfqpoint{0.011050in}{0.041667in}}{\pgfqpoint{0.000000in}{0.041667in}}%
\pgfpathcurveto{\pgfqpoint{-0.011050in}{0.041667in}}{\pgfqpoint{-0.021649in}{0.037276in}}{\pgfqpoint{-0.029463in}{0.029463in}}%
\pgfpathcurveto{\pgfqpoint{-0.037276in}{0.021649in}}{\pgfqpoint{-0.041667in}{0.011050in}}{\pgfqpoint{-0.041667in}{0.000000in}}%
\pgfpathcurveto{\pgfqpoint{-0.041667in}{-0.011050in}}{\pgfqpoint{-0.037276in}{-0.021649in}}{\pgfqpoint{-0.029463in}{-0.029463in}}%
\pgfpathcurveto{\pgfqpoint{-0.021649in}{-0.037276in}}{\pgfqpoint{-0.011050in}{-0.041667in}}{\pgfqpoint{0.000000in}{-0.041667in}}%
\pgfpathclose%
\pgfusepath{stroke,fill}%
}%
\begin{pgfscope}%
\pgfsys@transformshift{0.506574in}{0.278150in}%
\pgfsys@useobject{currentmarker}{}%
\end{pgfscope}%
\end{pgfscope}%
\begin{pgfscope}%
\pgfsetbuttcap%
\pgfsetroundjoin%
\definecolor{currentfill}{rgb}{1.000000,1.000000,1.000000}%
\pgfsetfillcolor{currentfill}%
\pgfsetlinewidth{1.003750pt}%
\definecolor{currentstroke}{rgb}{0.000000,0.000000,0.000000}%
\pgfsetstrokecolor{currentstroke}%
\pgfsetdash{}{0pt}%
\pgfsys@defobject{currentmarker}{\pgfqpoint{-0.041667in}{-0.041667in}}{\pgfqpoint{0.041667in}{0.041667in}}{%
\pgfpathmoveto{\pgfqpoint{0.000000in}{-0.041667in}}%
\pgfpathcurveto{\pgfqpoint{0.011050in}{-0.041667in}}{\pgfqpoint{0.021649in}{-0.037276in}}{\pgfqpoint{0.029463in}{-0.029463in}}%
\pgfpathcurveto{\pgfqpoint{0.037276in}{-0.021649in}}{\pgfqpoint{0.041667in}{-0.011050in}}{\pgfqpoint{0.041667in}{0.000000in}}%
\pgfpathcurveto{\pgfqpoint{0.041667in}{0.011050in}}{\pgfqpoint{0.037276in}{0.021649in}}{\pgfqpoint{0.029463in}{0.029463in}}%
\pgfpathcurveto{\pgfqpoint{0.021649in}{0.037276in}}{\pgfqpoint{0.011050in}{0.041667in}}{\pgfqpoint{0.000000in}{0.041667in}}%
\pgfpathcurveto{\pgfqpoint{-0.011050in}{0.041667in}}{\pgfqpoint{-0.021649in}{0.037276in}}{\pgfqpoint{-0.029463in}{0.029463in}}%
\pgfpathcurveto{\pgfqpoint{-0.037276in}{0.021649in}}{\pgfqpoint{-0.041667in}{0.011050in}}{\pgfqpoint{-0.041667in}{0.000000in}}%
\pgfpathcurveto{\pgfqpoint{-0.041667in}{-0.011050in}}{\pgfqpoint{-0.037276in}{-0.021649in}}{\pgfqpoint{-0.029463in}{-0.029463in}}%
\pgfpathcurveto{\pgfqpoint{-0.021649in}{-0.037276in}}{\pgfqpoint{-0.011050in}{-0.041667in}}{\pgfqpoint{0.000000in}{-0.041667in}}%
\pgfpathclose%
\pgfusepath{stroke,fill}%
}%
\begin{pgfscope}%
\pgfsys@transformshift{2.044132in}{0.278150in}%
\pgfsys@useobject{currentmarker}{}%
\end{pgfscope}%
\end{pgfscope}%
\begin{pgfscope}%
\pgfsetbuttcap%
\pgfsetroundjoin%
\definecolor{currentfill}{rgb}{0.000000,0.000000,0.000000}%
\pgfsetfillcolor{currentfill}%
\pgfsetlinewidth{1.003750pt}%
\definecolor{currentstroke}{rgb}{0.000000,0.000000,0.000000}%
\pgfsetstrokecolor{currentstroke}%
\pgfsetdash{}{0pt}%
\pgfsys@defobject{currentmarker}{\pgfqpoint{-0.041667in}{-0.041667in}}{\pgfqpoint{0.041667in}{0.041667in}}{%
\pgfpathmoveto{\pgfqpoint{0.000000in}{-0.041667in}}%
\pgfpathcurveto{\pgfqpoint{0.011050in}{-0.041667in}}{\pgfqpoint{0.021649in}{-0.037276in}}{\pgfqpoint{0.029463in}{-0.029463in}}%
\pgfpathcurveto{\pgfqpoint{0.037276in}{-0.021649in}}{\pgfqpoint{0.041667in}{-0.011050in}}{\pgfqpoint{0.041667in}{0.000000in}}%
\pgfpathcurveto{\pgfqpoint{0.041667in}{0.011050in}}{\pgfqpoint{0.037276in}{0.021649in}}{\pgfqpoint{0.029463in}{0.029463in}}%
\pgfpathcurveto{\pgfqpoint{0.021649in}{0.037276in}}{\pgfqpoint{0.011050in}{0.041667in}}{\pgfqpoint{0.000000in}{0.041667in}}%
\pgfpathcurveto{\pgfqpoint{-0.011050in}{0.041667in}}{\pgfqpoint{-0.021649in}{0.037276in}}{\pgfqpoint{-0.029463in}{0.029463in}}%
\pgfpathcurveto{\pgfqpoint{-0.037276in}{0.021649in}}{\pgfqpoint{-0.041667in}{0.011050in}}{\pgfqpoint{-0.041667in}{0.000000in}}%
\pgfpathcurveto{\pgfqpoint{-0.041667in}{-0.011050in}}{\pgfqpoint{-0.037276in}{-0.021649in}}{\pgfqpoint{-0.029463in}{-0.029463in}}%
\pgfpathcurveto{\pgfqpoint{-0.021649in}{-0.037276in}}{\pgfqpoint{-0.011050in}{-0.041667in}}{\pgfqpoint{0.000000in}{-0.041667in}}%
\pgfpathclose%
\pgfusepath{stroke,fill}%
}%
\begin{pgfscope}%
\pgfsys@transformshift{1.275353in}{0.868311in}%
\pgfsys@useobject{currentmarker}{}%
\end{pgfscope}%
\end{pgfscope}%
\end{pgfpicture}%
\makeatother%
\endgroup%

          \end{center}
          where the case with $(\beta,1-\beta)$ as a stable spiral is on the left, and the case with $(\beta,1-\beta)$ as a stable node is on the right. Note on the right that the trajectory from $(1,0)$ to $(\beta,1-\beta)$ does not go directly along the nullcline: if you zoom in, you can see it is slightly to the right (which makes sense, given that trajectories cross the nullcline vertically).
          %
          \commentary{Drawing the nullclines was done fine, but the trajectories in general were drawn very badly. We covered this extensively in Problems Class 2: if you are not showing up to problems classes, consider this your warning! Common errors leading to lost marks included:
            \begin{enumerate}[topsep=0pt,itemsep=0pt]
              \item Not drawing trajectories as continuous lines. This is not like drawing a vector field in AMV or a velocity field in Fluids, where you just draw little arrows to represent the direction/magnitude at some points: trajectories are paths and they have to start somewhere and end either at an equilibrium, or go off to infinity.
              \item{Not having the trajectory curves cross the nullclines either vertically or horizontally. This is \emph{crucial} to getting the right shape and you will lose marks in an exam if you don't draw these accurately.}
              \item{Having trajectory curves change either vertical or horizontal direction away from a nullcline. This is not allowed; that is the point of the nullclines: the locations where they change direction.
                }
              \item Not drawing a diagram that is big enough to contain all the detail you need. Go big or go home! In some cases, this led people to drawing trajectories which cross, which, you will remember from Problems Class 2, is not allowed.
              \item{Not observing there are two possibilities. The question gives you a clear steer on this by highlighting the two: that means two diagrams. The local structure of the linear system (node or spiral) has to dictate how you draw the curves near the equilibrium. }
          \end{enumerate}}
          %
        \item Since $(0,0)$ is a saddle point, there is only one with trajectory leaving it with $u>0$ and we can see that it corresponds to the horizontal nullcline, i.e.\ trajectories leave $(0,0)$ and go to $(1,0)$. Therefore, there is no direct trajectory from $(0,0)$ to $(\beta,1-\beta)$.
          %
          \commentary{There were quite a lot of misunderstandings about trajectories. A common error was saying that because you had seen the trajectories spiral in that a path exists. But, there are a couple of issues. First a travelling wave, the way we define it in class, has to go from an equilibrium to an equilibrium. If you are at $(0,0)$ the only way out is along the saddle arc, which is along $v=0$. Since this is a nullcline, $\fd{v}{t}=0$ always, and so the trajectory will end at $(1,0)$. It is fine to observe that in the sector between $(0,0)$ and $(\beta,1-\beta)$ the trajectories have to head downwards, so you can't go towards $(\beta,1-\beta)$ anyway. But then you have to explain why the trajectories can't leave the $u$-axis after the line $u=\beta$.}
      \end{enumerate}
    }

    % Q2
    \question{}{ % Question
      We now look for travelling wave solutions to the full system from $(1,0)$ to $(\beta,1-\beta)$. We place the stable equilibrium at $x=-\infty$ and the unstable equilibrium at $x=\infty$, and look for rightwards-travelling waves of the form $z=x-ct$, with $c>0$.
      \begin{enumerate}
        \item If we instead placed the stable equilibrium at $x=\infty$ and the unstable equilibrium at $x=-\infty$, what form of $z$ would we use?
        \item Show that, with $z=x-ct$ and $c>0$, and by introducing $w=\fd{v}{z}$, we arrive at the three-dimensional system,
          \begin{align}
            \fd{u}{z} &= \frac{u}{c}(u+v-1), \tag{\xmas{reindeer}}\\
            \fd{v}{z} &= w, \tag{\xmas{snowflake}}\\
            \fd{w}{z} &= \alpha v(\beta - u) - cw, \tag{\xmas{bauble}}
          \end{align}
          and write down the boundary conditions as $z \to \pm \infty$.
      \end{enumerate}
    }
    { % Solution
      \begin{enumerate}
        \item If the stable equilibrium is at $x=\infty$, we would expect this region to expand and therefore the wave to travel \emph{leftwards}, i.e.\ $z=x+ct$ with $c>0$.
          %
          \commentary{The main error was switching $x$ to $-x$. The locations of the particular equilibria are fixed in the question, so you can't swap: it has to be $c$ which changes sign.}
          %
        \item We look for solutions in $z$ where
          \begin{equation}
            u(x,t) = u(z), \quad v(x,t) = v(z), \quad z = x-ct.
          \end{equation}

          Substituting $z$ into (\xmas{santa}\xmas{christmas-tree}), we obtain
          \begin{subequations}
            \begin{align}
              -c\fd{u}{z} &= u(1-u-v),\\
              -c\fd{v}{z} &= \fdd{v}{z} + \alpha v (u-\beta).
            \end{align}
          \end{subequations}
          Introducing $w = \fd{v}{z}$ and rearranging, we arrive at the three-dimensional system
          \begin{subequations}
            \begin{align}
              \fd{u}{z} &= \frac{u}{c}(u+v-1),\\
              \fd{v}{z} &= w,\\
              \fd{w}{z} &= \alpha v(\beta - u) - cw.
            \end{align}
          \end{subequations}
          The boundary conditions are
          \begin{align}
            \label{trav-wave-interesting}
            &(u,v,w) \to (\beta,1-\beta,0) \text{ as } z \to -\infty, \\
            &(u,v,w) \to (1,0,0) \text{ as } z \to \infty.
          \end{align}
          %(We have that $w$ is 0 because these two boundary conditions are equilibria, i.e.\ when time derivatives are zero, and $w=\fd{v}{t}$.)
          %
          \commentary{Two main errors:
        \begin{enumerate}[itemsep=0pt,topsep=0pt]
        \item Putting the stable equilibrium at $z\rightarrow \infty$: the movement should be away from that as it slowly takes over.
        \item Only giving one boundary condition on $u$. You need to also give a boundary condition on $w$.
        \end{enumerate}}
      \end{enumerate}

    }

    % Q3
    \question{}{ % Questoon
      We are now going to perform a linear stability analysis on the equilibrium at $(u,v,w)=(1,0,0)$.
      \begin{enumerate}
        \item Write down the $3\times 3$ Jacobian of the system (\xmas{reindeer}\xmas{snowflake}\xmas{bauble}).
        \item Find the eigenvalues evaluated at $(1,0,0)$.
        \item Make an argument for why we cannot have complex eigenvalues and thus derive a constraint on the squared wavespeed, $c^2$.\label{c-constraint}
      \end{enumerate}
    }
    { % Solution
      \begin{enumerate}
        \item The Jacobian of the system (\xmas{reindeer}\xmas{snowflake}\xmas{bauble}) is
          \begin{align}
            \m{J} =
            \begin{pmatrix}
              (2u+v-1)/c & u/c & 0 \\
              0 & 0 & 1 \\
              -\alpha v & \alpha(\beta-u) & -c
            \end{pmatrix}.
          \end{align}
          %
        \item The Jacobian, evaluated at $(1,0,0)$ gives us
          \begin{align}
            \m{J} =
            \begin{pmatrix}
              1/c & 1/c & 0 \\
              0 & 0 & 1 \\
              0 & \alpha(\beta-1) & -c
            \end{pmatrix}.
          \end{align}
          This matrix leads to the eigenvalue equation
          \begin{equation}
            (\lambda-1/c)[\lambda(c+\lambda)-\alpha(\beta-1)] = 0,
          \end{equation}
          and hence the eigenvalues are
          \begin{equation}
            \lambda = \frac{1}{c}, \quad \lambda = \frac{1}{2}\left( -c \pm \sqrt{c^2 + 4\alpha(\beta-1)} \right).
          \end{equation}
          %
          \commentary{Some small errors in the matrix calculation; otherwise this was done well.}
          %
        \item The first eigenvalue is positive so we expect an unstable solution. But what type of unstable solution? From the second expression, we see that the solution will be oscillatory (a spiral) if $c^2 < 4\alpha(1-\beta)$. Since $v$ cannot be less than zero, these oscillatory solutions can be discarded and this gives us a constraint on the wavespeed, just like it did with the FK equation. Namely, the permissible wave speeds satisfy
          \begin{equation}
            c^2 \geq 4\alpha(1-\beta).
          \end{equation}
          %
          \commentary{Quite a few people missed the key point here. A travelling wave is really just a solution in the form $u(x-ct)$, with the linear coupling of the two independent variables. Its name comes from the wave equation,
    \begin{equation}
    	\pdd{u}{t} = c \pdd{u}{x},
    \end{equation}
    where any $f(x-ct)$ solves the system, whatever the shape of $f$. So there are no issues with the oscillatory/spiral shape in principle (mathematically). But on \emph{physical} grounds, spirals imply a negative population of $v$ as the spiral is around $(1,0,0)$ (so you should say that $v$ is the problem). This was the same argument we used in class for the FK equation.}
      \end{enumerate}

    }

    % Q4
    \question{}{ % Question
      We now look at the stability of the other equilibrium, $(\beta,1-\beta,0)$.
      \begin{enumerate}
        \item Evaluate the Jacobian at this equilibrium and show that the resulting eigenvalue equation is
          \begin{equation}
            \lambda^3 + \lambda^2\left(c-\frac{\beta}{c}\right) - \beta \lambda + \frac{\beta}{c}\alpha(1-\beta) = 0. \tag{\xmas{candle}}
            \label{candle}
          \end{equation}
        \item This is a positive cubic polynomial, so we might expect there to be a local minimum and maximum. Sketch the polynomial and show that the local minimum and maximum are at
          \begin{equation}
            \lambda_{\text{max}},\lambda_{\text{min}} = -\frac{1}{3}\left[c - \frac{\beta}{c} \pm \sqrt{\left( c- \frac{\beta}{c} \right)^2 + \xmas{snowman}} \right],
            \tag{\xmas{gingerbread-man}}
          \end{equation}
          where you should state the value of \xmas{snowman}. (Just to reiterate: these are \emph{not} necessarily eigenvalues, just the local maximum and minimum of the eigenvalue equation.)
        \item See that these values of the local maximum and minimum of the eigenvalue equation do not depend on the parameter $\alpha$. This suggests that increasing $\alpha$ just corresponds to shifting the cubic upwards. Convince yourself that this is the case. Are you convinced?
        \item Consider \cref{candle}, then, with $\alpha=0$ (its lowest possible value), and find the three real eigenvalues which satisfy this equation.
        \item Therefore we have that when $\alpha=0$, there are three real eigenvalues. As we increase $\alpha$, the cubic graph shifts upwards until $\alpha$ reaches some critical value $\alpha_c$, when the local minimum of the cubic goes above the $x$-axis, after which there is one real eigenvalue and two complex eigenvalues. In the squirrel (FK) problem, we disallowed complex eigenvalues. Should we also disallow them here?
      \end{enumerate}
      We haven't given names to types of stability in three dimensions; suffice to say that if $\alpha \leq \alpha_c$, we expect $(u(z),v(z))$ to vary monotonically with $z$, going from $(\beta, 1-\beta)$ to $(1,0)$ as $z$ increases. And if $\alpha > \alpha_c$, we expect the solution to oscillate as $z$ increases away from the equilibrium at $(\beta,1-\beta)$. The solutions look like \cref{fig:sols}.

      As for the wavespeed, this equilibrium puts no constraints on it. Coupled with the rule that waves travel at their slowest possible value, $c^2$ turns out to be equal to the value you specified in question (\ref{c-constraint}).
    }
    { % Solution
      \begin{enumerate}
        \item The Jacobian, evaluated at $(\beta,1-\beta,0)$ is
          \begin{align}
            \m{J} =
            \begin{pmatrix}
              \beta/c & \beta/c & 0 \\
              0 & 0 & 1 \\
              \alpha(\beta-1) & 0 & -c
            \end{pmatrix},
          \end{align}
          which has eigenvalue equation
          \begin{equation}
            \label{main}
            \lambda^3 + \lambda^2\left(c-\frac{\beta}{c}\right) - \beta \lambda + \frac{\beta}{c}\alpha(1-\beta) = 0.
          \end{equation}
          %
          \commentary{This was done fine.}
          %
        \item
          The cubic looks like
          \begin{center}
            %% Creator: Matplotlib, PGF backend
%%
%% To include the figure in your LaTeX document, write
%%   \input{<filename>.pgf}
%%
%% Make sure the required packages are loaded in your preamble
%%   \usepackage{pgf}
%%
%% and, on pdftex
%%   \usepackage[utf8]{inputenc}\DeclareUnicodeCharacter{2212}{-}
%%
%% or, on luatex and xetex
%%   \usepackage{unicode-math}
%%
%% Figures using additional raster images can only be included by \input if
%% they are in the same directory as the main LaTeX file. For loading figures
%% from other directories you can use the `import` package
%%   \usepackage{import}
%%
%% and then include the figures with
%%   \import{<path to file>}{<filename>.pgf}
%%
%% Matplotlib used the following preamble
%%   \usepackage{fontspec}
%%   \setmainfont{DejaVuSerif.ttf}[Path=/Users/adam/opt/anaconda3/lib/python3.8/site-packages/matplotlib/mpl-data/fonts/ttf/]
%%   \setsansfont{DejaVuSans.ttf}[Path=/Users/adam/opt/anaconda3/lib/python3.8/site-packages/matplotlib/mpl-data/fonts/ttf/]
%%   \setmonofont{DejaVuSansMono.ttf}[Path=/Users/adam/opt/anaconda3/lib/python3.8/site-packages/matplotlib/mpl-data/fonts/ttf/]
%%
\begingroup%
\makeatletter%
\begin{pgfpicture}%
\pgfpathrectangle{\pgfpointorigin}{\pgfqpoint{3.000000in}{2.200000in}}%
\pgfusepath{use as bounding box, clip}%
\begin{pgfscope}%
\pgfsetbuttcap%
\pgfsetmiterjoin%
\definecolor{currentfill}{rgb}{1.000000,1.000000,1.000000}%
\pgfsetfillcolor{currentfill}%
\pgfsetlinewidth{0.000000pt}%
\definecolor{currentstroke}{rgb}{1.000000,1.000000,1.000000}%
\pgfsetstrokecolor{currentstroke}%
\pgfsetdash{}{0pt}%
\pgfpathmoveto{\pgfqpoint{0.000000in}{0.000000in}}%
\pgfpathlineto{\pgfqpoint{3.000000in}{0.000000in}}%
\pgfpathlineto{\pgfqpoint{3.000000in}{2.200000in}}%
\pgfpathlineto{\pgfqpoint{0.000000in}{2.200000in}}%
\pgfpathclose%
\pgfusepath{fill}%
\end{pgfscope}%
\begin{pgfscope}%
\pgfsetbuttcap%
\pgfsetmiterjoin%
\definecolor{currentfill}{rgb}{1.000000,1.000000,1.000000}%
\pgfsetfillcolor{currentfill}%
\pgfsetlinewidth{0.000000pt}%
\definecolor{currentstroke}{rgb}{0.000000,0.000000,0.000000}%
\pgfsetstrokecolor{currentstroke}%
\pgfsetstrokeopacity{0.000000}%
\pgfsetdash{}{0pt}%
\pgfpathmoveto{\pgfqpoint{0.033333in}{0.033333in}}%
\pgfpathlineto{\pgfqpoint{2.966667in}{0.033333in}}%
\pgfpathlineto{\pgfqpoint{2.966667in}{2.166667in}}%
\pgfpathlineto{\pgfqpoint{0.033333in}{2.166667in}}%
\pgfpathclose%
\pgfusepath{fill}%
\end{pgfscope}%
\begin{pgfscope}%
\pgfsetbuttcap%
\pgfsetroundjoin%
\definecolor{currentfill}{rgb}{0.000000,0.000000,0.000000}%
\pgfsetfillcolor{currentfill}%
\pgfsetlinewidth{0.803000pt}%
\definecolor{currentstroke}{rgb}{0.000000,0.000000,0.000000}%
\pgfsetstrokecolor{currentstroke}%
\pgfsetdash{}{0pt}%
\pgfsys@defobject{currentmarker}{\pgfqpoint{0.000000in}{-0.048611in}}{\pgfqpoint{0.000000in}{0.000000in}}{%
\pgfpathmoveto{\pgfqpoint{0.000000in}{0.000000in}}%
\pgfpathlineto{\pgfqpoint{0.000000in}{-0.048611in}}%
\pgfusepath{stroke,fill}%
}%
\begin{pgfscope}%
\pgfsys@transformshift{1.500000in}{1.100000in}%
\pgfsys@useobject{currentmarker}{}%
\end{pgfscope}%
\end{pgfscope}%
\begin{pgfscope}%
\definecolor{textcolor}{rgb}{0.000000,0.000000,0.000000}%
\pgfsetstrokecolor{textcolor}%
\pgfsetfillcolor{textcolor}%
\pgftext[x=1.500000in,y=1.002778in,,top]{\color{textcolor}\rmfamily\fontsize{12.000000}{14.400000}\selectfont 0}%
\end{pgfscope}%
\begin{pgfscope}%
\pgfsetbuttcap%
\pgfsetroundjoin%
\definecolor{currentfill}{rgb}{0.000000,0.000000,0.000000}%
\pgfsetfillcolor{currentfill}%
\pgfsetlinewidth{0.803000pt}%
\definecolor{currentstroke}{rgb}{0.000000,0.000000,0.000000}%
\pgfsetstrokecolor{currentstroke}%
\pgfsetdash{}{0pt}%
\pgfsys@defobject{currentmarker}{\pgfqpoint{0.000000in}{-0.048611in}}{\pgfqpoint{0.000000in}{0.000000in}}{%
\pgfpathmoveto{\pgfqpoint{0.000000in}{0.000000in}}%
\pgfpathlineto{\pgfqpoint{0.000000in}{-0.048611in}}%
\pgfusepath{stroke,fill}%
}%
\begin{pgfscope}%
\pgfsys@transformshift{0.884160in}{1.100000in}%
\pgfsys@useobject{currentmarker}{}%
\end{pgfscope}%
\end{pgfscope}%
\begin{pgfscope}%
\definecolor{textcolor}{rgb}{0.000000,0.000000,0.000000}%
\pgfsetstrokecolor{textcolor}%
\pgfsetfillcolor{textcolor}%
\pgftext[x=0.884160in,y=1.002778in,,top]{\color{textcolor}\rmfamily\fontsize{12.000000}{14.400000}\selectfont \(\displaystyle -c\)}%
\end{pgfscope}%
\begin{pgfscope}%
\pgfsetbuttcap%
\pgfsetroundjoin%
\definecolor{currentfill}{rgb}{0.000000,0.000000,0.000000}%
\pgfsetfillcolor{currentfill}%
\pgfsetlinewidth{0.803000pt}%
\definecolor{currentstroke}{rgb}{0.000000,0.000000,0.000000}%
\pgfsetstrokecolor{currentstroke}%
\pgfsetdash{}{0pt}%
\pgfsys@defobject{currentmarker}{\pgfqpoint{0.000000in}{-0.048611in}}{\pgfqpoint{0.000000in}{0.000000in}}{%
\pgfpathmoveto{\pgfqpoint{0.000000in}{0.000000in}}%
\pgfpathlineto{\pgfqpoint{0.000000in}{-0.048611in}}%
\pgfusepath{stroke,fill}%
}%
\begin{pgfscope}%
\pgfsys@transformshift{2.013200in}{1.100000in}%
\pgfsys@useobject{currentmarker}{}%
\end{pgfscope}%
\end{pgfscope}%
\begin{pgfscope}%
\definecolor{textcolor}{rgb}{0.000000,0.000000,0.000000}%
\pgfsetstrokecolor{textcolor}%
\pgfsetfillcolor{textcolor}%
\pgftext[x=2.013200in,y=1.002778in,,top]{\color{textcolor}\rmfamily\fontsize{12.000000}{14.400000}\selectfont \(\displaystyle \beta/c\)}%
\end{pgfscope}%
\begin{pgfscope}%
\pgfsetbuttcap%
\pgfsetroundjoin%
\definecolor{currentfill}{rgb}{0.000000,0.000000,0.000000}%
\pgfsetfillcolor{currentfill}%
\pgfsetlinewidth{0.803000pt}%
\definecolor{currentstroke}{rgb}{0.000000,0.000000,0.000000}%
\pgfsetstrokecolor{currentstroke}%
\pgfsetdash{}{0pt}%
\pgfsys@defobject{currentmarker}{\pgfqpoint{-0.048611in}{0.000000in}}{\pgfqpoint{0.000000in}{0.000000in}}{%
\pgfpathmoveto{\pgfqpoint{0.000000in}{0.000000in}}%
\pgfpathlineto{\pgfqpoint{-0.048611in}{0.000000in}}%
\pgfusepath{stroke,fill}%
}%
\begin{pgfscope}%
\pgfsys@transformshift{1.500000in}{1.100000in}%
\pgfsys@useobject{currentmarker}{}%
\end{pgfscope}%
\end{pgfscope}%
\begin{pgfscope}%
\definecolor{textcolor}{rgb}{0.000000,0.000000,0.000000}%
\pgfsetstrokecolor{textcolor}%
\pgfsetfillcolor{textcolor}%
\pgftext[x=1.296739in, y=1.036686in, left, base]{\color{textcolor}\rmfamily\fontsize{12.000000}{14.400000}\selectfont 0}%
\end{pgfscope}%
\begin{pgfscope}%
\pgfpathrectangle{\pgfqpoint{0.033333in}{0.033333in}}{\pgfqpoint{2.933333in}{2.133333in}}%
\pgfusepath{clip}%
\pgfsetrectcap%
\pgfsetroundjoin%
\pgfsetlinewidth{1.505625pt}%
\definecolor{currentstroke}{rgb}{0.121569,0.466667,0.705882}%
\pgfsetstrokecolor{currentstroke}%
\pgfsetdash{}{0pt}%
\pgfpathmoveto{\pgfqpoint{0.604676in}{0.028333in}}%
\pgfpathlineto{\pgfqpoint{0.624579in}{0.141507in}}%
\pgfpathlineto{\pgfqpoint{0.651515in}{0.283549in}}%
\pgfpathlineto{\pgfqpoint{0.678451in}{0.414825in}}%
\pgfpathlineto{\pgfqpoint{0.705387in}{0.535689in}}%
\pgfpathlineto{\pgfqpoint{0.732323in}{0.646499in}}%
\pgfpathlineto{\pgfqpoint{0.759259in}{0.747610in}}%
\pgfpathlineto{\pgfqpoint{0.786195in}{0.839378in}}%
\pgfpathlineto{\pgfqpoint{0.813131in}{0.922161in}}%
\pgfpathlineto{\pgfqpoint{0.840067in}{0.996312in}}%
\pgfpathlineto{\pgfqpoint{0.867003in}{1.062190in}}%
\pgfpathlineto{\pgfqpoint{0.893939in}{1.120151in}}%
\pgfpathlineto{\pgfqpoint{0.920875in}{1.170549in}}%
\pgfpathlineto{\pgfqpoint{0.947811in}{1.213742in}}%
\pgfpathlineto{\pgfqpoint{0.974747in}{1.250086in}}%
\pgfpathlineto{\pgfqpoint{1.001684in}{1.279937in}}%
\pgfpathlineto{\pgfqpoint{1.028620in}{1.303651in}}%
\pgfpathlineto{\pgfqpoint{1.055556in}{1.321584in}}%
\pgfpathlineto{\pgfqpoint{1.082492in}{1.334093in}}%
\pgfpathlineto{\pgfqpoint{1.109428in}{1.341533in}}%
\pgfpathlineto{\pgfqpoint{1.136364in}{1.344261in}}%
\pgfpathlineto{\pgfqpoint{1.163300in}{1.342633in}}%
\pgfpathlineto{\pgfqpoint{1.190236in}{1.337005in}}%
\pgfpathlineto{\pgfqpoint{1.217172in}{1.327734in}}%
\pgfpathlineto{\pgfqpoint{1.244108in}{1.315175in}}%
\pgfpathlineto{\pgfqpoint{1.271044in}{1.299685in}}%
\pgfpathlineto{\pgfqpoint{1.297980in}{1.281620in}}%
\pgfpathlineto{\pgfqpoint{1.324916in}{1.261335in}}%
\pgfpathlineto{\pgfqpoint{1.351852in}{1.239188in}}%
\pgfpathlineto{\pgfqpoint{1.378788in}{1.215535in}}%
\pgfpathlineto{\pgfqpoint{1.405724in}{1.190731in}}%
\pgfpathlineto{\pgfqpoint{1.432660in}{1.165133in}}%
\pgfpathlineto{\pgfqpoint{1.459596in}{1.139096in}}%
\pgfpathlineto{\pgfqpoint{1.486532in}{1.112978in}}%
\pgfpathlineto{\pgfqpoint{1.513468in}{1.087135in}}%
\pgfpathlineto{\pgfqpoint{1.540404in}{1.061921in}}%
\pgfpathlineto{\pgfqpoint{1.567340in}{1.037695in}}%
\pgfpathlineto{\pgfqpoint{1.594276in}{1.014811in}}%
\pgfpathlineto{\pgfqpoint{1.621212in}{0.993626in}}%
\pgfpathlineto{\pgfqpoint{1.648148in}{0.974497in}}%
\pgfpathlineto{\pgfqpoint{1.675084in}{0.957779in}}%
\pgfpathlineto{\pgfqpoint{1.702020in}{0.943828in}}%
\pgfpathlineto{\pgfqpoint{1.728956in}{0.933002in}}%
\pgfpathlineto{\pgfqpoint{1.755892in}{0.925655in}}%
\pgfpathlineto{\pgfqpoint{1.782828in}{0.922144in}}%
\pgfpathlineto{\pgfqpoint{1.809764in}{0.922826in}}%
\pgfpathlineto{\pgfqpoint{1.836700in}{0.928056in}}%
\pgfpathlineto{\pgfqpoint{1.863636in}{0.938190in}}%
\pgfpathlineto{\pgfqpoint{1.890572in}{0.953586in}}%
\pgfpathlineto{\pgfqpoint{1.917508in}{0.974598in}}%
\pgfpathlineto{\pgfqpoint{1.944444in}{1.001584in}}%
\pgfpathlineto{\pgfqpoint{1.971380in}{1.034899in}}%
\pgfpathlineto{\pgfqpoint{1.998316in}{1.074900in}}%
\pgfpathlineto{\pgfqpoint{2.025253in}{1.121942in}}%
\pgfpathlineto{\pgfqpoint{2.052189in}{1.176382in}}%
\pgfpathlineto{\pgfqpoint{2.079125in}{1.238576in}}%
\pgfpathlineto{\pgfqpoint{2.106061in}{1.308881in}}%
\pgfpathlineto{\pgfqpoint{2.132997in}{1.387652in}}%
\pgfpathlineto{\pgfqpoint{2.159933in}{1.475245in}}%
\pgfpathlineto{\pgfqpoint{2.186869in}{1.572018in}}%
\pgfpathlineto{\pgfqpoint{2.213805in}{1.678325in}}%
\pgfpathlineto{\pgfqpoint{2.240741in}{1.794524in}}%
\pgfpathlineto{\pgfqpoint{2.267677in}{1.920969in}}%
\pgfpathlineto{\pgfqpoint{2.294613in}{2.058019in}}%
\pgfpathlineto{\pgfqpoint{2.315296in}{2.171667in}}%
\pgfusepath{stroke}%
\end{pgfscope}%
\begin{pgfscope}%
\pgfpathrectangle{\pgfqpoint{0.033333in}{0.033333in}}{\pgfqpoint{2.933333in}{2.133333in}}%
\pgfusepath{clip}%
\pgfsetrectcap%
\pgfsetroundjoin%
\pgfsetlinewidth{1.505625pt}%
\definecolor{currentstroke}{rgb}{1.000000,0.498039,0.054902}%
\pgfsetstrokecolor{currentstroke}%
\pgfsetdash{}{0pt}%
\pgfpathmoveto{\pgfqpoint{0.580006in}{0.028333in}}%
\pgfpathlineto{\pgfqpoint{0.597643in}{0.136142in}}%
\pgfpathlineto{\pgfqpoint{0.624579in}{0.289308in}}%
\pgfpathlineto{\pgfqpoint{0.651515in}{0.431351in}}%
\pgfpathlineto{\pgfqpoint{0.678451in}{0.562627in}}%
\pgfpathlineto{\pgfqpoint{0.705387in}{0.683491in}}%
\pgfpathlineto{\pgfqpoint{0.732323in}{0.794301in}}%
\pgfpathlineto{\pgfqpoint{0.759259in}{0.895412in}}%
\pgfpathlineto{\pgfqpoint{0.786195in}{0.987180in}}%
\pgfpathlineto{\pgfqpoint{0.813131in}{1.069962in}}%
\pgfpathlineto{\pgfqpoint{0.840067in}{1.144114in}}%
\pgfpathlineto{\pgfqpoint{0.867003in}{1.209992in}}%
\pgfpathlineto{\pgfqpoint{0.893939in}{1.267952in}}%
\pgfpathlineto{\pgfqpoint{0.920875in}{1.318351in}}%
\pgfpathlineto{\pgfqpoint{0.947811in}{1.361544in}}%
\pgfpathlineto{\pgfqpoint{0.974747in}{1.397888in}}%
\pgfpathlineto{\pgfqpoint{1.001684in}{1.427739in}}%
\pgfpathlineto{\pgfqpoint{1.028620in}{1.451453in}}%
\pgfpathlineto{\pgfqpoint{1.055556in}{1.469386in}}%
\pgfpathlineto{\pgfqpoint{1.082492in}{1.481894in}}%
\pgfpathlineto{\pgfqpoint{1.109428in}{1.489335in}}%
\pgfpathlineto{\pgfqpoint{1.136364in}{1.492063in}}%
\pgfpathlineto{\pgfqpoint{1.163300in}{1.490435in}}%
\pgfpathlineto{\pgfqpoint{1.190236in}{1.484807in}}%
\pgfpathlineto{\pgfqpoint{1.217172in}{1.475536in}}%
\pgfpathlineto{\pgfqpoint{1.244108in}{1.462977in}}%
\pgfpathlineto{\pgfqpoint{1.271044in}{1.447486in}}%
\pgfpathlineto{\pgfqpoint{1.297980in}{1.429421in}}%
\pgfpathlineto{\pgfqpoint{1.324916in}{1.409137in}}%
\pgfpathlineto{\pgfqpoint{1.351852in}{1.386990in}}%
\pgfpathlineto{\pgfqpoint{1.378788in}{1.363336in}}%
\pgfpathlineto{\pgfqpoint{1.405724in}{1.338533in}}%
\pgfpathlineto{\pgfqpoint{1.432660in}{1.312934in}}%
\pgfpathlineto{\pgfqpoint{1.459596in}{1.286898in}}%
\pgfpathlineto{\pgfqpoint{1.486532in}{1.260780in}}%
\pgfpathlineto{\pgfqpoint{1.513468in}{1.234936in}}%
\pgfpathlineto{\pgfqpoint{1.540404in}{1.209723in}}%
\pgfpathlineto{\pgfqpoint{1.567340in}{1.185497in}}%
\pgfpathlineto{\pgfqpoint{1.594276in}{1.162613in}}%
\pgfpathlineto{\pgfqpoint{1.621212in}{1.141428in}}%
\pgfpathlineto{\pgfqpoint{1.648148in}{1.122299in}}%
\pgfpathlineto{\pgfqpoint{1.675084in}{1.105581in}}%
\pgfpathlineto{\pgfqpoint{1.702020in}{1.091630in}}%
\pgfpathlineto{\pgfqpoint{1.728956in}{1.080803in}}%
\pgfpathlineto{\pgfqpoint{1.755892in}{1.073456in}}%
\pgfpathlineto{\pgfqpoint{1.782828in}{1.069946in}}%
\pgfpathlineto{\pgfqpoint{1.809764in}{1.070627in}}%
\pgfpathlineto{\pgfqpoint{1.836700in}{1.075857in}}%
\pgfpathlineto{\pgfqpoint{1.863636in}{1.085992in}}%
\pgfpathlineto{\pgfqpoint{1.890572in}{1.101388in}}%
\pgfpathlineto{\pgfqpoint{1.917508in}{1.122400in}}%
\pgfpathlineto{\pgfqpoint{1.944444in}{1.149386in}}%
\pgfpathlineto{\pgfqpoint{1.971380in}{1.182701in}}%
\pgfpathlineto{\pgfqpoint{1.998316in}{1.222701in}}%
\pgfpathlineto{\pgfqpoint{2.025253in}{1.269744in}}%
\pgfpathlineto{\pgfqpoint{2.052189in}{1.324184in}}%
\pgfpathlineto{\pgfqpoint{2.079125in}{1.386378in}}%
\pgfpathlineto{\pgfqpoint{2.106061in}{1.456683in}}%
\pgfpathlineto{\pgfqpoint{2.132997in}{1.535454in}}%
\pgfpathlineto{\pgfqpoint{2.159933in}{1.623047in}}%
\pgfpathlineto{\pgfqpoint{2.186869in}{1.719819in}}%
\pgfpathlineto{\pgfqpoint{2.213805in}{1.826127in}}%
\pgfpathlineto{\pgfqpoint{2.240741in}{1.942325in}}%
\pgfpathlineto{\pgfqpoint{2.267677in}{2.068771in}}%
\pgfpathlineto{\pgfqpoint{2.287900in}{2.171667in}}%
\pgfusepath{stroke}%
\end{pgfscope}%
\begin{pgfscope}%
\pgfpathrectangle{\pgfqpoint{0.033333in}{0.033333in}}{\pgfqpoint{2.933333in}{2.133333in}}%
\pgfusepath{clip}%
\pgfsetrectcap%
\pgfsetroundjoin%
\pgfsetlinewidth{1.505625pt}%
\definecolor{currentstroke}{rgb}{0.172549,0.627451,0.172549}%
\pgfsetstrokecolor{currentstroke}%
\pgfsetdash{}{0pt}%
\pgfpathmoveto{\pgfqpoint{0.556823in}{0.028333in}}%
\pgfpathlineto{\pgfqpoint{0.570707in}{0.119298in}}%
\pgfpathlineto{\pgfqpoint{0.597643in}{0.283944in}}%
\pgfpathlineto{\pgfqpoint{0.624579in}{0.437110in}}%
\pgfpathlineto{\pgfqpoint{0.651515in}{0.579153in}}%
\pgfpathlineto{\pgfqpoint{0.678451in}{0.710428in}}%
\pgfpathlineto{\pgfqpoint{0.705387in}{0.831293in}}%
\pgfpathlineto{\pgfqpoint{0.732323in}{0.942102in}}%
\pgfpathlineto{\pgfqpoint{0.759259in}{1.043213in}}%
\pgfpathlineto{\pgfqpoint{0.786195in}{1.134982in}}%
\pgfpathlineto{\pgfqpoint{0.813131in}{1.217764in}}%
\pgfpathlineto{\pgfqpoint{0.840067in}{1.291916in}}%
\pgfpathlineto{\pgfqpoint{0.867003in}{1.357794in}}%
\pgfpathlineto{\pgfqpoint{0.893939in}{1.415754in}}%
\pgfpathlineto{\pgfqpoint{0.920875in}{1.466153in}}%
\pgfpathlineto{\pgfqpoint{0.947811in}{1.509346in}}%
\pgfpathlineto{\pgfqpoint{0.974747in}{1.545690in}}%
\pgfpathlineto{\pgfqpoint{1.001684in}{1.575540in}}%
\pgfpathlineto{\pgfqpoint{1.028620in}{1.599254in}}%
\pgfpathlineto{\pgfqpoint{1.055556in}{1.617187in}}%
\pgfpathlineto{\pgfqpoint{1.082492in}{1.629696in}}%
\pgfpathlineto{\pgfqpoint{1.109428in}{1.637136in}}%
\pgfpathlineto{\pgfqpoint{1.136364in}{1.639864in}}%
\pgfpathlineto{\pgfqpoint{1.163300in}{1.638236in}}%
\pgfpathlineto{\pgfqpoint{1.190236in}{1.632609in}}%
\pgfpathlineto{\pgfqpoint{1.217172in}{1.623337in}}%
\pgfpathlineto{\pgfqpoint{1.244108in}{1.610778in}}%
\pgfpathlineto{\pgfqpoint{1.271044in}{1.595288in}}%
\pgfpathlineto{\pgfqpoint{1.297980in}{1.577223in}}%
\pgfpathlineto{\pgfqpoint{1.324916in}{1.556939in}}%
\pgfpathlineto{\pgfqpoint{1.351852in}{1.534792in}}%
\pgfpathlineto{\pgfqpoint{1.378788in}{1.511138in}}%
\pgfpathlineto{\pgfqpoint{1.405724in}{1.486334in}}%
\pgfpathlineto{\pgfqpoint{1.432660in}{1.460736in}}%
\pgfpathlineto{\pgfqpoint{1.459596in}{1.434700in}}%
\pgfpathlineto{\pgfqpoint{1.486532in}{1.408582in}}%
\pgfpathlineto{\pgfqpoint{1.513468in}{1.382738in}}%
\pgfpathlineto{\pgfqpoint{1.540404in}{1.357525in}}%
\pgfpathlineto{\pgfqpoint{1.567340in}{1.333298in}}%
\pgfpathlineto{\pgfqpoint{1.594276in}{1.310414in}}%
\pgfpathlineto{\pgfqpoint{1.621212in}{1.289230in}}%
\pgfpathlineto{\pgfqpoint{1.648148in}{1.270100in}}%
\pgfpathlineto{\pgfqpoint{1.675084in}{1.253382in}}%
\pgfpathlineto{\pgfqpoint{1.702020in}{1.239432in}}%
\pgfpathlineto{\pgfqpoint{1.728956in}{1.228605in}}%
\pgfpathlineto{\pgfqpoint{1.755892in}{1.221258in}}%
\pgfpathlineto{\pgfqpoint{1.782828in}{1.217747in}}%
\pgfpathlineto{\pgfqpoint{1.809764in}{1.218429in}}%
\pgfpathlineto{\pgfqpoint{1.836700in}{1.223659in}}%
\pgfpathlineto{\pgfqpoint{1.863636in}{1.233794in}}%
\pgfpathlineto{\pgfqpoint{1.890572in}{1.249189in}}%
\pgfpathlineto{\pgfqpoint{1.917508in}{1.270202in}}%
\pgfpathlineto{\pgfqpoint{1.944444in}{1.297187in}}%
\pgfpathlineto{\pgfqpoint{1.971380in}{1.330502in}}%
\pgfpathlineto{\pgfqpoint{1.998316in}{1.370503in}}%
\pgfpathlineto{\pgfqpoint{2.025253in}{1.417545in}}%
\pgfpathlineto{\pgfqpoint{2.052189in}{1.471985in}}%
\pgfpathlineto{\pgfqpoint{2.079125in}{1.534180in}}%
\pgfpathlineto{\pgfqpoint{2.106061in}{1.604484in}}%
\pgfpathlineto{\pgfqpoint{2.132997in}{1.683255in}}%
\pgfpathlineto{\pgfqpoint{2.159933in}{1.770849in}}%
\pgfpathlineto{\pgfqpoint{2.186869in}{1.867621in}}%
\pgfpathlineto{\pgfqpoint{2.213805in}{1.973928in}}%
\pgfpathlineto{\pgfqpoint{2.240741in}{2.090127in}}%
\pgfpathlineto{\pgfqpoint{2.258111in}{2.171667in}}%
\pgfusepath{stroke}%
\end{pgfscope}%
\begin{pgfscope}%
\pgfsetrectcap%
\pgfsetmiterjoin%
\pgfsetlinewidth{0.803000pt}%
\definecolor{currentstroke}{rgb}{0.000000,0.000000,0.000000}%
\pgfsetstrokecolor{currentstroke}%
\pgfsetdash{}{0pt}%
\pgfpathmoveto{\pgfqpoint{1.500000in}{0.033333in}}%
\pgfpathlineto{\pgfqpoint{1.500000in}{2.166667in}}%
\pgfusepath{stroke}%
\end{pgfscope}%
\begin{pgfscope}%
\pgfsetrectcap%
\pgfsetmiterjoin%
\pgfsetlinewidth{0.803000pt}%
\definecolor{currentstroke}{rgb}{0.000000,0.000000,0.000000}%
\pgfsetstrokecolor{currentstroke}%
\pgfsetdash{}{0pt}%
\pgfpathmoveto{\pgfqpoint{0.033333in}{1.100000in}}%
\pgfpathlineto{\pgfqpoint{2.966667in}{1.100000in}}%
\pgfusepath{stroke}%
\end{pgfscope}%
\begin{pgfscope}%
\definecolor{textcolor}{rgb}{0.000000,0.000000,0.000000}%
\pgfsetstrokecolor{textcolor}%
\pgfsetfillcolor{textcolor}%
\pgftext[x=2.966667in,y=1.002778in,right,top]{\color{textcolor}\rmfamily\fontsize{12.000000}{14.400000}\selectfont \(\displaystyle \lambda\)}%
\end{pgfscope}%
\end{pgfpicture}%
\makeatother%
\endgroup%

          \end{center}
          for three different values of $\alpha$. (You do not need to find the points where it crosses the $x$-axis at this stage.)

          Differentiating and setting to zero, we find
          \begin{equation}
            3\lambda^2+2\lambda\left(c-\frac{\beta}{c}\right)-\beta=0,
          \end{equation}
          and solving the quadratic we find
          \begin{equation}
            \lambda_{\text{max}},\lambda_{\text{min}} = -\frac{1}{3}\left[c - \frac{\beta}{c} \pm \sqrt{\left( c- \frac{\beta}{c} \right)^2 + 3\beta} \right].
          \end{equation}
          %
          \commentary{This was also done fine. Big up to those who drew snowfolk.}
          %
        \item I am convinced. Looking at \cref{main}, changing $\alpha$ is like doing `$f(x)+a$', translating the function vertically. (See picture above.)
          %
          \commentary{Many of you seemed convinced. Good for you: back yourselves!}
          %
        \item Equation (\xmas{candle}) with $\alpha=0$ has solutions $\lambda = 0$ and
          \begin{align}
            \lambda^2 + \lambda\left(c-\frac{\beta}{c}\right)-\beta &=0 \\
            \implies (\lambda+c)\left(\lambda-\frac{\beta}{c}\right) &= 0 \implies \lambda = -c, \frac{\beta}{c}.
          \end{align} The eigenvalues are therefore $\lambda = 0$, $-c$, $\beta/c$.
          %
          \commentary{The majority of you used the quadratic formula, which does simplify to these solutions. But you should always have a quick look as to whether you can factorise clearly: in this case all you need is $(a-b)^2 + 4ab = (a+b)^2$. The solutions from the quadratic equation do not simplify obviously and the clarity of the result is important for interpretation.}
          %
        \item Previously, we discarded these solutions on the grounds that they would yield negative values of a population size. In those instances, however, we were dealing with equilibria where one population size was zero. Here, we have $(\beta,1-\beta)$, so the solution may safely oscillate without ever becoming negative. Thus, we can permit oscillatory wavefronts. While this is a fundamental difference with the FK equation where oscillations were not possible, like the FK equation the wave speed is given by the minimum possible value, namely, $c^2 = 4\alpha(1-\beta)$.
          %
          \commentary{Here the issues were similar to the previous question about validity of travelling waves. The issue in the previous case was forbidding negative $v$; but that is not an issue here.}
      \end{enumerate}

    }
\end{enumerate}

\ifbool{show_solution}{}{
  \begin{figure}
    \centering
    %% Creator: Matplotlib, PGF backend
%%
%% To include the figure in your LaTeX document, write
%%   \input{<filename>.pgf}
%%
%% Make sure the required packages are loaded in your preamble
%%   \usepackage{pgf}
%%
%% and, on pdftex
%%   \usepackage[utf8]{inputenc}\DeclareUnicodeCharacter{2212}{-}
%%
%% or, on luatex and xetex
%%   \usepackage{unicode-math}
%%
%% Figures using additional raster images can only be included by \input if
%% they are in the same directory as the main LaTeX file. For loading figures
%% from other directories you can use the `import` package
%%   \usepackage{import}
%%
%% and then include the figures with
%%   \import{<path to file>}{<filename>.pgf}
%%
%% Matplotlib used the following preamble
%%   \usepackage{fontspec}
%%   \setmainfont{DejaVuSerif.ttf}[Path=/Users/adam/opt/anaconda3/lib/python3.8/site-packages/matplotlib/mpl-data/fonts/ttf/]
%%   \setsansfont{DejaVuSans.ttf}[Path=/Users/adam/opt/anaconda3/lib/python3.8/site-packages/matplotlib/mpl-data/fonts/ttf/]
%%   \setmonofont{DejaVuSansMono.ttf}[Path=/Users/adam/opt/anaconda3/lib/python3.8/site-packages/matplotlib/mpl-data/fonts/ttf/]
%%
\begingroup%
\makeatletter%
\begin{pgfpicture}%
\pgfpathrectangle{\pgfpointorigin}{\pgfqpoint{2.200000in}{2.200000in}}%
\pgfusepath{use as bounding box, clip}%
\begin{pgfscope}%
\pgfsetbuttcap%
\pgfsetmiterjoin%
\definecolor{currentfill}{rgb}{1.000000,1.000000,1.000000}%
\pgfsetfillcolor{currentfill}%
\pgfsetlinewidth{0.000000pt}%
\definecolor{currentstroke}{rgb}{1.000000,1.000000,1.000000}%
\pgfsetstrokecolor{currentstroke}%
\pgfsetdash{}{0pt}%
\pgfpathmoveto{\pgfqpoint{0.000000in}{0.000000in}}%
\pgfpathlineto{\pgfqpoint{2.200000in}{0.000000in}}%
\pgfpathlineto{\pgfqpoint{2.200000in}{2.200000in}}%
\pgfpathlineto{\pgfqpoint{0.000000in}{2.200000in}}%
\pgfpathclose%
\pgfusepath{fill}%
\end{pgfscope}%
\begin{pgfscope}%
\pgfsetbuttcap%
\pgfsetmiterjoin%
\definecolor{currentfill}{rgb}{1.000000,1.000000,1.000000}%
\pgfsetfillcolor{currentfill}%
\pgfsetlinewidth{0.000000pt}%
\definecolor{currentstroke}{rgb}{0.000000,0.000000,0.000000}%
\pgfsetstrokecolor{currentstroke}%
\pgfsetstrokeopacity{0.000000}%
\pgfsetdash{}{0pt}%
\pgfpathmoveto{\pgfqpoint{0.033333in}{0.278150in}}%
\pgfpathlineto{\pgfqpoint{2.166667in}{0.278150in}}%
\pgfpathlineto{\pgfqpoint{2.166667in}{2.166667in}}%
\pgfpathlineto{\pgfqpoint{0.033333in}{2.166667in}}%
\pgfpathclose%
\pgfusepath{fill}%
\end{pgfscope}%
\begin{pgfscope}%
\pgfpathrectangle{\pgfqpoint{0.033333in}{0.278150in}}{\pgfqpoint{2.133333in}{1.888517in}}%
\pgfusepath{clip}%
\pgfsetrectcap%
\pgfsetroundjoin%
\pgfsetlinewidth{1.505625pt}%
\definecolor{currentstroke}{rgb}{0.121569,0.466667,0.705882}%
\pgfsetstrokecolor{currentstroke}%
\pgfsetdash{}{0pt}%
\pgfpathmoveto{\pgfqpoint{0.028333in}{0.879001in}}%
\pgfpathlineto{\pgfqpoint{0.384449in}{0.879198in}}%
\pgfpathlineto{\pgfqpoint{0.464449in}{0.880273in}}%
\pgfpathlineto{\pgfqpoint{0.517783in}{0.881927in}}%
\pgfpathlineto{\pgfqpoint{0.571116in}{0.884930in}}%
\pgfpathlineto{\pgfqpoint{0.597783in}{0.887167in}}%
\pgfpathlineto{\pgfqpoint{0.624449in}{0.890053in}}%
\pgfpathlineto{\pgfqpoint{0.651116in}{0.893728in}}%
\pgfpathlineto{\pgfqpoint{0.677783in}{0.898354in}}%
\pgfpathlineto{\pgfqpoint{0.704449in}{0.904113in}}%
\pgfpathlineto{\pgfqpoint{0.731116in}{0.911208in}}%
\pgfpathlineto{\pgfqpoint{0.757783in}{0.919864in}}%
\pgfpathlineto{\pgfqpoint{0.784449in}{0.930323in}}%
\pgfpathlineto{\pgfqpoint{0.811116in}{0.942847in}}%
\pgfpathlineto{\pgfqpoint{0.837783in}{0.957712in}}%
\pgfpathlineto{\pgfqpoint{0.864449in}{0.975208in}}%
\pgfpathlineto{\pgfqpoint{0.891116in}{0.995630in}}%
\pgfpathlineto{\pgfqpoint{0.917783in}{1.019270in}}%
\pgfpathlineto{\pgfqpoint{0.944449in}{1.046411in}}%
\pgfpathlineto{\pgfqpoint{0.971116in}{1.077306in}}%
\pgfpathlineto{\pgfqpoint{0.997783in}{1.112169in}}%
\pgfpathlineto{\pgfqpoint{1.024449in}{1.151143in}}%
\pgfpathlineto{\pgfqpoint{1.051116in}{1.194284in}}%
\pgfpathlineto{\pgfqpoint{1.077783in}{1.241526in}}%
\pgfpathlineto{\pgfqpoint{1.104449in}{1.292658in}}%
\pgfpathlineto{\pgfqpoint{1.131116in}{1.347299in}}%
\pgfpathlineto{\pgfqpoint{1.157783in}{1.404882in}}%
\pgfpathlineto{\pgfqpoint{1.211116in}{1.525684in}}%
\pgfpathlineto{\pgfqpoint{1.264449in}{1.647211in}}%
\pgfpathlineto{\pgfqpoint{1.291116in}{1.705423in}}%
\pgfpathlineto{\pgfqpoint{1.317783in}{1.760487in}}%
\pgfpathlineto{\pgfqpoint{1.344449in}{1.811490in}}%
\pgfpathlineto{\pgfqpoint{1.371116in}{1.857735in}}%
\pgfpathlineto{\pgfqpoint{1.397783in}{1.898782in}}%
\pgfpathlineto{\pgfqpoint{1.424449in}{1.934456in}}%
\pgfpathlineto{\pgfqpoint{1.451116in}{1.964828in}}%
\pgfpathlineto{\pgfqpoint{1.477783in}{1.990177in}}%
\pgfpathlineto{\pgfqpoint{1.504449in}{2.010935in}}%
\pgfpathlineto{\pgfqpoint{1.531116in}{2.027629in}}%
\pgfpathlineto{\pgfqpoint{1.557783in}{2.040828in}}%
\pgfpathlineto{\pgfqpoint{1.584449in}{2.051099in}}%
\pgfpathlineto{\pgfqpoint{1.611116in}{2.058973in}}%
\pgfpathlineto{\pgfqpoint{1.637783in}{2.064927in}}%
\pgfpathlineto{\pgfqpoint{1.664449in}{2.069371in}}%
\pgfpathlineto{\pgfqpoint{1.691116in}{2.072649in}}%
\pgfpathlineto{\pgfqpoint{1.717783in}{2.075039in}}%
\pgfpathlineto{\pgfqpoint{1.771116in}{2.077998in}}%
\pgfpathlineto{\pgfqpoint{1.824449in}{2.079485in}}%
\pgfpathlineto{\pgfqpoint{1.904449in}{2.080410in}}%
\pgfpathlineto{\pgfqpoint{2.064449in}{2.080791in}}%
\pgfpathlineto{\pgfqpoint{2.171667in}{2.080819in}}%
\pgfpathlineto{\pgfqpoint{2.171667in}{2.080819in}}%
\pgfusepath{stroke}%
\end{pgfscope}%
\begin{pgfscope}%
\pgfpathrectangle{\pgfqpoint{0.033333in}{0.278150in}}{\pgfqpoint{2.133333in}{1.888517in}}%
\pgfusepath{clip}%
\pgfsetbuttcap%
\pgfsetroundjoin%
\pgfsetlinewidth{1.505625pt}%
\definecolor{currentstroke}{rgb}{1.000000,0.498039,0.054902}%
\pgfsetstrokecolor{currentstroke}%
\pgfsetdash{{5.550000pt}{2.400000pt}}{0.000000pt}%
\pgfpathmoveto{\pgfqpoint{0.028333in}{1.565793in}}%
\pgfpathlineto{\pgfqpoint{0.517783in}{1.565456in}}%
\pgfpathlineto{\pgfqpoint{0.571116in}{1.564373in}}%
\pgfpathlineto{\pgfqpoint{0.624449in}{1.562156in}}%
\pgfpathlineto{\pgfqpoint{0.651116in}{1.560386in}}%
\pgfpathlineto{\pgfqpoint{0.677783in}{1.558005in}}%
\pgfpathlineto{\pgfqpoint{0.704449in}{1.554854in}}%
\pgfpathlineto{\pgfqpoint{0.731116in}{1.550741in}}%
\pgfpathlineto{\pgfqpoint{0.757783in}{1.545444in}}%
\pgfpathlineto{\pgfqpoint{0.784449in}{1.538706in}}%
\pgfpathlineto{\pgfqpoint{0.811116in}{1.530234in}}%
\pgfpathlineto{\pgfqpoint{0.837783in}{1.519704in}}%
\pgfpathlineto{\pgfqpoint{0.864449in}{1.506757in}}%
\pgfpathlineto{\pgfqpoint{0.891116in}{1.491012in}}%
\pgfpathlineto{\pgfqpoint{0.917783in}{1.472070in}}%
\pgfpathlineto{\pgfqpoint{0.944449in}{1.449529in}}%
\pgfpathlineto{\pgfqpoint{0.971116in}{1.423000in}}%
\pgfpathlineto{\pgfqpoint{0.997783in}{1.392133in}}%
\pgfpathlineto{\pgfqpoint{1.024449in}{1.356639in}}%
\pgfpathlineto{\pgfqpoint{1.051116in}{1.316330in}}%
\pgfpathlineto{\pgfqpoint{1.077783in}{1.271148in}}%
\pgfpathlineto{\pgfqpoint{1.104449in}{1.221202in}}%
\pgfpathlineto{\pgfqpoint{1.131116in}{1.166800in}}%
\pgfpathlineto{\pgfqpoint{1.157783in}{1.108471in}}%
\pgfpathlineto{\pgfqpoint{1.184449in}{1.046974in}}%
\pgfpathlineto{\pgfqpoint{1.237783in}{0.918536in}}%
\pgfpathlineto{\pgfqpoint{1.264449in}{0.854014in}}%
\pgfpathlineto{\pgfqpoint{1.291116in}{0.791016in}}%
\pgfpathlineto{\pgfqpoint{1.317783in}{0.730793in}}%
\pgfpathlineto{\pgfqpoint{1.344449in}{0.674447in}}%
\pgfpathlineto{\pgfqpoint{1.371116in}{0.622863in}}%
\pgfpathlineto{\pgfqpoint{1.397783in}{0.576648in}}%
\pgfpathlineto{\pgfqpoint{1.424449in}{0.536118in}}%
\pgfpathlineto{\pgfqpoint{1.451116in}{0.501308in}}%
\pgfpathlineto{\pgfqpoint{1.477783in}{0.472005in}}%
\pgfpathlineto{\pgfqpoint{1.504449in}{0.447809in}}%
\pgfpathlineto{\pgfqpoint{1.531116in}{0.428193in}}%
\pgfpathlineto{\pgfqpoint{1.557783in}{0.412561in}}%
\pgfpathlineto{\pgfqpoint{1.584449in}{0.400306in}}%
\pgfpathlineto{\pgfqpoint{1.611116in}{0.390842in}}%
\pgfpathlineto{\pgfqpoint{1.637783in}{0.383637in}}%
\pgfpathlineto{\pgfqpoint{1.664449in}{0.378222in}}%
\pgfpathlineto{\pgfqpoint{1.691116in}{0.374203in}}%
\pgfpathlineto{\pgfqpoint{1.717783in}{0.371254in}}%
\pgfpathlineto{\pgfqpoint{1.744449in}{0.369113in}}%
\pgfpathlineto{\pgfqpoint{1.797783in}{0.366478in}}%
\pgfpathlineto{\pgfqpoint{1.851116in}{0.365164in}}%
\pgfpathlineto{\pgfqpoint{1.931116in}{0.364352in}}%
\pgfpathlineto{\pgfqpoint{2.117783in}{0.364010in}}%
\pgfpathlineto{\pgfqpoint{2.171667in}{0.363999in}}%
\pgfpathlineto{\pgfqpoint{2.171667in}{0.363999in}}%
\pgfusepath{stroke}%
\end{pgfscope}%
\begin{pgfscope}%
\pgfsetbuttcap%
\pgfsetroundjoin%
\definecolor{currentfill}{rgb}{0.000000,0.000000,0.000000}%
\pgfsetfillcolor{currentfill}%
\pgfsetlinewidth{0.803000pt}%
\definecolor{currentstroke}{rgb}{0.000000,0.000000,0.000000}%
\pgfsetstrokecolor{currentstroke}%
\pgfsetdash{}{0pt}%
\pgfsys@defobject{currentmarker}{\pgfqpoint{0.000000in}{-0.048611in}}{\pgfqpoint{0.000000in}{0.000000in}}{%
\pgfpathmoveto{\pgfqpoint{0.000000in}{0.000000in}}%
\pgfpathlineto{\pgfqpoint{0.000000in}{-0.048611in}}%
\pgfusepath{stroke,fill}%
}%
\begin{pgfscope}%
\pgfsys@transformshift{1.366667in}{0.363991in}%
\pgfsys@useobject{currentmarker}{}%
\end{pgfscope}%
\end{pgfscope}%
\begin{pgfscope}%
\definecolor{textcolor}{rgb}{0.000000,0.000000,0.000000}%
\pgfsetstrokecolor{textcolor}%
\pgfsetfillcolor{textcolor}%
\pgftext[x=1.366667in,y=0.266769in,,top]{\color{textcolor}\rmfamily\fontsize{12.000000}{14.400000}\selectfont 0}%
\end{pgfscope}%
\begin{pgfscope}%
\pgfsetbuttcap%
\pgfsetroundjoin%
\definecolor{currentfill}{rgb}{0.000000,0.000000,0.000000}%
\pgfsetfillcolor{currentfill}%
\pgfsetlinewidth{0.803000pt}%
\definecolor{currentstroke}{rgb}{0.000000,0.000000,0.000000}%
\pgfsetstrokecolor{currentstroke}%
\pgfsetdash{}{0pt}%
\pgfsys@defobject{currentmarker}{\pgfqpoint{-0.048611in}{0.000000in}}{\pgfqpoint{-0.000000in}{0.000000in}}{%
\pgfpathmoveto{\pgfqpoint{-0.000000in}{0.000000in}}%
\pgfpathlineto{\pgfqpoint{-0.048611in}{0.000000in}}%
\pgfusepath{stroke,fill}%
}%
\begin{pgfscope}%
\pgfsys@transformshift{1.366667in}{0.363991in}%
\pgfsys@useobject{currentmarker}{}%
\end{pgfscope}%
\end{pgfscope}%
\begin{pgfscope}%
\definecolor{textcolor}{rgb}{0.000000,0.000000,0.000000}%
\pgfsetstrokecolor{textcolor}%
\pgfsetfillcolor{textcolor}%
\pgftext[x=1.163406in, y=0.300678in, left, base]{\color{textcolor}\rmfamily\fontsize{12.000000}{14.400000}\selectfont 0}%
\end{pgfscope}%
\begin{pgfscope}%
\pgfsetbuttcap%
\pgfsetroundjoin%
\definecolor{currentfill}{rgb}{0.000000,0.000000,0.000000}%
\pgfsetfillcolor{currentfill}%
\pgfsetlinewidth{0.803000pt}%
\definecolor{currentstroke}{rgb}{0.000000,0.000000,0.000000}%
\pgfsetstrokecolor{currentstroke}%
\pgfsetdash{}{0pt}%
\pgfsys@defobject{currentmarker}{\pgfqpoint{-0.048611in}{0.000000in}}{\pgfqpoint{-0.000000in}{0.000000in}}{%
\pgfpathmoveto{\pgfqpoint{-0.000000in}{0.000000in}}%
\pgfpathlineto{\pgfqpoint{-0.048611in}{0.000000in}}%
\pgfusepath{stroke,fill}%
}%
\begin{pgfscope}%
\pgfsys@transformshift{1.366667in}{0.879041in}%
\pgfsys@useobject{currentmarker}{}%
\end{pgfscope}%
\end{pgfscope}%
\begin{pgfscope}%
\definecolor{textcolor}{rgb}{0.000000,0.000000,0.000000}%
\pgfsetstrokecolor{textcolor}%
\pgfsetfillcolor{textcolor}%
\pgftext[x=1.168076in, y=0.815728in, left, base]{\color{textcolor}\rmfamily\fontsize{12.000000}{14.400000}\selectfont \(\displaystyle \beta\)}%
\end{pgfscope}%
\begin{pgfscope}%
\pgfsetbuttcap%
\pgfsetroundjoin%
\definecolor{currentfill}{rgb}{0.000000,0.000000,0.000000}%
\pgfsetfillcolor{currentfill}%
\pgfsetlinewidth{0.803000pt}%
\definecolor{currentstroke}{rgb}{0.000000,0.000000,0.000000}%
\pgfsetstrokecolor{currentstroke}%
\pgfsetdash{}{0pt}%
\pgfsys@defobject{currentmarker}{\pgfqpoint{-0.048611in}{0.000000in}}{\pgfqpoint{-0.000000in}{0.000000in}}{%
\pgfpathmoveto{\pgfqpoint{-0.000000in}{0.000000in}}%
\pgfpathlineto{\pgfqpoint{-0.048611in}{0.000000in}}%
\pgfusepath{stroke,fill}%
}%
\begin{pgfscope}%
\pgfsys@transformshift{1.366667in}{1.565775in}%
\pgfsys@useobject{currentmarker}{}%
\end{pgfscope}%
\end{pgfscope}%
\begin{pgfscope}%
\definecolor{textcolor}{rgb}{0.000000,0.000000,0.000000}%
\pgfsetstrokecolor{textcolor}%
\pgfsetfillcolor{textcolor}%
\pgftext[x=0.882777in, y=1.502461in, left, base]{\color{textcolor}\rmfamily\fontsize{12.000000}{14.400000}\selectfont \(\displaystyle 1-\beta\)}%
\end{pgfscope}%
\begin{pgfscope}%
\pgfsetbuttcap%
\pgfsetroundjoin%
\definecolor{currentfill}{rgb}{0.000000,0.000000,0.000000}%
\pgfsetfillcolor{currentfill}%
\pgfsetlinewidth{0.803000pt}%
\definecolor{currentstroke}{rgb}{0.000000,0.000000,0.000000}%
\pgfsetstrokecolor{currentstroke}%
\pgfsetdash{}{0pt}%
\pgfsys@defobject{currentmarker}{\pgfqpoint{-0.048611in}{0.000000in}}{\pgfqpoint{-0.000000in}{0.000000in}}{%
\pgfpathmoveto{\pgfqpoint{-0.000000in}{0.000000in}}%
\pgfpathlineto{\pgfqpoint{-0.048611in}{0.000000in}}%
\pgfusepath{stroke,fill}%
}%
\begin{pgfscope}%
\pgfsys@transformshift{1.366667in}{2.080825in}%
\pgfsys@useobject{currentmarker}{}%
\end{pgfscope}%
\end{pgfscope}%
\begin{pgfscope}%
\definecolor{textcolor}{rgb}{0.000000,0.000000,0.000000}%
\pgfsetstrokecolor{textcolor}%
\pgfsetfillcolor{textcolor}%
\pgftext[x=1.163406in, y=2.017511in, left, base]{\color{textcolor}\rmfamily\fontsize{12.000000}{14.400000}\selectfont 1}%
\end{pgfscope}%
\begin{pgfscope}%
\pgfsetrectcap%
\pgfsetmiterjoin%
\pgfsetlinewidth{0.803000pt}%
\definecolor{currentstroke}{rgb}{0.000000,0.000000,0.000000}%
\pgfsetstrokecolor{currentstroke}%
\pgfsetdash{}{0pt}%
\pgfpathmoveto{\pgfqpoint{1.366667in}{0.278150in}}%
\pgfpathlineto{\pgfqpoint{1.366667in}{2.166667in}}%
\pgfusepath{stroke}%
\end{pgfscope}%
\begin{pgfscope}%
\pgfsetrectcap%
\pgfsetmiterjoin%
\pgfsetlinewidth{0.803000pt}%
\definecolor{currentstroke}{rgb}{0.000000,0.000000,0.000000}%
\pgfsetstrokecolor{currentstroke}%
\pgfsetdash{}{0pt}%
\pgfpathmoveto{\pgfqpoint{0.033333in}{0.363991in}}%
\pgfpathlineto{\pgfqpoint{2.166667in}{0.363991in}}%
\pgfusepath{stroke}%
\end{pgfscope}%
\begin{pgfscope}%
\definecolor{textcolor}{rgb}{0.000000,0.000000,0.000000}%
\pgfsetstrokecolor{textcolor}%
\pgfsetfillcolor{textcolor}%
\pgftext[x=2.166667in,y=0.180927in,right,top]{\color{textcolor}\rmfamily\fontsize{12.000000}{14.400000}\selectfont \(\displaystyle z\)}%
\end{pgfscope}%
\begin{pgfscope}%
\pgfsetrectcap%
\pgfsetroundjoin%
\pgfsetlinewidth{1.505625pt}%
\definecolor{currentstroke}{rgb}{0.121569,0.466667,0.705882}%
\pgfsetstrokecolor{currentstroke}%
\pgfsetdash{}{0pt}%
\pgfpathmoveto{\pgfqpoint{0.033333in}{2.098372in}}%
\pgfpathlineto{\pgfqpoint{0.366667in}{2.098372in}}%
\pgfusepath{stroke}%
\end{pgfscope}%
\begin{pgfscope}%
\definecolor{textcolor}{rgb}{0.000000,0.000000,0.000000}%
\pgfsetstrokecolor{textcolor}%
\pgfsetfillcolor{textcolor}%
\pgftext[x=0.500000in,y=2.040039in,left,base]{\color{textcolor}\rmfamily\fontsize{12.000000}{14.400000}\selectfont \(\displaystyle u\)}%
\end{pgfscope}%
\begin{pgfscope}%
\pgfsetbuttcap%
\pgfsetroundjoin%
\pgfsetlinewidth{1.505625pt}%
\definecolor{currentstroke}{rgb}{1.000000,0.498039,0.054902}%
\pgfsetstrokecolor{currentstroke}%
\pgfsetdash{{5.550000pt}{2.400000pt}}{0.000000pt}%
\pgfpathmoveto{\pgfqpoint{0.033333in}{1.853744in}}%
\pgfpathlineto{\pgfqpoint{0.366667in}{1.853744in}}%
\pgfusepath{stroke}%
\end{pgfscope}%
\begin{pgfscope}%
\definecolor{textcolor}{rgb}{0.000000,0.000000,0.000000}%
\pgfsetstrokecolor{textcolor}%
\pgfsetfillcolor{textcolor}%
\pgftext[x=0.500000in,y=1.795410in,left,base]{\color{textcolor}\rmfamily\fontsize{12.000000}{14.400000}\selectfont \(\displaystyle v\)}%
\end{pgfscope}%
\end{pgfpicture}%
\makeatother%
\endgroup%
\qquad
    %% Creator: Matplotlib, PGF backend
%%
%% To include the figure in your LaTeX document, write
%%   \input{<filename>.pgf}
%%
%% Make sure the required packages are loaded in your preamble
%%   \usepackage{pgf}
%%
%% and, on pdftex
%%   \usepackage[utf8]{inputenc}\DeclareUnicodeCharacter{2212}{-}
%%
%% or, on luatex and xetex
%%   \usepackage{unicode-math}
%%
%% Figures using additional raster images can only be included by \input if
%% they are in the same directory as the main LaTeX file. For loading figures
%% from other directories you can use the `import` package
%%   \usepackage{import}
%%
%% and then include the figures with
%%   \import{<path to file>}{<filename>.pgf}
%%
%% Matplotlib used the following preamble
%%   \usepackage{fontspec}
%%   \setmainfont{DejaVuSerif.ttf}[Path=/Users/adam/opt/anaconda3/lib/python3.8/site-packages/matplotlib/mpl-data/fonts/ttf/]
%%   \setsansfont{DejaVuSans.ttf}[Path=/Users/adam/opt/anaconda3/lib/python3.8/site-packages/matplotlib/mpl-data/fonts/ttf/]
%%   \setmonofont{DejaVuSansMono.ttf}[Path=/Users/adam/opt/anaconda3/lib/python3.8/site-packages/matplotlib/mpl-data/fonts/ttf/]
%%
\begingroup%
\makeatletter%
\begin{pgfpicture}%
\pgfpathrectangle{\pgfpointorigin}{\pgfqpoint{2.200000in}{2.200000in}}%
\pgfusepath{use as bounding box, clip}%
\begin{pgfscope}%
\pgfsetbuttcap%
\pgfsetmiterjoin%
\definecolor{currentfill}{rgb}{1.000000,1.000000,1.000000}%
\pgfsetfillcolor{currentfill}%
\pgfsetlinewidth{0.000000pt}%
\definecolor{currentstroke}{rgb}{1.000000,1.000000,1.000000}%
\pgfsetstrokecolor{currentstroke}%
\pgfsetdash{}{0pt}%
\pgfpathmoveto{\pgfqpoint{0.000000in}{0.000000in}}%
\pgfpathlineto{\pgfqpoint{2.200000in}{0.000000in}}%
\pgfpathlineto{\pgfqpoint{2.200000in}{2.200000in}}%
\pgfpathlineto{\pgfqpoint{0.000000in}{2.200000in}}%
\pgfpathclose%
\pgfusepath{fill}%
\end{pgfscope}%
\begin{pgfscope}%
\pgfsetbuttcap%
\pgfsetmiterjoin%
\definecolor{currentfill}{rgb}{1.000000,1.000000,1.000000}%
\pgfsetfillcolor{currentfill}%
\pgfsetlinewidth{0.000000pt}%
\definecolor{currentstroke}{rgb}{0.000000,0.000000,0.000000}%
\pgfsetstrokecolor{currentstroke}%
\pgfsetstrokeopacity{0.000000}%
\pgfsetdash{}{0pt}%
\pgfpathmoveto{\pgfqpoint{0.033333in}{0.278150in}}%
\pgfpathlineto{\pgfqpoint{2.166667in}{0.278150in}}%
\pgfpathlineto{\pgfqpoint{2.166667in}{2.166667in}}%
\pgfpathlineto{\pgfqpoint{0.033333in}{2.166667in}}%
\pgfpathclose%
\pgfusepath{fill}%
\end{pgfscope}%
\begin{pgfscope}%
\pgfpathrectangle{\pgfqpoint{0.033333in}{0.278150in}}{\pgfqpoint{2.133333in}{1.888517in}}%
\pgfusepath{clip}%
\pgfsetrectcap%
\pgfsetroundjoin%
\pgfsetlinewidth{1.505625pt}%
\definecolor{currentstroke}{rgb}{0.121569,0.466667,0.705882}%
\pgfsetstrokecolor{currentstroke}%
\pgfsetdash{}{0pt}%
\pgfpathmoveto{\pgfqpoint{0.028333in}{0.880030in}}%
\pgfpathlineto{\pgfqpoint{0.120131in}{0.880145in}}%
\pgfpathlineto{\pgfqpoint{0.200131in}{0.879273in}}%
\pgfpathlineto{\pgfqpoint{0.280131in}{0.877348in}}%
\pgfpathlineto{\pgfqpoint{0.386798in}{0.874454in}}%
\pgfpathlineto{\pgfqpoint{0.440131in}{0.874009in}}%
\pgfpathlineto{\pgfqpoint{0.493465in}{0.875090in}}%
\pgfpathlineto{\pgfqpoint{0.546798in}{0.878199in}}%
\pgfpathlineto{\pgfqpoint{0.600131in}{0.883436in}}%
\pgfpathlineto{\pgfqpoint{0.733465in}{0.899607in}}%
\pgfpathlineto{\pgfqpoint{0.760131in}{0.901234in}}%
\pgfpathlineto{\pgfqpoint{0.786798in}{0.901452in}}%
\pgfpathlineto{\pgfqpoint{0.813465in}{0.899882in}}%
\pgfpathlineto{\pgfqpoint{0.840131in}{0.896211in}}%
\pgfpathlineto{\pgfqpoint{0.866798in}{0.890246in}}%
\pgfpathlineto{\pgfqpoint{0.893465in}{0.881963in}}%
\pgfpathlineto{\pgfqpoint{0.920131in}{0.871539in}}%
\pgfpathlineto{\pgfqpoint{0.946798in}{0.859372in}}%
\pgfpathlineto{\pgfqpoint{1.026798in}{0.819326in}}%
\pgfpathlineto{\pgfqpoint{1.053465in}{0.807800in}}%
\pgfpathlineto{\pgfqpoint{1.080131in}{0.798875in}}%
\pgfpathlineto{\pgfqpoint{1.106798in}{0.793615in}}%
\pgfpathlineto{\pgfqpoint{1.133465in}{0.793121in}}%
\pgfpathlineto{\pgfqpoint{1.160131in}{0.798569in}}%
\pgfpathlineto{\pgfqpoint{1.186798in}{0.811271in}}%
\pgfpathlineto{\pgfqpoint{1.213465in}{0.832743in}}%
\pgfpathlineto{\pgfqpoint{1.240131in}{0.864762in}}%
\pgfpathlineto{\pgfqpoint{1.266798in}{0.909366in}}%
\pgfpathlineto{\pgfqpoint{1.293465in}{0.968736in}}%
\pgfpathlineto{\pgfqpoint{1.320131in}{1.044879in}}%
\pgfpathlineto{\pgfqpoint{1.346798in}{1.138999in}}%
\pgfpathlineto{\pgfqpoint{1.373465in}{1.250567in}}%
\pgfpathlineto{\pgfqpoint{1.400131in}{1.376283in}}%
\pgfpathlineto{\pgfqpoint{1.453465in}{1.640816in}}%
\pgfpathlineto{\pgfqpoint{1.480131in}{1.760251in}}%
\pgfpathlineto{\pgfqpoint{1.506798in}{1.860056in}}%
\pgfpathlineto{\pgfqpoint{1.533465in}{1.936775in}}%
\pgfpathlineto{\pgfqpoint{1.560131in}{1.991340in}}%
\pgfpathlineto{\pgfqpoint{1.586798in}{2.027574in}}%
\pgfpathlineto{\pgfqpoint{1.613465in}{2.050272in}}%
\pgfpathlineto{\pgfqpoint{1.640131in}{2.063822in}}%
\pgfpathlineto{\pgfqpoint{1.666798in}{2.071601in}}%
\pgfpathlineto{\pgfqpoint{1.693465in}{2.075926in}}%
\pgfpathlineto{\pgfqpoint{1.720131in}{2.078269in}}%
\pgfpathlineto{\pgfqpoint{1.746798in}{2.079512in}}%
\pgfpathlineto{\pgfqpoint{1.800131in}{2.080491in}}%
\pgfpathlineto{\pgfqpoint{1.933465in}{2.080816in}}%
\pgfpathlineto{\pgfqpoint{2.171667in}{2.080825in}}%
\pgfpathlineto{\pgfqpoint{2.171667in}{2.080825in}}%
\pgfusepath{stroke}%
\end{pgfscope}%
\begin{pgfscope}%
\pgfpathrectangle{\pgfqpoint{0.033333in}{0.278150in}}{\pgfqpoint{2.133333in}{1.888517in}}%
\pgfusepath{clip}%
\pgfsetbuttcap%
\pgfsetroundjoin%
\pgfsetlinewidth{1.505625pt}%
\definecolor{currentstroke}{rgb}{1.000000,0.498039,0.054902}%
\pgfsetstrokecolor{currentstroke}%
\pgfsetdash{{5.550000pt}{2.400000pt}}{0.000000pt}%
\pgfpathmoveto{\pgfqpoint{0.028333in}{1.565516in}}%
\pgfpathlineto{\pgfqpoint{0.253465in}{1.562998in}}%
\pgfpathlineto{\pgfqpoint{0.306798in}{1.563981in}}%
\pgfpathlineto{\pgfqpoint{0.360131in}{1.566166in}}%
\pgfpathlineto{\pgfqpoint{0.413465in}{1.569466in}}%
\pgfpathlineto{\pgfqpoint{0.520131in}{1.576780in}}%
\pgfpathlineto{\pgfqpoint{0.546798in}{1.577867in}}%
\pgfpathlineto{\pgfqpoint{0.573465in}{1.578255in}}%
\pgfpathlineto{\pgfqpoint{0.600131in}{1.577729in}}%
\pgfpathlineto{\pgfqpoint{0.626798in}{1.576094in}}%
\pgfpathlineto{\pgfqpoint{0.653465in}{1.573194in}}%
\pgfpathlineto{\pgfqpoint{0.680131in}{1.568940in}}%
\pgfpathlineto{\pgfqpoint{0.706798in}{1.563331in}}%
\pgfpathlineto{\pgfqpoint{0.733465in}{1.556480in}}%
\pgfpathlineto{\pgfqpoint{0.760131in}{1.548637in}}%
\pgfpathlineto{\pgfqpoint{0.840131in}{1.523767in}}%
\pgfpathlineto{\pgfqpoint{0.866798in}{1.517220in}}%
\pgfpathlineto{\pgfqpoint{0.893465in}{1.512860in}}%
\pgfpathlineto{\pgfqpoint{0.920131in}{1.511521in}}%
\pgfpathlineto{\pgfqpoint{0.946798in}{1.513982in}}%
\pgfpathlineto{\pgfqpoint{0.973465in}{1.520916in}}%
\pgfpathlineto{\pgfqpoint{1.000131in}{1.532828in}}%
\pgfpathlineto{\pgfqpoint{1.026798in}{1.550007in}}%
\pgfpathlineto{\pgfqpoint{1.053465in}{1.572466in}}%
\pgfpathlineto{\pgfqpoint{1.080131in}{1.599891in}}%
\pgfpathlineto{\pgfqpoint{1.106798in}{1.631572in}}%
\pgfpathlineto{\pgfqpoint{1.133465in}{1.666317in}}%
\pgfpathlineto{\pgfqpoint{1.160131in}{1.702353in}}%
\pgfpathlineto{\pgfqpoint{1.186798in}{1.737193in}}%
\pgfpathlineto{\pgfqpoint{1.213465in}{1.767510in}}%
\pgfpathlineto{\pgfqpoint{1.240131in}{1.789018in}}%
\pgfpathlineto{\pgfqpoint{1.266798in}{1.796442in}}%
\pgfpathlineto{\pgfqpoint{1.293465in}{1.783683in}}%
\pgfpathlineto{\pgfqpoint{1.320131in}{1.744341in}}%
\pgfpathlineto{\pgfqpoint{1.346798in}{1.672786in}}%
\pgfpathlineto{\pgfqpoint{1.373465in}{1.565875in}}%
\pgfpathlineto{\pgfqpoint{1.400131in}{1.425037in}}%
\pgfpathlineto{\pgfqpoint{1.426798in}{1.257924in}}%
\pgfpathlineto{\pgfqpoint{1.453465in}{1.078268in}}%
\pgfpathlineto{\pgfqpoint{1.480131in}{0.903075in}}%
\pgfpathlineto{\pgfqpoint{1.506798in}{0.747864in}}%
\pgfpathlineto{\pgfqpoint{1.533465in}{0.622393in}}%
\pgfpathlineto{\pgfqpoint{1.560131in}{0.529146in}}%
\pgfpathlineto{\pgfqpoint{1.586798in}{0.464806in}}%
\pgfpathlineto{\pgfqpoint{1.613465in}{0.423141in}}%
\pgfpathlineto{\pgfqpoint{1.640131in}{0.397552in}}%
\pgfpathlineto{\pgfqpoint{1.666798in}{0.382505in}}%
\pgfpathlineto{\pgfqpoint{1.693465in}{0.373967in}}%
\pgfpathlineto{\pgfqpoint{1.720131in}{0.369261in}}%
\pgfpathlineto{\pgfqpoint{1.746798in}{0.366728in}}%
\pgfpathlineto{\pgfqpoint{1.773465in}{0.365392in}}%
\pgfpathlineto{\pgfqpoint{1.826798in}{0.364345in}}%
\pgfpathlineto{\pgfqpoint{1.960131in}{0.364001in}}%
\pgfpathlineto{\pgfqpoint{2.171667in}{0.363991in}}%
\pgfpathlineto{\pgfqpoint{2.171667in}{0.363991in}}%
\pgfusepath{stroke}%
\end{pgfscope}%
\begin{pgfscope}%
\pgfsetbuttcap%
\pgfsetroundjoin%
\definecolor{currentfill}{rgb}{0.000000,0.000000,0.000000}%
\pgfsetfillcolor{currentfill}%
\pgfsetlinewidth{0.803000pt}%
\definecolor{currentstroke}{rgb}{0.000000,0.000000,0.000000}%
\pgfsetstrokecolor{currentstroke}%
\pgfsetdash{}{0pt}%
\pgfsys@defobject{currentmarker}{\pgfqpoint{0.000000in}{-0.048611in}}{\pgfqpoint{0.000000in}{0.000000in}}{%
\pgfpathmoveto{\pgfqpoint{0.000000in}{0.000000in}}%
\pgfpathlineto{\pgfqpoint{0.000000in}{-0.048611in}}%
\pgfusepath{stroke,fill}%
}%
\begin{pgfscope}%
\pgfsys@transformshift{1.366667in}{0.363991in}%
\pgfsys@useobject{currentmarker}{}%
\end{pgfscope}%
\end{pgfscope}%
\begin{pgfscope}%
\definecolor{textcolor}{rgb}{0.000000,0.000000,0.000000}%
\pgfsetstrokecolor{textcolor}%
\pgfsetfillcolor{textcolor}%
\pgftext[x=1.366667in,y=0.266769in,,top]{\color{textcolor}\rmfamily\fontsize{12.000000}{14.400000}\selectfont 0}%
\end{pgfscope}%
\begin{pgfscope}%
\pgfsetbuttcap%
\pgfsetroundjoin%
\definecolor{currentfill}{rgb}{0.000000,0.000000,0.000000}%
\pgfsetfillcolor{currentfill}%
\pgfsetlinewidth{0.803000pt}%
\definecolor{currentstroke}{rgb}{0.000000,0.000000,0.000000}%
\pgfsetstrokecolor{currentstroke}%
\pgfsetdash{}{0pt}%
\pgfsys@defobject{currentmarker}{\pgfqpoint{-0.048611in}{0.000000in}}{\pgfqpoint{-0.000000in}{0.000000in}}{%
\pgfpathmoveto{\pgfqpoint{-0.000000in}{0.000000in}}%
\pgfpathlineto{\pgfqpoint{-0.048611in}{0.000000in}}%
\pgfusepath{stroke,fill}%
}%
\begin{pgfscope}%
\pgfsys@transformshift{1.366667in}{0.363991in}%
\pgfsys@useobject{currentmarker}{}%
\end{pgfscope}%
\end{pgfscope}%
\begin{pgfscope}%
\definecolor{textcolor}{rgb}{0.000000,0.000000,0.000000}%
\pgfsetstrokecolor{textcolor}%
\pgfsetfillcolor{textcolor}%
\pgftext[x=1.163406in, y=0.300678in, left, base]{\color{textcolor}\rmfamily\fontsize{12.000000}{14.400000}\selectfont 0}%
\end{pgfscope}%
\begin{pgfscope}%
\pgfsetbuttcap%
\pgfsetroundjoin%
\definecolor{currentfill}{rgb}{0.000000,0.000000,0.000000}%
\pgfsetfillcolor{currentfill}%
\pgfsetlinewidth{0.803000pt}%
\definecolor{currentstroke}{rgb}{0.000000,0.000000,0.000000}%
\pgfsetstrokecolor{currentstroke}%
\pgfsetdash{}{0pt}%
\pgfsys@defobject{currentmarker}{\pgfqpoint{-0.048611in}{0.000000in}}{\pgfqpoint{-0.000000in}{0.000000in}}{%
\pgfpathmoveto{\pgfqpoint{-0.000000in}{0.000000in}}%
\pgfpathlineto{\pgfqpoint{-0.048611in}{0.000000in}}%
\pgfusepath{stroke,fill}%
}%
\begin{pgfscope}%
\pgfsys@transformshift{1.366667in}{0.879041in}%
\pgfsys@useobject{currentmarker}{}%
\end{pgfscope}%
\end{pgfscope}%
\begin{pgfscope}%
\definecolor{textcolor}{rgb}{0.000000,0.000000,0.000000}%
\pgfsetstrokecolor{textcolor}%
\pgfsetfillcolor{textcolor}%
\pgftext[x=1.168076in, y=0.815728in, left, base]{\color{textcolor}\rmfamily\fontsize{12.000000}{14.400000}\selectfont \(\displaystyle \beta\)}%
\end{pgfscope}%
\begin{pgfscope}%
\pgfsetbuttcap%
\pgfsetroundjoin%
\definecolor{currentfill}{rgb}{0.000000,0.000000,0.000000}%
\pgfsetfillcolor{currentfill}%
\pgfsetlinewidth{0.803000pt}%
\definecolor{currentstroke}{rgb}{0.000000,0.000000,0.000000}%
\pgfsetstrokecolor{currentstroke}%
\pgfsetdash{}{0pt}%
\pgfsys@defobject{currentmarker}{\pgfqpoint{-0.048611in}{0.000000in}}{\pgfqpoint{-0.000000in}{0.000000in}}{%
\pgfpathmoveto{\pgfqpoint{-0.000000in}{0.000000in}}%
\pgfpathlineto{\pgfqpoint{-0.048611in}{0.000000in}}%
\pgfusepath{stroke,fill}%
}%
\begin{pgfscope}%
\pgfsys@transformshift{1.366667in}{1.565775in}%
\pgfsys@useobject{currentmarker}{}%
\end{pgfscope}%
\end{pgfscope}%
\begin{pgfscope}%
\definecolor{textcolor}{rgb}{0.000000,0.000000,0.000000}%
\pgfsetstrokecolor{textcolor}%
\pgfsetfillcolor{textcolor}%
\pgftext[x=0.882777in, y=1.502461in, left, base]{\color{textcolor}\rmfamily\fontsize{12.000000}{14.400000}\selectfont \(\displaystyle 1-\beta\)}%
\end{pgfscope}%
\begin{pgfscope}%
\pgfsetbuttcap%
\pgfsetroundjoin%
\definecolor{currentfill}{rgb}{0.000000,0.000000,0.000000}%
\pgfsetfillcolor{currentfill}%
\pgfsetlinewidth{0.803000pt}%
\definecolor{currentstroke}{rgb}{0.000000,0.000000,0.000000}%
\pgfsetstrokecolor{currentstroke}%
\pgfsetdash{}{0pt}%
\pgfsys@defobject{currentmarker}{\pgfqpoint{-0.048611in}{0.000000in}}{\pgfqpoint{-0.000000in}{0.000000in}}{%
\pgfpathmoveto{\pgfqpoint{-0.000000in}{0.000000in}}%
\pgfpathlineto{\pgfqpoint{-0.048611in}{0.000000in}}%
\pgfusepath{stroke,fill}%
}%
\begin{pgfscope}%
\pgfsys@transformshift{1.366667in}{2.080825in}%
\pgfsys@useobject{currentmarker}{}%
\end{pgfscope}%
\end{pgfscope}%
\begin{pgfscope}%
\definecolor{textcolor}{rgb}{0.000000,0.000000,0.000000}%
\pgfsetstrokecolor{textcolor}%
\pgfsetfillcolor{textcolor}%
\pgftext[x=1.163406in, y=2.017511in, left, base]{\color{textcolor}\rmfamily\fontsize{12.000000}{14.400000}\selectfont 1}%
\end{pgfscope}%
\begin{pgfscope}%
\pgfsetrectcap%
\pgfsetmiterjoin%
\pgfsetlinewidth{0.803000pt}%
\definecolor{currentstroke}{rgb}{0.000000,0.000000,0.000000}%
\pgfsetstrokecolor{currentstroke}%
\pgfsetdash{}{0pt}%
\pgfpathmoveto{\pgfqpoint{1.366667in}{0.278150in}}%
\pgfpathlineto{\pgfqpoint{1.366667in}{2.166667in}}%
\pgfusepath{stroke}%
\end{pgfscope}%
\begin{pgfscope}%
\pgfsetrectcap%
\pgfsetmiterjoin%
\pgfsetlinewidth{0.803000pt}%
\definecolor{currentstroke}{rgb}{0.000000,0.000000,0.000000}%
\pgfsetstrokecolor{currentstroke}%
\pgfsetdash{}{0pt}%
\pgfpathmoveto{\pgfqpoint{0.033333in}{0.363991in}}%
\pgfpathlineto{\pgfqpoint{2.166667in}{0.363991in}}%
\pgfusepath{stroke}%
\end{pgfscope}%
\begin{pgfscope}%
\definecolor{textcolor}{rgb}{0.000000,0.000000,0.000000}%
\pgfsetstrokecolor{textcolor}%
\pgfsetfillcolor{textcolor}%
\pgftext[x=2.166667in,y=0.180927in,right,top]{\color{textcolor}\rmfamily\fontsize{12.000000}{14.400000}\selectfont \(\displaystyle z\)}%
\end{pgfscope}%
\begin{pgfscope}%
\pgfsetrectcap%
\pgfsetroundjoin%
\pgfsetlinewidth{1.505625pt}%
\definecolor{currentstroke}{rgb}{0.121569,0.466667,0.705882}%
\pgfsetstrokecolor{currentstroke}%
\pgfsetdash{}{0pt}%
\pgfpathmoveto{\pgfqpoint{0.033333in}{2.098372in}}%
\pgfpathlineto{\pgfqpoint{0.366667in}{2.098372in}}%
\pgfusepath{stroke}%
\end{pgfscope}%
\begin{pgfscope}%
\definecolor{textcolor}{rgb}{0.000000,0.000000,0.000000}%
\pgfsetstrokecolor{textcolor}%
\pgfsetfillcolor{textcolor}%
\pgftext[x=0.500000in,y=2.040039in,left,base]{\color{textcolor}\rmfamily\fontsize{12.000000}{14.400000}\selectfont \(\displaystyle u\)}%
\end{pgfscope}%
\begin{pgfscope}%
\pgfsetbuttcap%
\pgfsetroundjoin%
\pgfsetlinewidth{1.505625pt}%
\definecolor{currentstroke}{rgb}{1.000000,0.498039,0.054902}%
\pgfsetstrokecolor{currentstroke}%
\pgfsetdash{{5.550000pt}{2.400000pt}}{0.000000pt}%
\pgfpathmoveto{\pgfqpoint{0.033333in}{1.853744in}}%
\pgfpathlineto{\pgfqpoint{0.366667in}{1.853744in}}%
\pgfusepath{stroke}%
\end{pgfscope}%
\begin{pgfscope}%
\definecolor{textcolor}{rgb}{0.000000,0.000000,0.000000}%
\pgfsetstrokecolor{textcolor}%
\pgfsetfillcolor{textcolor}%
\pgftext[x=0.500000in,y=1.795410in,left,base]{\color{textcolor}\rmfamily\fontsize{12.000000}{14.400000}\selectfont \(\displaystyle v\)}%
\end{pgfscope}%
\end{pgfpicture}%
\makeatother%
\endgroup%

    \caption{\textbf{Left:} Monotonic approach of $(u, v)$ to $(\beta, 1-\beta)$ when $\alpha \leq \alpha_c$. \\ \textbf{Right:} Oscillatory approach to the steady state $(\beta, 1-\beta)$, when $\alpha>\alpha_c$.}
    \label{fig:sols}
  \end{figure}
}

\end{document}
